\documentclass[12pt,a4paper]{article}

%%% ——— Packages ———
\usepackage[utf8]{inputenc}
\usepackage[T1]{fontenc}
\usepackage{amsmath,amssymb,amsthm}
\usepackage{mathrsfs}
\usepackage{hyperref}
\usepackage{geometry}
\geometry{margin=1in}
\usepackage{cite}
\usepackage{graphicx}

%%% ——— Theorem environments ———
\newtheorem{theorem}{Theorem}
\newtheorem{proposition}[theorem]{Proposition}
\newtheorem{corollary}[theorem]{Corollary}
\newtheorem{lemma}[theorem]{Lemma}

%%% ——— Title ———
\title{\textbf{Complex Gravitational Coupling from Arithmetic Quantum Gravity Noise}}
\author{Ryan O.\ Van Gelder\thanks{A Citizen Gardens Research Initiative.  Correspondence: Bridgeville, Pennsylvania, USA.}}
\date{February 2026}

\begin{document}
\maketitle

\begin{abstract}
The gravitational coupling $\kappa^{2}=16\pi G_{N}$ has been treated as a real, positive constant since Einstein wrote his field equations in 1915.
We show that this assumption fails when gravity is an open quantum system.
Starting from the Feynman--Vernon influence functional for metric perturbations coupled to a causal quantum environment, we prove that the retarded self-energy dresses $\kappa$ into a complex, frequency-dependent effective coupling
$\kappa_{\mathrm{eff}}(k) = \kappa\bigl/\bigl[1-\kappa\,\tilde{D}_{R}(k)/\mathcal{O}(k)\bigr]$,
whose imaginary part is fixed by the noise kernel and whose real and imaginary parts are linked by a Kramers--Kronig relation rooted in spacetime causality.
The result is a general theorem of stochastic gravity; its physical content becomes sharp when the environment is the Arithmetic--Langevin Group Field Theory (AL-GFT), whose noise kernel is a \emph{Zeta-Comb}---a discrete superposition of log-periodic modes at frequencies set by the imaginary parts $\gamma_{n}$ of the non-trivial Riemann zeta zeros.
In this realisation the beat spectrum of $\operatorname{Im}\kappa_{\mathrm{eff}}$ inherits the GUE pair-correlation statistics of the zeros, providing a binary diagnostic that requires no fitting.
We map the two free parameters $(\varepsilon,\sigma)$ onto a joint parameter space constrained from above by Planck 2018 oscillatory-feature bounds and from below by BBO/DECIGO sensitivity curves, identifying a sweet spot at $\sigma\sim 10^{-3}$, $\varepsilon\sim 10^{-2}$ where both the GUE cross-check and the gravitational-wave signal are accessible.
A self-consistent bootstrap condition---requiring the environment that generates the noise to be compatible with the geometry it modifies---takes the form of a nonlinear eigenvalue problem whose solution, if it exists, would connect the Riemann zero spectrum to the causal structure of spacetime via the Berry--Keating programme.
\end{abstract}


%% ================================================================
%%  SECTION 1 — INTRODUCTION
%% ================================================================
\section{Introduction}\label{sec:intro}

The gravitational coupling $\kappa^{2}=16\pi G_{N}$ has been treated as a real, positive constant since Einstein wrote his field equations in 1915.
Every test of general relativity, from perihelion precession to the LIGO--Virgo--KAGRA catalogue, is consistent with this assumption.
Yet the assumption is never derived: it is imported from classical physics, where $G_{N}$ enters Newton's law as a measured proportionality constant with no internal structure.
The moment one promotes the gravitational field to a quantum degree of freedom---or even treats it classically while allowing it to interact with quantum matter---the coupling to an environment generically dresses every real parameter into a complex, frequency-dependent effective parameter.
This paper asks what happens to $\kappa$ when that generic fact is taken seriously.

%% ——— Motivation 1: the open-system perspective ———
\subsection{Motivation 1: Gravity as an open quantum system}\label{sec:mot1}

Semiclassical gravity replaces the right-hand side of the Einstein equations with the expectation value $\langle\hat{T}_{\mu\nu}\rangle$ of the stress-energy tensor of quantum matter fields \cite{HuVerdaguer2008}.
This mean-field description is adequate when stress-energy fluctuations are small compared with the mean, but it discards all information about correlations.
Stochastic gravity---developed by Hu, Verdaguer, Calzetta and others \cite{HuVerdaguer2004,CalzettaHu1994}---restores the leading fluctuation corrections by deriving, via the Feynman--Vernon influence functional \cite{FeynmonVernon1963}, an Einstein--Langevin equation:
\begin{equation}\label{eq:EL}
  G_{\mu\nu}[g+h] \;=\; \kappa^{2}\,\bigl(\langle\hat{T}_{\mu\nu}\rangle + \xi_{\mu\nu}\bigr),
\end{equation}
where $\xi_{\mu\nu}$ is a stochastic tensor whose two-point function is the \emph{noise kernel}
\begin{equation}\label{eq:noisekernel}
  \langle \xi_{\mu\nu}(x)\,\xi_{\alpha\beta}(x')\rangle \;=\; N_{\mu\nu\alpha\beta}(x,x')
  \;=\; \tfrac{1}{2}\bigl\langle\bigl\{\hat{T}_{\mu\nu}(x),\,\hat{T}_{\alpha\beta}(x')\bigr\}\bigr\rangle.
\end{equation}
The noise kernel is real, symmetric, and positive semi-definite; it encodes the quantum state of the environment.
Its existence is a theorem of quantum field theory in curved spacetime, not a modelling assumption.

The influence functional that generates (\ref{eq:EL}) also contains a \emph{dissipation kernel} $D_{R}$, the retarded part of the stress-energy two-point function, related to the noise kernel by a fluctuation--dissipation relation (FDR) \cite{HuVerdaguer2008}.
In any causal linear-response theory, the retarded response is analytic in the upper half of the complex frequency plane; its real and imaginary parts are therefore linked by the Kramers--Kronig relations \cite{KramersKronig}.
The key observation of this paper is that these two kernels---noise and dissipation---dress the bare gravitational coupling $\kappa$ into a \emph{complex, frequency-dependent effective coupling} $\kappa_{\mathrm{eff}}(k)$, in exact analogy with the way a dielectric medium dresses the bare electric charge into a complex permittivity.

This is not a new physical effect; it is an inescapable consequence of combining two established results:
\begin{enumerate}
\item The influence functional for gravity yields both a noise kernel and a dissipation kernel \cite{HuVerdaguer2004}.
\item Kramers--Kronig causality ties the real and imaginary parts of any retarded response function \cite{KramersKronig}.
\end{enumerate}
What \emph{is} new is the recognition that the resulting $\kappa_{\mathrm{eff}}(k)$ carries physical content---specifically, that its imaginary part and its beat-frequency structure are observable in principle, and that for a particular class of environments they encode number-theoretic information.


%% ——— Motivation 2: the conformal-factor problem ———
\subsection{Motivation 2: The conformal-factor problem and the sign of fluctuations}\label{sec:mot2}

A second, independent motivation comes from the long-standing \emph{conformal-factor problem} of Euclidean quantum gravity \cite{GibbonsHawkingPerry1978,Mazur1990}.
The Einstein--Hilbert action, after a conformal decomposition $g_{\mu\nu}=\Omega^{2}\tilde{g}_{\mu\nu}$, contains a kinetic term for the conformal factor with the ``wrong'' sign:
\begin{equation}\label{eq:SEconformal}
  S_{E} \;=\; -\frac{1}{16\pi G}\int d^{4}x\,\sqrt{\tilde{g}}\,\bigl(\Omega^{2}\tilde{R}+6(\nabla\Omega)^{2}\bigr).
\end{equation}
The negative-definite kinetic energy makes the Euclidean path integral $\int[Dg]\,e^{-S_{E}}$ unbounded from below, rendering the standard Wick rotation ill-defined \cite{ConformalInstability2022}.
Proposed remedies include restricting to on-shell configurations, rotating the conformal-factor contour, or constraining before quantising \cite{GravPI_UCSB}.

Our framework offers a complementary perspective.
In the Lorentzian Schwinger--Keldysh formalism, there is no Euclidean rotation, and the conformal instability manifests instead as a \emph{sign ambiguity in the dissipation kernel} $D_{R}$.
When $D_{R}$ has the ``wrong'' sign relative to the noise kernel, the FDR forces $\operatorname{Im}\kappa_{\mathrm{eff}}<0$---corresponding to \emph{anti-dissipation}, i.e., gravitational amplification of fluctuations.
This is precisely the instability that the conformal-factor problem encodes in Euclidean language.
The advantage of the Lorentzian open-system treatment is that the instability is automatically regulated by the causal structure: $\kappa_{\mathrm{eff}}(k)$ is analytic in the upper half-plane and approaches the real value $\kappa$ at asymptotically high frequencies, so the theory remains well-defined mode by mode.
The conformal-factor problem thus re-emerges as a \emph{spectral feature} of the complex coupling rather than a global pathology of the path integral.


%% ——— Motivation 3: Riemann zeros and quantum gravity ———
\subsection{Motivation 3: Riemann zeros, GUE statistics, and number-theoretic spacetime}\label{sec:mot3}

The third motivation is the oldest open question in the intersection of physics and number theory: the spectral interpretation of the Riemann zeta zeros.

Since Montgomery's 1973 pair-correlation conjecture \cite{Montgomery1973} and its stunning numerical verification by Odlyzko \cite{Odlyzko1987}, the non-trivial zeros of $\zeta(s)$ on the critical line $\operatorname{Re}(s)=\tfrac{1}{2}$ have been known to exhibit the same local statistics as the eigenvalues of large random Hermitian matrices drawn from the Gaussian Unitary Ensemble (GUE) \cite{MontgomeryWiki,MathWorldMontgomery}.
Berry and Keating conjectured that the zeros \emph{are} eigenvalues of a semiclassical Hamiltonian $\hat{H}=\tfrac{1}{2}(\hat{x}\hat{p}+\hat{p}\hat{x})$ \cite{BerryKeating1999}, and Bender, Brody and M\"uller \cite{BBM2017} constructed a $\mathcal{PT}$-symmetric Hamiltonian whose eigenvalues, under a suitable boundary condition, coincide with the zeros.
Rodgers proved \cite{Rodgers2013} that the full GUE conjecture for the $n$-level correlations is logically equivalent to a precise statement about the distribution of products of primes---making the spectral statistics of the zeros a window into the arithmetic structure of the integers.

Despite this remarkable convergence of evidence, no physical system has been identified whose energy spectrum \emph{is} the Riemann zero spectrum in a measurable sense.
This paper does not claim to solve the Riemann hypothesis.
What it does claim is the following: \emph{if the quantum gravitational environment has the spectral content of the Riemann zeros}---as it does in the Arithmetic--Langevin Group Field Theory (AL-GFT) constructed in our companion work \cite{ALGFTGate1}---then the beat spectrum of the imaginary part of $\kappa_{\mathrm{eff}}$ inherits GUE pair-correlation statistics, and this correlation structure is in principle observable in the stochastic gravitational-wave background at decihertz frequencies.

The conditional nature of this claim is essential.
We are not asserting that the Riemann zeros govern gravity.
We are proving that \emph{if they do}, the consequence is a specific, falsifiable pattern in a specific observable.
This transforms the Berry--Keating programme from a conjecture about an unknown Hamiltonian into a prediction about a measurable spectrum.


%% ——— The answer in brief ———
\subsection{Summary of results}\label{sec:summary}

The paper establishes the following chain of results:

\begin{enumerate}
\item \textbf{General theorem (Section~\ref{sec:complexkappa}).}
  For \emph{any} causal quantum environment coupled to linearised gravity via the Feynman--Vernon influence functional, the bare coupling $\kappa$ is dressed into a complex effective coupling
  \begin{equation}\label{eq:keff_intro}
    \kappa_{\mathrm{eff}}(k) \;=\; \frac{\kappa}{1 - \kappa\,\tilde{D}_{R}(k)\big/\mathcal{O}(k)},
  \end{equation}
  whose imaginary part is determined by the noise kernel via a generalised fluctuation--dissipation theorem, and whose real and imaginary parts obey a Kramers--Kronig relation.
  The proof uses only standard results from influence-functional theory and the analyticity of causal response functions; no assumption about the microscopic theory of quantum gravity is required.
  A Ward identity ensures that $\nabla^{\mu}\langle\hat{T}_{\mu\nu}\rangle=0$ is preserved by the dressing.

\item \textbf{Zeta-Comb realisation (Section~\ref{sec:zetacomb}).}
  When the environment is the AL-GFT---a Group Field Theory whose noise kernel is a discrete superposition of log-periodic modes at frequencies $\gamma_{n}=\operatorname{Im}(\rho_{n})$, the imaginary parts of the non-trivial Riemann zeta zeros---the signal amplitude takes the form
  \begin{equation}\label{eq:signal_intro}
    \left|\frac{\delta\kappa}{\kappa}\right| \;=\; \varepsilon^{2}\,e^{-2\sigma\gamma_{1}^{2}},
  \end{equation}
  where $\varepsilon$ is the multiplicity coupling strength, $\sigma$ the soft-resonance width, and $\gamma_{1}=14.1347\ldots$ is the first zeta zero.
  At the sweet-spot parameters $\sigma\sim 10^{-3}$, $\varepsilon\sim 10^{-2}$, this gives $|\delta\kappa/\kappa|\sim 7\times 10^{-5}$.

\item \textbf{GUE diagnostic (Section~\ref{sec:GUE}).}
  The pairwise beat frequencies $\omega_{nm}=|\gamma_{n}-\gamma_{m}|$ of the Zeta-Comb, when normalised by the local mean spacing, reproduce the GUE two-point correlation function $R_{2}(u)=1-(\sin\pi u/\pi u)^{2}$ \cite{MontgomeryWiki}.
  In particular, GUE level repulsion forces the smallest normalised spacing $u_{\min}\sim 0.2$--$0.5$, whereas a generic (Poisson) environment would populate $u\to 0$.
  This binary yes/no test requires no parameter fitting and distinguishes the Riemann-zero environment from any non-arithmetic noise source.

\item \textbf{Detection roadmap (Section~\ref{sec:detection}).}
  The parameter space $(\varepsilon,\sigma)$ is bounded from above by Planck 2018 constraints on oscillatory features in the primordial power spectrum \cite{Planck2018X} and from below by the projected strain sensitivities of BBO \cite{BBO2011} and DECIGO \cite{DECIGO2021}.
  The sweet spot where the GUE cross-check is visible ($\sigma\lesssim 0.002$, requiring $N_{\mathrm{eff}}\geq 2$ contributing zeros) overlaps with the region of maximal signal amplitude, resolving a potential tension between discriminating power and detectability.

\item \textbf{Bootstrap condition (Section~\ref{sec:bootstrap}).}
  Self-consistency---requiring the noise kernel that generates $\kappa_{\mathrm{eff}}$ to be compatible with the geometry it modifies---yields a nonlinear eigenvalue problem
  \begin{equation}\label{eq:bootstrap_intro}
    K(\sigma)\cdot\mathbf{v}(\sigma) = 0, \qquad v_{n}=e^{-2\sigma\gamma_{n}^{2}},
  \end{equation}
  whose structure parallels the Berry--Keating programme: the zeros $\gamma_{n}$ are simultaneously the data that define the environment and the eigenvalues that the bootstrap must reproduce.
  The existence or non-existence of a solution at physical $\sigma$ values would constitute a non-trivial test of the arithmetic structure of spacetime.
\end{enumerate}

\medskip

The remainder of the paper is organised as follows.
Section~\ref{sec:complexkappa} derives the complex effective coupling from the influence functional and establishes the Kramers--Kronig and Ward-identity constraints.
Section~\ref{sec:zetacomb} specialises to the AL-GFT Zeta-Comb environment and derives the signal amplitude and beat-frequency structure.
Section~\ref{sec:GUE} presents the GUE cross-check as a binary diagnostic.
Section~\ref{sec:detection} maps out the parameter space and the detection roadmap.
Section~\ref{sec:FDT} derives the generalised fluctuation--dissipation theorem with beat-frequency corrections.
Section~\ref{sec:bootstrap} formulates the self-consistent bootstrap condition.
Section~\ref{sec:discussion} connects the results to the Berry--Keating programme and discusses falsifiability.

\medskip

\noindent\textit{Notation and conventions.}
We work in natural units $\hbar=c=1$ unless otherwise stated, with metric signature $(-,+,+,+)$.
The gravitational coupling is $\kappa^{2}=16\pi G_{N}$.
Riemann zeta zeros on the critical line are written $\rho_{n}=\tfrac{1}{2}+i\gamma_{n}$ with $\gamma_{n}>0$ labelled in ascending order; the first few values are $\gamma_{1}\approx 14.135$, $\gamma_{2}\approx 21.022$, $\gamma_{3}\approx 25.011$.
Fourier transforms use the convention $\tilde{f}(k)=\int d^{4}x\,e^{ik\cdot x}f(x)$.



%% ================================================================
%%  SECTION 2 — COMPLEX κ FROM STOCHASTIC GRAVITY
%% ================================================================
\section{Complex $\kappa$ from stochastic gravity}\label{sec:complexkappa}

In this section we prove that the gravitational coupling $\kappa$ is generically dressed into a complex, frequency-dependent effective coupling $\kappa_{\mathrm{eff}}(k)$ whenever gravity is treated as an open quantum system.
The derivation uses only three ingredients: (i) the Schwinger--Keldysh closed-time-path (CTP) formalism for the gravitational field, (ii) the Feynman--Vernon influence functional obtained by tracing over a causal quantum environment, and (iii) the analyticity of retarded response functions in the upper half of the complex frequency plane.
No assumption about the microscopic theory of quantum gravity is required; the result holds for any environment whose stress-energy correlators are well-defined distributions.

%% ——— 2.1: Setup ———
\subsection{Setup: linearised gravity as an open system}\label{sec:setup}

Consider a globally hyperbolic spacetime $(\mathcal{M},g_{ab})$ satisfying the semiclassical Einstein equation
\begin{equation}\label{eq:SCE}
  G_{ab}[g] + \Lambda\, g_{ab} - \alpha\, A_{ab}[g] - \beta\, B_{ab}[g]
  \;=\; \kappa^{2}\,\bigl\langle\hat{T}^{R}_{ab}[g]\bigr\rangle,
\end{equation}
where $\kappa^{2}=16\pi G_{N}$, $\Lambda$ is the cosmological constant, $\alpha$ and $\beta$ are the renormalised higher-derivative couplings, and the right-hand side is the renormalised expectation value of the stress-energy tensor of quantum matter in a physically acceptable (Hadamard) state $|\psi[g]\rangle$ \cite{HuVerdaguer2008,HuVerdaguer2004,Wald1994}.

We treat the gravitational field as the \emph{system} and the quantum matter field as the \emph{environment}, and consider metric perturbations $h_{ab}$ around the semiclassical background:
\begin{equation}\label{eq:metricpert}
  g_{ab} \;\longrightarrow\; g_{ab} + h_{ab}.
\end{equation}
On the Schwinger--Keldysh closed time path, we double the degrees of freedom: the forward branch carries $h^{+}_{ab}$ and the backward branch carries $h^{-}_{ab}$.
It is convenient to pass to the Keldysh basis
\begin{equation}\label{eq:Keldysh}
  \bar{h}_{ab} \;=\; \tfrac{1}{2}\bigl(h^{+}_{ab}+h^{-}_{ab}\bigr), \qquad
  \Delta h_{ab} \;=\; h^{+}_{ab}-h^{-}_{ab},
\end{equation}
where $\bar{h}$ is the classical (average) field and $\Delta h$ is the quantum (difference) field that vanishes on shell.

%% ——— 2.2: Influence functional ———
\subsection{The influence functional}\label{sec:IF}

The CTP generating functional for the metric perturbations, after integrating out the quantum matter fields in their initial state $\hat{\rho}_{\mathrm{env}}$, takes the form \cite{HuVerdaguer2004,CalzettaHu1994,FeynmonVernon1963}
\begin{equation}\label{eq:ZCTP}
  Z[J^{+},J^{-}]
  \;=\; \int \mathcal{D}h^{+}\,\mathcal{D}h^{-}\;
  e^{\,i\bigl(S_{\mathrm{grav}}[g+h^{+}]-S_{\mathrm{grav}}[g+h^{-}]\bigr)
     \;+\; i\,S_{\mathrm{IF}}[h^{+},h^{-}]
     \;+\; i\,(J^{+}\cdot h^{+} - J^{-}\cdot h^{-})},
\end{equation}
where $S_{\mathrm{grav}}$ is the gravitational action and $S_{\mathrm{IF}}$ is the \emph{influence action}, defined via
\begin{equation}\label{eq:IFdef}
  e^{\,i\,S_{\mathrm{IF}}[h^{+},h^{-}]}
  \;=\;
  \mathrm{Tr}_{\mathrm{env}}\Bigl[
    \hat{U}[g+h^{+}]\;\hat{\rho}_{\mathrm{env}}\;\hat{U}^{\dagger}[g+h^{-}]
  \Bigr],
\end{equation}
with $\hat{U}[g+h]$ the unitary evolution operator for the matter field on the perturbed background.

For a linear coupling of the metric perturbation to the stress-energy tensor (which is the physical coupling of gravity to matter at leading order), the influence action to second order in $h$ takes the standard Feynman--Vernon form \cite{HuVerdaguer2008,HuVerdaguer2004}.
In the Keldysh basis (\ref{eq:Keldysh}), it reads
\begin{equation}\label{eq:SIF}
  S_{\mathrm{IF}}[\bar{h},\Delta h]
  \;=\;
  \int d^{4}x\,d^{4}x'\;
  \Bigl[
    \Delta h^{ab}(x)\,D^{R}_{abcd}(x,x')\,\bar{h}^{cd}(x')
    \;+\;
    \tfrac{i}{2}\,\Delta h^{ab}(x)\,N_{abcd}(x,x')\,\Delta h^{cd}(x')
  \Bigr],
\end{equation}
where:
\begin{itemize}
\item The \textbf{dissipation kernel}
\begin{equation}\label{eq:DRdef}
  D^{R}_{abcd}(x,x')
  \;=\;
  i\,\theta(x^{0}-x'^{0})\,
  \bigl\langle\bigl[\hat{T}_{ab}(x),\,\hat{T}_{cd}(x')\bigr]\bigr\rangle
\end{equation}
is the retarded part of the stress-energy two-point function.
It is \emph{causal}: $D^{R}_{abcd}(x,x')=0$ whenever $x$ lies outside the causal future of $x'$.

\item The \textbf{noise kernel}
\begin{equation}\label{eq:Ndef}
  N_{abcd}(x,x')
  \;=\;
  \tfrac{1}{2}\,\bigl\langle\bigl\{\hat{T}_{ab}(x),\,\hat{T}_{cd}(x')\bigr\}\bigr\rangle
  \;-\;
  \bigl\langle\hat{T}_{ab}(x)\bigr\rangle\,\bigl\langle\hat{T}_{cd}(x')\bigr\rangle
\end{equation}
is the symmetrised, connected two-point function of the stress-energy fluctuations.
It is real, symmetric in its arguments, and positive semi-definite \cite{HuVerdaguer2004}.
\end{itemize}

\noindent
The first term in (\ref{eq:SIF}) is the \emph{dissipative} part of the influence action: it modifies the retarded equations of motion for $\bar{h}$.
The second term is the \emph{noise} part: it generates stochastic fluctuations.

\medskip

\noindent
\textbf{Remark.}
Equation~(\ref{eq:SIF}) is the gravitational analogue of the Caldeira--Leggett influence action for quantum Brownian motion \cite{CaldeiraLeggett1983}.
The tensor indices make the expressions more involved, but the algebraic structure is identical: a real, retarded dissipation kernel plus an imaginary, symmetric noise kernel.
This universality is the reason the complex-$\kappa$ result holds for any environment.

%% ——— 2.3: Effective equation of motion ———
\subsection{Effective equation of motion and the dressed propagator}\label{sec:effeom}

The equation of motion for the classical (average) metric perturbation $\bar{h}_{ab}$ is obtained by extremising the CTP effective action with respect to $\Delta h_{ab}$ and then setting $\Delta h=0$ (the physical-limit condition of the Schwinger--Keldysh formalism).
From the gravitational action expanded to second order, the kinetic operator for $h_{ab}$ on a background $g_{ab}$ is the Lichnerowicz operator $\hat{\mathcal{G}}$, defined so that the linearised Einstein tensor is
\begin{equation}\label{eq:Lichnerowicz}
  \delta G_{ab} \;=\; -\tfrac{1}{2}\,\hat{\mathcal{G}}_{ab}{}^{cd}\,h_{cd}.
\end{equation}
Including the contribution from the influence action~(\ref{eq:SIF}), the full linearised equation of motion in Fourier space becomes\footnote{We suppress tensor indices for clarity, writing the equation schematically; the full tensorial structure is restored in Appendix~A.}
\begin{equation}\label{eq:eomfull}
  \Bigl[\hat{\mathcal{G}}(k) \;-\; \kappa^{2}\,\tilde{D}^{R}(k)\Bigr]\,\tilde{\bar{h}}(k)
  \;=\;
  \kappa^{2}\,\tilde{\xi}(k),
\end{equation}
where $\tilde{D}^{R}(k)$ is the Fourier transform of the dissipation kernel and $\tilde{\xi}(k)$ is the stochastic source with correlator $\langle\tilde{\xi}(k)\,\tilde{\xi}(k')\rangle = \tilde{N}(k)\,(2\pi)^{4}\delta^{(4)}(k+k')$.

Equation~(\ref{eq:eomfull}) has the structure of a Dyson equation: the bare gravitational propagator $G_{0}(k)=\hat{\mathcal{G}}^{-1}(k)$ is dressed by the retarded self-energy $\Sigma^{R}(k)=\kappa^{2}\,\tilde{D}^{R}(k)$.
The full (dressed) retarded propagator is
\begin{equation}\label{eq:dressedprop}
  G^{R}(k)
  \;=\;
  \bigl[\hat{\mathcal{G}}(k)-\kappa^{2}\,\tilde{D}^{R}(k)\bigr]^{-1}.
\end{equation}

%% ——— 2.4: Resummation into κ_eff ———
\subsection{Resummation: the complex effective coupling}\label{sec:resum}

The key observation is that the self-energy $\Sigma^{R}(k)=\kappa^{2}\,\tilde{D}^{R}(k)$ can be absorbed into a redefinition of the gravitational coupling.
To see this, we factor the kinetic operator as $\hat{\mathcal{G}}(k)=\kappa^{-2}\,\mathcal{O}(k)$, where $\mathcal{O}(k)$ is the differential operator whose eigenvalues are the squares of the graviton mode frequencies (this factorisation is standard in the linearised Einstein equations, where $\kappa^{2}$ multiplies the source, not the kinetic term; see e.g.~\cite{HuVerdaguer2004}).
Then (\ref{eq:dressedprop}) becomes
\begin{equation}\label{eq:factor}
  G^{R}(k)
  \;=\;
  \frac{\kappa^{2}}{\mathcal{O}(k)\bigl[1-\kappa^{2}\,\tilde{D}^{R}(k)/\mathcal{O}(k)\bigr]}
  \;=\;
  \frac{\kappa_{\mathrm{eff}}^{2}(k)}{\mathcal{O}(k)},
\end{equation}
where we have defined the \textbf{effective gravitational coupling}:

\begin{equation}\label{eq:keff}
  \boxed{
  \kappa_{\mathrm{eff}}^{2}(k)
  \;\equiv\;
  \frac{\kappa^{2}}{1 \;-\; \kappa^{2}\,\tilde{D}^{R}(k)\big/\mathcal{O}(k)}
  }
\end{equation}

\medskip

\noindent
This is the central equation of the paper.
Its structure is a geometric resummation, formally identical to the relation between bare and screened charge in quantum electrodynamics:
\begin{equation}\label{eq:QEDanalogy}
  e_{\mathrm{eff}}^{2}(\omega)
  \;=\;
  \frac{e^{2}}{1-e^{2}\,\Pi^{R}(\omega)},
\end{equation}
where $\Pi^{R}$ is the retarded vacuum polarisation.
The analogy is exact: $\kappa$ plays the role of $e$, and the retarded dissipation kernel $\tilde{D}^{R}(k)$ plays the role of the vacuum polarisation $\Pi^{R}(\omega)$.

Since $\tilde{D}^{R}(k)$ is \emph{complex-valued} for real $k$ (its imaginary part is the spectral density of stress-energy fluctuations), the effective coupling $\kappa_{\mathrm{eff}}^{2}(k)$ is generically complex:
\begin{equation}\label{eq:keffdecomp}
  \kappa_{\mathrm{eff}}^{2}(k) \;=\;
  \operatorname{Re}\kappa_{\mathrm{eff}}^{2}(k)
  \;+\;
  i\,\operatorname{Im}\kappa_{\mathrm{eff}}^{2}(k).
\end{equation}
The real part is the dispersive (reactive) correction to Newton's constant---a running gravitational coupling in the Wilsonian sense.
The imaginary part encodes \emph{dissipation}: the rate at which gravitational modes lose (or gain) energy to the quantum environment.

%% ——— 2.5: Kramers-Kronig proof ———
\subsection{Kramers--Kronig relations for $\kappa_{\mathrm{eff}}$}\label{sec:KK}

The dissipation kernel $D^{R}_{abcd}(x,x')$ is retarded by construction: it vanishes when $x$ lies outside the causal future of $x'$.
In Fourier space, this causal support property implies that $\tilde{D}^{R}(k)$, viewed as a function of the complex frequency $k^{0}=\omega+i\eta$, is \emph{analytic in the upper half-plane} $\eta>0$ \cite{KramersKronig,KKproof}.

\begin{proposition}[Kramers--Kronig for $\tilde{D}^{R}$]\label{prop:KKDR}
Let $\tilde{D}^{R}(\omega,\mathbf{k})$ be the Fourier transform of a retarded kernel satisfying
\begin{enumerate}
  \item[(i)] $D^{R}(x,x')=0$ for $x^{0}<x'^{0}$ \textup{(causality)},
  \item[(ii)] $|\tilde{D}^{R}(\omega,\mathbf{k})|\to 0$ as $|\omega|\to\infty$ in the upper half-plane \textup{(regularity at infinity)}.
\end{enumerate}
Then for real $\omega$:
\begin{align}
  \operatorname{Re}\tilde{D}^{R}(\omega,\mathbf{k})
  &\;=\;
  \frac{1}{\pi}\,\mathcal{P}\!\!\int_{-\infty}^{\infty}\frac{d\omega'}{\omega'-\omega}\,
  \operatorname{Im}\tilde{D}^{R}(\omega',\mathbf{k}), \label{eq:KK1}\\[6pt]
  \operatorname{Im}\tilde{D}^{R}(\omega,\mathbf{k})
  &\;=\;
  -\frac{1}{\pi}\,\mathcal{P}\!\!\int_{-\infty}^{\infty}\frac{d\omega'}{\omega'-\omega}\,
  \operatorname{Re}\tilde{D}^{R}(\omega',\mathbf{k}), \label{eq:KK2}
\end{align}
where $\mathcal{P}$ denotes the Cauchy principal value.
\end{proposition}

\begin{proof}
Since $D^{R}(t,\mathbf{x};\,t',\mathbf{x}')=0$ for $t<t'$, its Fourier transform in $\tau=t-t'$ is
\begin{equation}
  \tilde{D}^{R}(\omega,\mathbf{k})
  \;=\;
  \int_{0}^{\infty}d\tau\;e^{i\omega\tau}\,D^{R}(\tau,\mathbf{k}).
\end{equation}
For $\omega$ in the upper half-plane ($\mathrm{Im}\,\omega>0$), the exponential $e^{i\omega\tau}=e^{i\,\mathrm{Re}\,\omega\,\tau}\,e^{-\mathrm{Im}\,\omega\,\tau}$ is damped for $\tau>0$, so the integral converges and defines an analytic function of $\omega$ in $\mathrm{Im}\,\omega>0$.

Condition~(ii) ensures that $\tilde{D}^{R}/(\omega'-\omega)$ vanishes on the semicircular arc at infinity in the upper half-plane.
Closing the contour in the standard way (real axis plus upper semicircle, with an infinitesimal indentation around $\omega'=\omega$) and applying the Cauchy integral theorem gives
\begin{equation}
  0 \;=\; \mathcal{P}\!\!\int_{-\infty}^{\infty}\frac{d\omega'}{\omega'-\omega}\,\tilde{D}^{R}(\omega',\mathbf{k})
  \;-\; i\pi\,\tilde{D}^{R}(\omega,\mathbf{k}).
\end{equation}
Separating real and imaginary parts yields (\ref{eq:KK1})--(\ref{eq:KK2}).
\end{proof}

\medskip

Since $\kappa_{\mathrm{eff}}^{2}(k)$ is a rational function of $\tilde{D}^{R}(k)$ (Eq.~\ref{eq:keff}), and the denominator $1-\kappa^{2}\tilde{D}^{R}/\mathcal{O}$ is analytic wherever $\tilde{D}^{R}$ is (i.e.\ in the upper half-plane), $\kappa_{\mathrm{eff}}^{2}(k)$ inherits upper-half-plane analyticity provided there are no zeros of the denominator there.
The absence of such zeros is guaranteed in the perturbative regime $|\kappa^{2}\tilde{D}^{R}/\mathcal{O}|\ll 1$, which is the regime where the semiclassical expansion is valid.

\begin{corollary}[Kramers--Kronig for $\kappa_{\mathrm{eff}}$]\label{cor:KKkeff}
In the perturbative regime, $\kappa_{\mathrm{eff}}^{2}(\omega,\mathbf{k})$ is analytic in $\mathrm{Im}\,\omega>0$, and its real and imaginary parts are related by the dispersion relations
\begin{align}
  \operatorname{Re}\kappa_{\mathrm{eff}}^{2}(\omega,\mathbf{k})
  &\;=\;
  \kappa^{2} \;+\; \frac{1}{\pi}\,\mathcal{P}\!\!\int_{-\infty}^{\infty}
  \frac{d\omega'}{\omega'-\omega}\,
  \operatorname{Im}\kappa_{\mathrm{eff}}^{2}(\omega',\mathbf{k}), \label{eq:KKkeff1}\\[6pt]
  \operatorname{Im}\kappa_{\mathrm{eff}}^{2}(\omega,\mathbf{k})
  &\;=\;
  -\frac{1}{\pi}\,\mathcal{P}\!\!\int_{-\infty}^{\infty}
  \frac{d\omega'}{\omega'-\omega}\,
  \bigl[\operatorname{Re}\kappa_{\mathrm{eff}}^{2}(\omega',\mathbf{k})-\kappa^{2}\bigr]. \label{eq:KKkeff2}
\end{align}
\end{corollary}

\noindent
The subtraction constant $\kappa^{2}$ in (\ref{eq:KKkeff1}) arises because $\kappa_{\mathrm{eff}}^{2}\to\kappa^{2}$ as $|\omega|\to\infty$; this is the standard once-subtracted dispersion relation.

\medskip

\noindent
\textbf{Physical content.}
The Kramers--Kronig relations~(\ref{eq:KKkeff1})--(\ref{eq:KKkeff2}) state that the \emph{dispersive} correction to $G_{N}$ (real part) and the \emph{dissipative} correction (imaginary part) are not independent: measuring one determines the other.
This is a direct consequence of spacetime causality---the fact that gravitational perturbations propagate on or inside the light cone.
In particular, $\operatorname{Im}\kappa_{\mathrm{eff}}^{2}$ \emph{cannot} be set to zero unless $\operatorname{Re}\kappa_{\mathrm{eff}}^{2}$ is frequency-independent, i.e.\ unless the environment is trivial.
Any non-trivial quantum environment that modifies the running of $G_{N}$ necessarily generates a non-zero imaginary part.

%% ——— 2.6: FDR ———
\subsection{Fluctuation--dissipation relation}\label{sec:FDR}

The noise and dissipation kernels are not independent.
For a thermal environment at temperature $T=1/\beta$, they satisfy the Kubo--Martin--Schwinger (KMS) condition, which in Fourier space reads \cite{HuVerdaguer2008}
\begin{equation}\label{eq:FDR}
  \tilde{N}(\omega,\mathbf{k})
  \;=\;
  \coth\!\Bigl(\frac{\beta\omega}{2}\Bigr)\,\operatorname{Im}\tilde{D}^{R}(\omega,\mathbf{k}).
\end{equation}
At zero temperature ($\beta\to\infty$), this reduces to
\begin{equation}\label{eq:FDR0}
  \tilde{N}(\omega,\mathbf{k})
  \;=\;
  \operatorname{sgn}(\omega)\,\operatorname{Im}\tilde{D}^{R}(\omega,\mathbf{k}).
\end{equation}
In the de Sitter background relevant to inflation, the Gibbons--Hawking temperature $T_{\mathrm{dS}}=H/(2\pi)$ plays the role of $T$; the FDR then links the noise kernel (which is the observable encoding the environment's quantum state) to the imaginary part of the dissipation kernel (which is the quantity that enters $\kappa_{\mathrm{eff}}$).

Combining (\ref{eq:FDR}) with the definition~(\ref{eq:keff}), we obtain a \emph{direct link between the noise kernel and the imaginary part of the effective coupling}:
\begin{equation}\label{eq:ImkeffN}
  \operatorname{Im}\kappa_{\mathrm{eff}}^{2}(\omega,\mathbf{k})
  \;=\;
  \frac{\kappa^{4}\,\tilde{N}(\omega,\mathbf{k})\,\tanh(\beta\omega/2)}
       {\bigl|\mathcal{O}(k)-\kappa^{2}\,\tilde{D}^{R}(k)\bigr|^{2}}\;\mathcal{O}(k),
\end{equation}
valid to all orders in the perturbative resummation.
This equation is the precise sense in which the noise kernel---the observable encoding the environment's quantum state---\emph{determines} the imaginary part of the gravitational coupling.
Its physical content becomes sharp in Section~\ref{sec:zetacomb}, where the noise kernel is the Zeta-Comb.

%% ——— 2.7: Ward identity ———
\subsection{Ward identity: conservation of stress-energy}\label{sec:Ward}

A natural concern is whether the dressing $\kappa\to\kappa_{\mathrm{eff}}(k)$ is consistent with the conservation of the stress-energy tensor, which underpins the Bianchi identity $\nabla^{a}G_{ab}=0$ and hence the consistency of general relativity.

The answer is affirmative, and follows from a Ward identity of the CTP effective action.
Diffeomorphism invariance of the full theory (gravity + matter) implies the contracted Bianchi identity at the quantum level: the divergence of the stress-energy operator vanishes as an operator equation,
\begin{equation}\label{eq:Tcons}
  \nabla^{a}\hat{T}_{ab}[g+h] \;=\; 0.
\end{equation}
Taking the two-point function of both sides, this implies
\begin{align}
  \nabla^{a}_{x}\,D^{R}_{abcd}(x,x') &\;=\; 0, \label{eq:divDR}\\[4pt]
  \nabla^{a}_{x}\,N_{abcd}(x,x') &\;=\; 0. \label{eq:divN}
\end{align}
In Fourier space, these become transversality conditions:
\begin{equation}\label{eq:transverse}
  k^{a}\,\tilde{D}^{R}_{abcd}(k) \;=\; 0, \qquad
  k^{a}\,\tilde{N}_{abcd}(k) \;=\; 0.
\end{equation}

\begin{proposition}[Conservation under dressing]\label{prop:Ward}
The dressed propagator $G^{R}(k)$ defined by~\textup{(\ref{eq:dressedprop})} satisfies the transversality condition
\begin{equation}\label{eq:Wardprop}
  k^{a}\,\bigl[G^{R}(k)\bigr]^{-1}_{abcd} \;=\; k^{a}\,\hat{\mathcal{G}}_{abcd}(k),
\end{equation}
i.e.\ the dressing does not generate longitudinal (unphysical) graviton modes.
Equivalently, the effective Einstein equation with $\kappa_{\mathrm{eff}}(k)$ in place of $\kappa$ is automatically divergence-free.
\end{proposition}

\begin{proof}
From (\ref{eq:dressedprop}),
$[G^{R}]^{-1}_{abcd} = \hat{\mathcal{G}}_{abcd} - \kappa^{2}\tilde{D}^{R}_{abcd}$.
Contracting with $k^{a}$ and using the transversality (\ref{eq:transverse}):
\[
  k^{a}\bigl[\hat{\mathcal{G}}_{abcd}(k) - \kappa^{2}\tilde{D}^{R}_{abcd}(k)\bigr]
  \;=\;
  k^{a}\hat{\mathcal{G}}_{abcd}(k) - \kappa^{2}\,\underbrace{k^{a}\tilde{D}^{R}_{abcd}(k)}_{=\,0}
  \;=\;
  k^{a}\hat{\mathcal{G}}_{abcd}(k).
\]
The last expression is the linearised Bianchi identity for the bare kinetic operator, which vanishes identically.
Hence $k^{a}[G^{R}]^{-1}_{abcd}=0$, proving transversality.
\end{proof}

\noindent
The stochastic source $\xi_{ab}$ inherits the same conservation: $\nabla^{a}\xi_{ab}=0$ with probability one, as shown in~\cite{HuVerdaguer2004}.
Thus the full Einstein--Langevin equation with $\kappa_{\mathrm{eff}}(k)$ preserves all the gauge structure of general relativity.

%% ——— 2.8: Summary theorem ———
\subsection{Summary: the general theorem}\label{sec:theorem}

We collect the results of this section into a single statement.

\begin{theorem}[Complex gravitational coupling from causal environments]\label{thm:main}
Let $(\mathcal{M},g_{ab})$ be a globally hyperbolic spacetime satisfying the semiclassical Einstein equation~\textup{(\ref{eq:SCE})}, and let the quantum matter field be in any physically acceptable \textup{(Hadamard)} state.
Then:
\begin{enumerate}
\item[\textup{(a)}]
The Feynman--Vernon influence functional for metric perturbations yields a dissipation kernel $D^{R}_{abcd}$ and a noise kernel $N_{abcd}$ satisfying the fluctuation--dissipation relation~\textup{(\ref{eq:FDR})}.
\item[\textup{(b)}]
The dressed retarded graviton propagator defines a complex, frequency-dependent effective gravitational coupling
\begin{equation}\label{eq:keffthm}
  \kappa_{\mathrm{eff}}^{2}(k)
  \;=\;
  \frac{\kappa^{2}}{1-\kappa^{2}\,\tilde{D}^{R}(k)/\mathcal{O}(k)},
\end{equation}
whose imaginary part is determined by the noise kernel via~\textup{(\ref{eq:ImkeffN})} and is generically non-zero.
\item[\textup{(c)}]
$\kappa_{\mathrm{eff}}^{2}(\omega,\mathbf{k})$ is analytic in $\mathrm{Im}\,\omega>0$ \textup{(}in the perturbative regime\textup{)}, and its real and imaginary parts obey the once-subtracted Kramers--Kronig dispersion relations~\textup{(\ref{eq:KKkeff1})}--\textup{(\ref{eq:KKkeff2})}.
\item[\textup{(d)}]
The Ward identity~\textup{(\ref{eq:transverse})} ensures that the dressing preserves $\nabla^{a}G_{ab}=0$; no unphysical longitudinal graviton modes are generated.
\end{enumerate}
\end{theorem}

\noindent
The theorem establishes that a complex gravitational coupling is not an exotic possibility but an \emph{unavoidable consequence} of treating gravity as an open quantum system---which it necessarily is, since every graviton interacts with quantum matter.
The physical content of $\kappa_{\mathrm{eff}}$ depends entirely on the quantum state of the environment, encoded in the noise kernel $N_{abcd}$.
In Section~\ref{sec:zetacomb} we specialise to the environment provided by the Arithmetic--Langevin Group Field Theory, where the noise kernel is a Zeta-Comb and $\kappa_{\mathrm{eff}}$ acquires number-theoretic structure.

%% ——— New bibliography entries needed for §2 ———
%% (to be appended to the existing bibliography)
%%
%% \bibitem{Wald1994}
%% R.~M.~Wald,
%% \emph{Quantum Field Theory in Curved Spacetime and Black Hole Thermodynamics},
%% University of Chicago Press (1994).
%%
%% \bibitem{CaldeiraLeggett1983}
%% A.~O.~Caldeira and A.~J.~Leggett,
%% ``Path integral approach to quantum Brownian motion,''
%% Physica A \textbf{121}, 587 (1983).
%%
%% \bibitem{KKproof}
%% A.~Rustagi,
%% ``Kramers-Kronig Relations,''
%% Lecture Notes, NCSU (2018),
%% \url{http://www.phys.ufl.edu/~avinash/Notes/KramersKronig/KramersKronig.pdf}.
%%


%% §2 of "Complex Gravitational Coupling from Arithmetic Quantum Gravity Noise"
%% Requires: amsmath, amssymb, amsthm, mathrsfs, hyperref
%% Place after §1 (Introduction/Motivation)

%%---------------------------------------------------------------------
\section{Complex $\kappa$ from Stochastic Gravity}\label{sec:complex-kappa}
%%---------------------------------------------------------------------

The gravitational coupling $\kappa^{2}=16\pi G_{N}$ enters Einstein's
field equations as a real, positive constant.  We show in this section
that treating gravity as an \emph{open} quantum system---coupled to any
causal quantum environment---forces $\kappa$ to acquire an imaginary
part.  The argument proceeds in three stages: (i)~we set up the
influence-functional formalism for linearised gravity coupled to an
arbitrary environment; (ii)~we resum the dressed gravitational
propagator and identify the effective coupling
$\kappa_{\mathrm{eff}}(k)$; (iii)~we invoke the Kramers--Kronig
relations to prove that $\kappa_{\mathrm{eff}}$ is necessarily complex
whenever the environment is causal and dissipative.  The result is a
\emph{theorem}, independent of the specific environment; its
specialisation to the Arithmetic-Langevin environment is deferred to
\S\ref{sec:zeta-comb}.

%%---------------------------------------------------------------------
\subsection{Gravity as an open quantum system}\label{ssec:open-gravity}
%%---------------------------------------------------------------------

We work in linearised gravity on a cosmological (FLRW) background.
Let $\phi(\eta,\mathbf{k})$ denote the gauge-invariant curvature
perturbation in Fourier space, with conformal time~$\eta$ and comoving
wavevector~$\mathbf{k}$.  Following the stochastic-gravity programme
of Hu and Verdaguer~\cite{HuVerdaguer2008}, we partition the total
system into

\begin{itemize}
\item \textbf{System $\boldsymbol{S}$}: the long-wavelength metric
  perturbation $\phi$, governed by a quadratic action
  %
  \begin{equation}\label{eq:Ssys}
    S_{\mathrm{sys}}[\phi]
    = \frac{1}{2}\int d\eta\,\frac{d^{3}k}{(2\pi)^{3}}\;
      a^{2}\!\left[\dot{\phi}_{\mathbf{k}}^{2}
      - c_{s}^{2}\,k^{2}\,\phi_{\mathbf{k}}^{2}\right],
  \end{equation}
  %
  where $a(\eta)$ is the scale factor and $c_{s}$ the adiabatic sound
  speed.

\item \textbf{Environment $\boldsymbol{E}$}: a collection of
  short-wavelength or internal degrees of freedom
  $\{\chi_{n,\mathbf{k}}\}$, labelled by a discrete index~$n$, with
  quadratic action
  %
  \begin{equation}\label{eq:Senv}
    S_{\mathrm{env}}[\chi]
    = \frac{1}{2}\int d\eta\,\frac{d^{3}k}{(2\pi)^{3}}\;
      \sum_{n}\left[\dot{\chi}_{n,\mathbf{k}}^{2}
      - \Omega_{n}^{2}(k,\eta)\,\chi_{n,\mathbf{k}}^{2}\right],
  \end{equation}
  %
  where the dispersion $\Omega_{n}(k,\eta)$ encodes the microscopic
  physics of the environment.  No assumption is made about its origin
  at this stage; it could arise from quantum matter loops,
  spin-foam/GFT degrees of freedom, or any other UV completion.

\item \textbf{Interaction}: a linear coupling of $\phi$ to a
  collective environment operator
  $\mathcal{O}(\eta,\mathbf{k})\equiv\sum_{n}g_{n}\,\chi_{n,\mathbf{k}}$,
  %
  \begin{equation}\label{eq:Sint}
    S_{\mathrm{int}}[\phi,\chi]
    = \kappa\int d\eta\,\frac{d^{3}k}{(2\pi)^{3}}\;
      a^{2}\,\phi_{\mathbf{k}}\,
      \mathcal{O}_{-\mathbf{k}}\,,
  \end{equation}
  %
  where $\kappa$ is the bare gravitational coupling and $g_{n}$ are
  dimensionless weights carrying the environment's spectral
  information.
\end{itemize}

\noindent
The total action
$S_{\mathrm{tot}}=S_{\mathrm{sys}}+S_{\mathrm{env}}+S_{\mathrm{int}}$
is at most quadratic in~$\chi$, so the environment can be integrated
out \emph{exactly}.

%%---------------------------------------------------------------------
\subsection{Influence functional and
  Schwinger--Keldysh action}\label{ssec:influence-functional}
%%---------------------------------------------------------------------

On the closed time path (CTP), we double the fields:
$\phi^{\pm}$, $\chi^{\pm}$.  The CTP generating functional is
\begin{equation}\label{eq:CTP-Z}
  Z[J^{+},J^{-}]
  = \int\!\mathcal{D}\phi^{+}\mathcal{D}\phi^{-}\;
    e^{\,i\bigl(S_{\mathrm{sys}}[\phi^{+}]
      -S_{\mathrm{sys}}[\phi^{-}]\bigr)
      +\,i\,S_{\mathrm{IF}}[\phi^{+},\phi^{-}]
      +\,i\!\int\!(J^{+}\phi^{+}-J^{-}\phi^{-})}\,,
\end{equation}
where the \textbf{Feynman--Vernon influence functional}~\cite{FeynmanVernon1963}
\begin{equation}\label{eq:IF-def}
  \mathcal{F}[\phi^{+},\phi^{-}]
  \;\equiv\; e^{\,i\,S_{\mathrm{IF}}[\phi^{+},\phi^{-}]}
  \;=\; \int\!\mathcal{D}\chi^{+}\mathcal{D}\chi^{-}\;
  e^{\,i\bigl(S_{\mathrm{env+int}}[\phi^{+},\chi^{+}]
    -S_{\mathrm{env+int}}[\phi^{-},\chi^{-}]\bigr)}\,
  \hat{\rho}_{\mathrm{env}}[\chi^{+},\chi^{-}]
\end{equation}
encapsulates the entire back-reaction of the environment on the
system.

Because $S_{\mathrm{env}}+S_{\mathrm{int}}$ is quadratic in~$\chi$,
the path integral over $\chi^{\pm}$ is Gaussian and yields an
\emph{exact} influence action.  Introducing the Keldysh variables
\begin{equation}\label{eq:Keldysh-vars}
  \phi_{c}\equiv\tfrac{1}{2}(\phi^{+}+\phi^{-})\,,\qquad
  \Delta\phi\equiv\phi^{+}-\phi^{-}\,,
\end{equation}
the standard result~\cite{CalzettaHu2008,HuVerdaguer2008} is
\begin{equation}\label{eq:SIF}
  \boxed{
  S_{\mathrm{IF}}[\phi_{c},\Delta\phi]
  = \int\!d\eta\,d\eta'\,\frac{d^{3}k}{(2\pi)^{3}}\;
    \biggl[
      \Delta\phi_{\mathbf{k}}(\eta)\;
      D_{R}(\eta,\eta';k)\;\phi_{c,-\mathbf{k}}(\eta')
      + \frac{i}{2}\,\Delta\phi_{\mathbf{k}}(\eta)\;
      N(\eta,\eta';k)\;\Delta\phi_{-\mathbf{k}}(\eta')
    \biggr],}
\end{equation}
with two kernels:

\begin{enumerate}
\item The \textbf{dissipation kernel}
  (retarded, real):
  \begin{equation}\label{eq:DR-def}
    D_{R}(\eta,\eta';k)
    = i\,\kappa^{2}\,\Theta(\eta-\eta')\;
      \bigl\langle\bigl[\mathcal{O}_{\mathbf{k}}(\eta),\;
      \mathcal{O}_{-\mathbf{k}}(\eta')\bigr]\bigr\rangle_{\mathrm{env}},
  \end{equation}
  which governs the causal, non-conservative part of the gravitational
  response.

\item The \textbf{noise kernel}
  (symmetric, positive semi-definite):
  \begin{equation}\label{eq:N-def}
    N(\eta,\eta';k)
    = \frac{\kappa^{2}}{2}\;
      \bigl\langle\bigl\{\mathcal{O}_{\mathbf{k}}(\eta),\;
      \mathcal{O}_{-\mathbf{k}}(\eta')\bigr\}\bigr\rangle_{\mathrm{env}},
  \end{equation}
  which sources stochastic fluctuations in the metric.
\end{enumerate}

\noindent
These two kernels are \emph{not} independent: by construction they
satisfy the fluctuation--dissipation relation (FDR)
\begin{equation}\label{eq:FDR-basic}
  N(\omega,k)
  = \coth\!\biggl(\frac{\omega}{2T_{\mathrm{eff}}}\biggr)\,
    \mathrm{Im}\,D_{R}(\omega,k)\,,
\end{equation}
where $T_{\mathrm{eff}}$ is the effective temperature of the
environment (e.g.\ the Gibbons--Hawking temperature $H/2\pi$ for a
de~Sitter background).

%%---------------------------------------------------------------------
\subsection{Dressed gravitational propagator and effective coupling}%
\label{ssec:resummation}
%%---------------------------------------------------------------------

The influence action~\eqref{eq:SIF} modifies the inverse propagator of
$\phi$.  In Fourier space $(\omega,k)$, the bare inverse propagator
from $S_{\mathrm{sys}}$ is
\begin{equation}\label{eq:G0inv}
  G_{0}^{-1}(\omega,k)
  = a^{2}\bigl(\omega^{2}-c_{s}^{2}\,k^{2}\bigr)\,,
\end{equation}
and the dissipation kernel contributes an additional self-energy.
The full dressed inverse propagator is
\begin{equation}\label{eq:Ginv-dressed}
  G^{-1}(\omega,k)
  = G_{0}^{-1}(\omega,k) - \kappa\,\widetilde{D}_{R}(\omega,k)\,,
\end{equation}
where $\widetilde{D}_{R}$ is the Fourier transform of $D_{R}$.
Factoring out the bare kinetic structure, we write
\begin{equation}\label{eq:Geff}
  G^{-1}(\omega,k)
  = a^{2}\bigl(\omega^{2}-c_{s}^{2}\,k^{2}\bigr)
    \left[
      1 - \frac{\kappa\,\widetilde{D}_{R}(\omega,k)}
               {a^{2}\bigl(\omega^{2}-c_{s}^{2}\,k^{2}\bigr)}
    \right].
\end{equation}

This invites the identification of an \textbf{effective gravitational
coupling} via a geometric resummation of the self-energy insertions.
Denoting $\Pi(\omega,k)\equiv
\widetilde{D}_{R}(\omega,k)/\bigl[a^{2}(\omega^{2}-c_{s}^{2}k^{2})\bigr]$,
the resummed propagator has the Dyson form
\begin{equation}\label{eq:Dyson}
  G(\omega,k)
  = \frac{G_{0}(\omega,k)}{1-\kappa\,\Pi(\omega,k)}\,,
\end{equation}
which defines
\begin{equation}\label{eq:kappa-eff}
  \boxed{
  \kappa_{\mathrm{eff}}(\omega,k)
  \;\equiv\;
  \frac{\kappa}{1-\kappa\,\Pi(\omega,k)}\,.}
\end{equation}

\noindent
\textbf{Key observation:} The self-energy $\Pi(\omega,k)$ is a
retarded response function.  It is therefore \emph{generically complex}
for real $\omega$.  Consequently, $\kappa_{\mathrm{eff}}$ is complex:
its real part describes the renormalised gravitational strength, and
its imaginary part describes gravitational dissipation into the
environment.

\begin{remark}[Perturbative regime]
Throughout this work, we operate in the regime
$|\kappa\,\Pi|\ll 1$, where the Dyson resummation is well-controlled
and the perturbative expansion
\begin{equation}\label{eq:kappa-pert}
  \kappa_{\mathrm{eff}}
  \approx \kappa\bigl(1+\kappa\,\Pi+\kappa^{2}\Pi^{2}+\cdots\bigr)
\end{equation}
converges.  At leading non-trivial order,
$\delta\kappa/\kappa=\kappa\,\Pi(\omega,k)$.
\end{remark}

%%---------------------------------------------------------------------
\subsection{Kramers--Kronig proof: complex $\kappa$ is
  mandatory}\label{ssec:KK-proof}
%%---------------------------------------------------------------------

We now elevate the observation of \S\ref{ssec:resummation} to a
\emph{theorem}.

\begin{theorem}[Mandatory complexity of $\kappa_{\mathrm{eff}}$]%
\label{thm:complex-kappa}
Let $\Pi(\omega,k)$ be the retarded self-energy of the gravitational
perturbation induced by \emph{any} causal quantum environment with a
non-trivial spectral density.  Then
$\mathrm{Im}\,\kappa_{\mathrm{eff}}(\omega,k)\neq 0$ for all
$\omega$ in the support of the environment's spectral function.
\end{theorem}

\begin{proof}
The proof consists of three steps.

\medskip\noindent\textbf{Step 1: Causality $\Rightarrow$ analyticity.}
The dissipation kernel $D_{R}(\eta,\eta';k)$ is retarded
[Eq.~\eqref{eq:DR-def}], i.e.\
$D_{R}(\eta,\eta';k)=0$ for $\eta<\eta'$.  By Titchmarsh's
theorem~\cite{Titchmarsh1948}, its Fourier transform
$\widetilde{D}_{R}(\omega,k)$ is analytic in the upper half of the
complex $\omega$-plane:
\begin{equation}\label{eq:analyticity}
  \widetilde{D}_{R}(\omega,k)\;\text{analytic for}\;
  \mathrm{Im}\,\omega>0\,.
\end{equation}
Since $\Pi=\widetilde{D}_{R}/G_{0}^{-1}$ inherits this analyticity
(the bare propagator $G_{0}^{-1}$ is a polynomial in $\omega$), the
self-energy $\Pi(\omega,k)$ is also analytic in the upper half-plane
and vanishes as $|\omega|\to\infty$ (the environment cannot respond
faster than it is driven).

\medskip\noindent\textbf{Step 2: Kramers--Kronig relations.}
From the analyticity and asymptotic decay of~$\Pi$, the standard
Kramers--Kronig (KK) dispersion relations~\cite{KramersKronig,Nussenzveig1972}
hold:
\begin{subequations}\label{eq:KK}
\begin{align}
  \mathrm{Re}\,\Pi(\omega,k)
  &= \frac{1}{\pi}\,\mathcal{P}\!\!\int_{-\infty}^{\infty}
     \frac{\mathrm{Im}\,\Pi(\omega',k)}{\omega'-\omega}\;d\omega'\,,
  \label{eq:KK-Re}\\[4pt]
  \mathrm{Im}\,\Pi(\omega,k)
  &= -\frac{1}{\pi}\,\mathcal{P}\!\!\int_{-\infty}^{\infty}
     \frac{\mathrm{Re}\,\Pi(\omega',k)}{\omega'-\omega}\;d\omega'\,,
  \label{eq:KK-Im}
\end{align}
\end{subequations}
where $\mathcal{P}$ denotes the Cauchy principal value.  The two
relations are the Hilbert-transform pair that encodes the content of
causality in the frequency domain.

\medskip\noindent\textbf{Step 3: Non-trivial environment
  $\Rightarrow$ $\mathrm{Im}\,\Pi\neq 0$.}
A non-trivial environment has a spectral density $\rho(\omega,k)>0$
on some open interval of~$\omega$.  By the optical theorem for the
self-energy (or equivalently, by unitarity of the underlying quantum
theory),
\begin{equation}\label{eq:optical}
  \mathrm{Im}\,\Pi(\omega,k)
  = -\pi\,\kappa^{2}\,\rho(\omega,k)\,,
\end{equation}
where $\rho(\omega,k)=
(2\pi)^{-1}\sum_{n}|g_{n}|^{2}\,\delta(\omega-\Omega_{n}(k))$ is the
spectral function of $\mathcal{O}$.  Since $\rho>0$ on the support of
the environment, $\mathrm{Im}\,\Pi\neq 0$, and by
Eq.~\eqref{eq:KK-Re}, $\mathrm{Re}\,\Pi$ is non-trivially determined
by the Hilbert transform of $\mathrm{Im}\,\Pi$.

From Eq.~\eqref{eq:kappa-eff},
\begin{equation}\label{eq:Im-kappa}
  \mathrm{Im}\,\kappa_{\mathrm{eff}}
  = \frac{\kappa^{2}\,\mathrm{Im}\,\Pi}
         {|1-\kappa\,\Pi|^{2}}
  \;\neq\; 0\,.
\end{equation}
This completes the proof.
\end{proof}

\begin{remark}[Generality]
Theorem~\ref{thm:complex-kappa} makes no reference to the specific
nature of the environment.  It applies equally to quantum matter
fields in semiclassical gravity~\cite{HuVerdaguer2008}, to GFT/spin-foam
degrees of freedom~\cite{Oriti2012}, or to the Arithmetic-Langevin
environment introduced in \S\ref{sec:zeta-comb}.  The only input is
\emph{causality} and \emph{non-trivial coupling}---both of which are
physically unavoidable for any environment interacting with gravity.
\end{remark}

\begin{remark}[Physical content of $\mathrm{Im}\,\kappa_{\mathrm{eff}}$]
The imaginary part of the effective coupling has a direct physical
interpretation.  It controls:
\begin{enumerate}
  \item \emph{Gravitational dissipation}: energy leaks from
    long-wavelength metric perturbations into the environment at a
    rate $\Gamma\propto\mathrm{Im}\,\kappa_{\mathrm{eff}}$.
  \item \emph{Decoherence}: the noise kernel $N$ associated with the
    dissipation via the FDR~\eqref{eq:FDR-basic} destroys quantum
    coherence of superposed geometries.
  \item \emph{Spectral imprint}: the frequency dependence of
    $\mathrm{Im}\,\kappa_{\mathrm{eff}}(\omega,k)$ is a direct
    fingerprint of the environment's spectral density---and therefore
    a potential observable.
\end{enumerate}
\end{remark}

%%---------------------------------------------------------------------
\subsection{Ward identity and conservation}\label{ssec:ward}
%%---------------------------------------------------------------------

A legitimate concern is whether the complexification of $\kappa$
violates the conservation laws that follow from diffeomorphism
invariance.  We show that it does not.

The linearised diffeomorphism Ward identity requires the divergence of
the effective gravitational equations to vanish:
\begin{equation}\label{eq:ward-div}
  k_{\mu}\,\Gamma^{\mu\nu}[\phi]=0\,,
\end{equation}
where $\Gamma^{\mu\nu}$ is the 1PI effective action vertex.  In the
influence-functional formalism, this identity is preserved
order-by-order in the CTP perturbation theory, because the
coupling~\eqref{eq:Sint} is constructed from the gauge-invariant
perturbation~$\phi$.  The environment-induced self-energy $\Pi$ enters
as a scalar function of $(\omega,k)$ multiplying the
\emph{transverse-traceless} projection of the graviton propagator.
Consequently,
\begin{equation}\label{eq:ward-preserved}
  k_{\mu}\bigl[\kappa_{\mathrm{eff}}(\omega,k)\,T^{\mu\nu}\bigr]
  = \kappa_{\mathrm{eff}}(\omega,k)\,k_{\mu}\,T^{\mu\nu}
  = 0\,,
\end{equation}
since stress-energy conservation $k_{\mu}T^{\mu\nu}=0$ is a
property of the matter sector, independent of the value---real or
complex---of the coupling.

More formally, the influence functional preserves the BRST symmetry of
the gravitational sector.  The Slavnov--Taylor identities of the
dressed theory are identical in form to those of the bare theory; the
only modification is the replacement
$\kappa\to\kappa_{\mathrm{eff}}(\omega,k)$, which is a
$c$-number function and commutes with all
generators~\cite{CalzettaHu2008}.

\begin{proposition}[Conservation guarantee]\label{prop:conservation}
The effective Einstein equation with complex
$\kappa_{\mathrm{eff}}(\omega,k)$ satisfies the linearised contracted
Bianchi identity, $\nabla_{\mu}G^{\mu\nu}_{\mathrm{eff}}=0$, to the
same order as the perturbative expansion of the self-energy.
\end{proposition}

%%---------------------------------------------------------------------
\subsection{Summary of the general result}\label{ssec:sec2-summary}
%%---------------------------------------------------------------------

The chain of logic is:
\begin{equation}\label{eq:logic-chain}
  \underbrace{\text{Causal environment}}_{\text{retarded }D_{R}}
  \;\;\xRightarrow{\;\text{Titchmarsh}\;}
  \;\;\underbrace{\Pi\;\text{analytic in UHP}}_{\text{KK relations}}
  \;\;\xRightarrow{\;\rho>0\;}
  \;\;\underbrace{\mathrm{Im}\,\kappa_{\mathrm{eff}}\neq 0}_{\text{complex coupling}}\,.
\end{equation}
This is a \emph{model-independent} result.  The only freedom is in
the choice of spectral density $\rho(\omega,k)$, which determines the
\emph{pattern}---but not the \emph{existence}---of the imaginary part.
In the next section, we specialise to the Arithmetic-Langevin
environment, where $\rho$ is a Zeta-Comb built from the non-trivial
zeros of the Riemann zeta function, and derive the specific
observational signatures that follow.


%% §3 of "Complex Gravitational Coupling from Arithmetic Quantum Gravity Noise"
%% Requires: amsmath, amssymb, amsthm, mathrsfs, hyperref
%% Place after §2 (Complex κ from Stochastic Gravity)

%%=====================================================================
\section{The Zeta-Comb Environment}\label{sec:zeta-comb}
%%=====================================================================

Section~\ref{sec:complex-kappa} established that \emph{any} causal
quantum environment renders the gravitational coupling complex.  We now
specialise to the environment dictated by the Arithmetic-Langevin
Group Field Theory (AL-GFT): a discrete tower of oscillator modes
whose frequencies and couplings are built from the non-trivial zeros
of the Riemann zeta function.  This choice is not
\emph{ad hoc}---it is the unique environment that arises when the
microscopic GFT vertex amplitudes carry arithmetic (prime-indexed)
weights, as prescribed by Multiplicity Theory.

The section is organised as follows.
In \S\ref{ssec:zeta-environment} we define the environment and
compute the noise kernel explicitly.
In \S\ref{ssec:gen-FDT} we derive the generalised
fluctuation--dissipation relation that connects noise and dissipation
in the Zeta-Comb setting, revealing a family of \emph{beat
  frequencies} between distinct Riemann zeros.
In \S\ref{ssec:beat-spectrum} we characterise the resulting beat
spectrum and show that it carries a GUE fingerprint testable against
random-matrix theory.

%%---------------------------------------------------------------------
\subsection{The Arithmetic-Langevin environment}\label{ssec:zeta-environment}
%%---------------------------------------------------------------------

Recall from \S\ref{ssec:open-gravity} that the environment consists of
modes $\chi_{n,\mathbf{k}}$ with collective operator
$\mathcal{O}=\sum_{n}g_{n}\,\chi_{n}$.  In the AL-GFT realisation, the
index~$n$ labels the non-trivial zeros $\rho_{n}=\tfrac{1}{2}+i\gamma_{n}$
of the Riemann zeta function, with $\gamma_{1}=14.1347\ldots$,
$\gamma_{2}=21.0220\ldots$, etc.  The coupling constants are
\begin{equation}\label{eq:gn}
  \boxed{
  g_{n}
  = \varepsilon\;
    e^{-\sigma\,\gamma_{n}^{2}}\;
    e^{\,i\,\varphi_{n}}\,,
  \qquad
  \varphi_{n}=\tfrac{1}{2}+i\gamma_{n}\,,}
\end{equation}
where $\varepsilon\ll 1$ is the overall multiplicity coupling strength
and $\sigma>0$ is the soft-resonance width that provides a Gaussian
damping at large~$n$.  The phase $\varphi_{n}$ is precisely the
non-trivial zero itself, evaluated on the critical line.

\begin{remark}[Motivation for the couplings]\label{rem:gn-motivation}
The exponential $e^{-\sigma\gamma_{n}^{2}}$ ensures convergence of
all sums over zeros and defines the effective number of contributing
zeros,
\begin{equation}\label{eq:Neff}
  N_{\mathrm{eff}}(\sigma)
  \equiv \sum_{n=1}^{\infty}e^{-2\sigma\,\gamma_{n}^{2}}\,.
\end{equation}
For $\sigma\sim 10^{-3}$, which we will identify as the
phenomenological sweet spot in \S\ref{sec:parameter-space},
$N_{\mathrm{eff}}\approx 4.2$ zeros contribute, generating 7~pairwise
beat frequencies.  The phase $\varphi_{n}$ is \emph{not} a free
parameter; it is fixed by the requirement that the noise kernel
reproduce the Zeta-Comb modulation of the primordial power spectrum
derived in Gate~1 of the AL-GFT programme.
\end{remark}

%%---------------------------------------------------------------------
\subsubsection{Noise kernel}\label{sssec:noise-kernel}
%%---------------------------------------------------------------------

The noise kernel~\eqref{eq:N-def} evaluated for the AL-GFT
environment is
\begin{equation}\label{eq:N-ALGFT}
  N(\eta,\eta';k)
  = \frac{\kappa^{2}}{2}\sum_{n,m}g_{n}\,g_{m}^{*}\;
    \bigl\langle\bigl\{\chi_{n,\mathbf{k}}(\eta),\;
    \chi_{m,-\mathbf{k}}(\eta')\bigr\}\bigr\rangle_{\mathrm{env}}\,.
\end{equation}
For independent harmonic-oscillator modes in the Bunch--Davies vacuum
on a quasi-de~Sitter background, the equal-time limit
$\eta\to\eta'$ yields (after the standard conformal-time-to-$k$
mapping $\eta_{*}\approx -1/k$)
\begin{equation}\label{eq:Nk-sum}
  N(k)
  = \frac{\kappa^{2}\varepsilon^{2}}{2}\;
    \sum_{n=1}^{N_{\mathrm{cut}}}
    e^{-2\sigma\gamma_{n}^{2}}\;
    \left|\mathcal{A}_{n}(k)\right|^{2}\,,
\end{equation}
where $\mathcal{A}_{n}(k)$ is the slowly varying Hankel-function
amplitude evaluated near horizon crossing (see
Appendix~\ref{app:mode-functions}), and $N_{\mathrm{cut}}$ is a
technical cutoff that can be sent to infinity for $\sigma>0$.

The full $k$-dependent noise spectrum---including the oscillatory
phases from the mode functions---takes the form
\begin{equation}\label{eq:Nk-full}
  \boxed{
  N(k)
  = N_{0}(k)\;\mathcal{M}(k)\,,}
\end{equation}
where $N_{0}(k)$ is the standard scale-invariant noise spectrum of
slow-roll inflation, and the \textbf{Zeta-Comb modulation} is
\begin{equation}\label{eq:Mk}
  \boxed{
  \mathcal{M}(k)
  = 1 + 2\varepsilon^{2}\sum_{n=1}^{\infty}
    e^{-2\sigma\gamma_{n}^{2}}\;
    \cos\!\left(\gamma_{n}\ln\frac{k}{k_{*}}+\varphi_{n}\right)
    + \mathcal{O}(\varepsilon^{4})\,.}
\end{equation}
This is a \emph{log-periodic} modulation of the primordial power
spectrum at frequencies set by the imaginary parts of Riemann zeta
zeros.  In the limit $\varepsilon\to 0$, the environment decouples
and $\mathcal{M}\to 1$, recovering the standard $\Lambda$CDM
spectrum.

\begin{remark}[Primordial power spectrum]
The primordial curvature power spectrum inherits the Zeta-Comb
structure via the relation $\mathcal{P}(k)=\mathcal{P}_{0}(k)\,\mathcal{M}(k)$,
where $\mathcal{P}_{0}(k)=A_{s}(k/k_{*})^{n_{s}-1}$ is the standard
nearly scale-invariant spectrum.  The modulation amplitude is
controlled by $\varepsilon^{2}$, which is constrained by Planck~2018
data to satisfy $\varepsilon<0.21$ at the 95\% confidence level for
$\sigma=10^{-3}$.
\end{remark}

%%---------------------------------------------------------------------
\subsubsection{Dissipation kernel}\label{sssec:dissipation-kernel}
%%---------------------------------------------------------------------

The dissipation kernel~\eqref{eq:DR-def} for the AL-GFT environment
is similarly computed from the commutator of the collective operator:
\begin{equation}\label{eq:DR-ALGFT}
  D_{R}(\omega,k)
  = \kappa^{2}\varepsilon^{2}\sum_{n=1}^{\infty}
    e^{-2\sigma\gamma_{n}^{2}}\;
    \frac{\omega}{\omega^{2}-\Omega_{n}^{2}(k)+i\,0^{+}}\,,
\end{equation}
where $\Omega_{n}(k)$ is the effective frequency of the $n$-th
environment mode (dependent on the inflationary background through the
scale factor).  The retarded $i\,0^{+}$ prescription ensures
causality.  The spectral density of the environment is therefore a
\textbf{Zeta-Comb}:
\begin{equation}\label{eq:rho-zeta}
  \rho(\omega,k)
  = \varepsilon^{2}\sum_{n=1}^{\infty}
    e^{-2\sigma\gamma_{n}^{2}}\;
    \delta\bigl(\omega-\Omega_{n}(k)\bigr)\,,
\end{equation}
a discrete, arithmetically-structured spectral density with peaks at
frequencies determined by the Riemann zeros.

%%---------------------------------------------------------------------
\subsection{Generalised fluctuation--dissipation theorem}%
\label{ssec:gen-FDT}
%%---------------------------------------------------------------------

In the standard FDT~\eqref{eq:FDR-basic}, the noise and dissipation
kernels are related at each frequency independently.  For the
Zeta-Comb environment, the discrete, multi-frequency structure of
$\rho(\omega,k)$ induces \emph{cross-frequency} correlations that
enrich the standard relation.

%%---------------------------------------------------------------------
\subsubsection{Derivation}\label{sssec:FDT-derivation}
%%---------------------------------------------------------------------

Substituting the explicit forms~\eqref{eq:Nk-full}
and~\eqref{eq:DR-ALGFT} into the definition of
$\kappa_{\mathrm{eff}}$~\eqref{eq:kappa-eff}, the imaginary part of
the effective coupling reads
\begin{equation}\label{eq:Im-kappa-ALGFT}
  \mathrm{Im}\,\kappa_{\mathrm{eff}}(\omega,k)
  = \frac{\kappa^{3}\varepsilon^{2}}
         {\bigl|1-\kappa\,\Pi(\omega,k)\bigr|^{2}}\;
    \sum_{n=1}^{\infty}
    e^{-2\sigma\gamma_{n}^{2}}\;
    \pi\,\delta\bigl(\omega-\Omega_{n}(k)\bigr)
    + \mathcal{O}(\varepsilon^{4})\,.
\end{equation}

However, the physically relevant quantity is not
$\mathrm{Im}\,\kappa_{\mathrm{eff}}$ at a single frequency, but the
\emph{two-point function} of the noise that drives the stochastic
Einstein--Langevin equation.  The generalised FDT for the Zeta-Comb
reads:
\begin{equation}\label{eq:gen-FDT}
  \boxed{
  N(\omega,k)
  = \coth\!\biggl(\frac{\omega}{2T_{\mathrm{eff}}}\biggr)\,
    \mathrm{Im}\,D_{R}(\omega,k)
  + \varepsilon^{2}\!\!\!
    \sum_{\substack{n,m=1\\n\neq m}}^{\infty}
    \!\!\!
    e^{-\sigma(\gamma_{n}^{2}+\gamma_{m}^{2})}\;
    \mathcal{H}_{nm}(\omega,k)\,,}
\end{equation}
where $\mathcal{H}_{nm}$ is the \textbf{Hilbert interference term}
arising from the cross-correlation of the $n$-th and $m$-th
environment modes:
\begin{equation}\label{eq:Hnm}
  \mathcal{H}_{nm}(\omega,k)
  = \frac{1}{\pi}\,\mathcal{P}\!\!\int_{-\infty}^{\infty}
    \frac{\mathrm{Im}\bigl[
      \widetilde{G}_{n}(\omega',k)\,
      \widetilde{G}_{m}^{*}(\omega',k)
    \bigr]}{\omega'-\omega}\;d\omega'\,.
\end{equation}
Here $\widetilde{G}_{n}(\omega,k)$ is the retarded Green function of
the $n$-th environment mode.

\begin{remark}[Structure of the generalised FDT]
Equation~\eqref{eq:gen-FDT} has two terms:
\begin{enumerate}
  \item The \textbf{diagonal term} ($n=m$) reproduces the standard FDT
    at each Zeta-Comb frequency individually.
  \item The \textbf{off-diagonal terms} ($n\neq m$) are the
    \emph{beat-frequency} contributions.  They arise because the
    Zeta-Comb environment has a \emph{discrete, multi-line} spectral
    density---the noise at frequency~$\omega$ receives contributions
    from interference between lines at $\Omega_{n}$ and $\Omega_{m}$.
    These terms are suppressed by
    $e^{-\sigma(\gamma_{n}^{2}+\gamma_{m}^{2})}$ and are therefore
    sensitive to the pairwise \emph{differences}
    $\gamma_{n}-\gamma_{m}$.
\end{enumerate}
\end{remark}

%%---------------------------------------------------------------------
\subsubsection{Justification of the Hilbert-transform form}%
\label{sssec:Hilbert-justification}
%%---------------------------------------------------------------------

The appearance of the Hilbert transform in~\eqref{eq:Hnm} is not
incidental.  It follows from the same Kramers--Kronig logic used in
\S\ref{ssec:KK-proof}: because each $\widetilde{G}_{n}$ is a retarded
response function, the product
$\widetilde{G}_{n}\widetilde{G}_{m}^{*}$ is analytic in the upper
half-plane for fixed~$m$, and its real and imaginary parts are
therefore Hilbert-transform pairs.  The $\mathcal{H}_{nm}$ terms
extract the part of the noise that is \emph{coherent} between modes
$n$ and~$m$, which is precisely the beat-frequency content.

%%---------------------------------------------------------------------
\subsection{Beat-frequency spectrum and GUE
  fingerprint}\label{ssec:beat-spectrum}
%%---------------------------------------------------------------------

The off-diagonal terms in the generalised
FDT~\eqref{eq:gen-FDT} oscillate at the \textbf{beat frequencies}
\begin{equation}\label{eq:beat-freq}
  \omega_{nm} \equiv \gamma_{n}-\gamma_{m}\,,
  \qquad n>m\,.
\end{equation}
The number of distinct beat frequencies from $N_{\mathrm{eff}}$
effective zeros is $\binom{N_{\mathrm{eff}}}{2}$; for the sweet-spot
value $\sigma\sim 10^{-3}$, the first $\sim 4$ zeros contribute,
yielding $\sim 7$ pairwise beats.

%%---------------------------------------------------------------------
\subsubsection{Observable modulation in $\kappa_{\mathrm{eff}}$}%
\label{sssec:observable-modulation}
%%---------------------------------------------------------------------

Including the beat-frequency terms, the fractional shift of the
gravitational coupling acquires a characteristic oscillatory structure
in $\ln(k/k_{*})$:
\begin{equation}\label{eq:delta-kappa-full}
  \frac{\delta\kappa}{\kappa}(\omega,k)
  = \varepsilon^{2}\sum_{n}
    e^{-2\sigma\gamma_{n}^{2}}\;
    e^{i\gamma_{n}\ln(k/k_{*})}
  + \varepsilon^{2}\!\!\!
    \sum_{n\neq m}
    e^{-\sigma(\gamma_{n}^{2}+\gamma_{m}^{2})}\;
    e^{i\,\omega_{nm}\ln(k/k_{*})}
  + \mathcal{O}(\varepsilon^{4})\,.
\end{equation}
The first sum is the \textbf{fundamental Zeta-Comb}
(oscillations at each $\gamma_{n}$); the second is the
\textbf{beat-frequency comb} (oscillations at each
$\gamma_{n}-\gamma_{m}$).  The signal amplitude at the sweet spot is
\begin{equation}\label{eq:signal-amplitude}
  \left|\frac{\delta\kappa}{\kappa}\right|_{\mathrm{sweet\;spot}}
  = \varepsilon^{2}\,e^{-2\sigma\gamma_{1}^{2}}
  \approx 6.7\times 10^{-5}\,,
  \qquad
  (\varepsilon=0.01,\;\sigma=0.001)\,.
\end{equation}

%%---------------------------------------------------------------------
\subsubsection{The GUE cross-check}\label{sssec:GUE-crosscheck}
%%---------------------------------------------------------------------

The distribution of the normalised beat frequencies
\begin{equation}\label{eq:u-def}
  u_{nm}
  \equiv \frac{\gamma_{n}-\gamma_{m}}{\langle\Delta\gamma\rangle}\,,
  \qquad
  \langle\Delta\gamma\rangle
  = \frac{2\pi}{\ln(\gamma_{\mathrm{mid}}/2\pi)}\,,
\end{equation}
(where $\langle\Delta\gamma\rangle$ is the local mean spacing from
the Riemann--von~Mangoldt formula) is predicted by the
\textbf{Montgomery--Odlyzko pair correlation conjecture} to follow
the GUE distribution:
\begin{equation}\label{eq:R2-GUE}
  \boxed{
  R_{2}(u)
  = 1-\left(\frac{\sin\pi u}{\pi u}\right)^{2}\,.}
\end{equation}
This is \emph{the same} pair-correlation function that governs the
eigenvalue statistics of random matrices drawn from the Gaussian
Unitary Ensemble (GUE), as first noted by Dyson in a celebrated
conversation with Montgomery.

The GUE prediction is a \textbf{binary diagnostic} for the
Zeta-Comb hypothesis.  It makes two sharp predictions:
\begin{enumerate}
  \item \textbf{Level repulsion at $u\to 0$}: $R_{2}(u)\to 0$
    quadratically as $u\to 0$, meaning the smallest normalised
    beat-frequency spacings are suppressed.  For the first 10~zeros,
    the smallest normalised spacing is
    $u_{\mathrm{min}}\sim 0.2$--$0.5$.
    A Poisson (random, uncorrelated) environment would produce
    $R_{2}(u)\to 1$ as $u\to 0$---no level repulsion.

  \item \textbf{Oscillatory approach to unity}: $R_{2}(u)$ overshoots
    and oscillates around~1 for $u\gtrsim 1$, a distinctive
    ``ringing'' signature absent in Poisson statistics.
\end{enumerate}

If the beat frequencies extracted from a gravitational-wave or CMB
dataset match the GUE pattern~\eqref{eq:R2-GUE}, it constitutes
evidence for an arithmetic environment; if they match Poisson, the
Zeta-Comb hypothesis is falsified regardless of parameter tuning.

%%---------------------------------------------------------------------
\subsubsection{Connection to Rodgers' theorem and prime
  distributions}\label{sssec:rodgers}
%%---------------------------------------------------------------------

The diagnostic above acquires deeper significance through a theorem of
Rodgers, who proved (conditional on the Riemann
Hypothesis) that the GUE conjecture for zeta zeros is \emph{logically
  equivalent} to a specific statement about the distribution of
products of prime numbers in short intervals.  Concretely: the
beat spectrum of $\kappa_{\mathrm{eff}}$ traces GUE statistics
\emph{if and only if} the primes satisfy the Hardy--Littlewood-type
correlation laws encoded in the von~Mangoldt function.

This transforms an \emph{observational} measurement---the pattern of
oscillations in the gravitational-wave or CMB power spectrum---into a
\emph{number-theoretic} test.  A successful detection of the GUE
pattern would:
\begin{enumerate}
  \item Confirm the Zeta-Comb as the physical environment.
  \item Provide empirical evidence for the GUE conjecture
    (equivalently, for the Montgomery pair correlation conjecture).
  \item Via Rodgers' theorem, test a specific prediction about
    the distribution of prime powers.
\end{enumerate}

%%---------------------------------------------------------------------
\subsection{Beat-frequency visibility and the
  $N_{\mathrm{eff}}$ threshold}\label{ssec:beat-visibility}
%%---------------------------------------------------------------------

The beat-frequency terms in~\eqref{eq:delta-kappa-full} are
observable only if $N_{\mathrm{eff}}\geq 2$, i.e.\ at least two
zeros contribute with appreciable amplitude.  This requirement sets an
\emph{upper bound} on the resonance width:
\begin{equation}\label{eq:sigma-bound}
  \sigma
  \lesssim \frac{1}{2\gamma_{2}^{2}}
  \approx \frac{1}{2\times(21.02)^{2}}
  \approx 0.0011\,,
\end{equation}
or equivalently $\sigma\lesssim 0.002$ when we require
$N_{\mathrm{eff}}\geq 2$ to one-$e$-fold accuracy.

Simultaneously, the signal amplitude~\eqref{eq:signal-amplitude}
is maximised for small~$\sigma$ (since the exponential damping is
weakest).  This means that the requirement for beat-frequency
visibility \emph{overlaps} with the requirement for maximal
signal-to-noise in gravitational-wave detectors---a non-trivial
structural coincidence that we elevate to the status of
\textbf{theoretical naturalness} in
\S\ref{sec:parameter-space}.

\begin{remark}[Sweet spot]
The intersection of the beat-visibility constraint
($\sigma\lesssim 0.002$) with the Planck~2018 upper bound on
$\varepsilon$ and the BBO/DECIGO sensitivity curves defines a
\emph{sweet spot} at $\sigma\sim 10^{-3}$, $\varepsilon\sim 10^{-2}$.
At this point the signal amplitude is
$|\delta\kappa/\kappa|\approx 6.7\times 10^{-5}$, the effective
number of zeros is $N_{\mathrm{eff}}\approx 4.2$, and the GUE
cross-check is operational with 7~beat frequencies.
\end{remark}

%%---------------------------------------------------------------------
\subsection{Summary: from general theorem to arithmetic
  prediction}\label{ssec:sec3-summary}
%%---------------------------------------------------------------------

The logical flow of this section is:
\[
  \underbrace{\text{AL-GFT couplings}~g_{n}}_%
    {\text{Eq.~\eqref{eq:gn}}}
  \;\;\xrightarrow{\;\text{SK trace}\;}
  \;\;\underbrace{N(k)=N_{0}\,\mathcal{M}(k)}_%
    {\text{Zeta-Comb}}
  \;\;\xrightarrow{\;\text{gen.\ FDT}\;}
  \\\;\;\underbrace{\omega_{nm}=\gamma_{n}-\gamma_{m}}_%
    {\text{beat frequencies}}

\;\;\xrightarrow{\;\text{Montgomery}\;}
  \;\;\underbrace{R_{2}(u)\stackrel{?}{=}1-\mathrm{sinc}^{2}\!\pi u}_%
    {\text{GUE test}}\,.
\]

\noindent
The chain has a single free spectral-density choice
(Zeta-Comb vs.\ anything else); \emph{every subsequent prediction}
is fixed.  The GUE cross-check is binary and requires \emph{no
  fitting}: it tests the hypothesis at the level of the
environment's identity, not its parameters.


\section{Conservation and Consistency}
\label{sec:conservation}

The promotion of the gravitational coupling $\kappa^2 = 16\pi G_N$ from a real constant
to a complex, scale-dependent effective coupling $\kappa_{\mathrm{eff}}(k)$ raises an
immediate question of principle: does the modified theory remain consistent with the
diffeomorphism invariance that underpins general relativity?  In this section we show
that the answer is yes, provided the noise and dissipation kernels satisfy a set of
Ward--Takahashi identities inherited from the conservation of the stress-energy tensor.
We establish three results: (i) the stochastic source $\xi_{ab}$ is covariantly
conserved on the semiclassical background; (ii) the contracted Bianchi identity
$\nabla^\mu G_{\mu\nu} = 0$ is preserved at each order in $\varepsilon$; and (iii) the
energy exchanged between the gravitational ``system'' and the Zeta-Comb environment is
governed by a fluctuation-dissipation relation whose structure is dictated by causality
and unitarity.

%======================================================================
\subsection{Ward identity for the stochastic source}
\label{subsec:ward}
%======================================================================

The starting point is the operator identity for the renormalized stress-energy tensor
on the semiclassical background $({\mathcal M}, g_{ab})$:
%
\begin{equation}
  \nabla^a \hat{T}^{\mathrm{R}}_{ab}[g;\, x) = 0\,.
  \label{eq:Tabconserved}
\end{equation}
%
This is the quantum Ward identity expressing diffeomorphism invariance of the matter
sector.  The noise kernel, defined as the connected symmetrized two-point function of
$\hat{T}^{\mathrm{R}}_{ab}$,
%
\begin{equation}
  N_{abcd}(x,y)
  \;=\;
  \tfrac{1}{2}\,
  \bigl\langle
    \bigl\{\hat{t}_{ab}(x),\, \hat{t}_{cd}(y)\bigr\}
  \bigr\rangle\,,
  \qquad
  \hat{t}_{ab} \;\equiv\; \hat{T}^{\mathrm{R}}_{ab}
    - \langle \hat{T}^{\mathrm{R}}_{ab}\rangle\,,
  \label{eq:noisekernel}
\end{equation}
%
inherits this conservation as a \emph{distributional} identity on each index group
separately~\cite{HuVerdaguer2008}:
%
\begin{equation}
  \nabla^a_x\, N_{abcd}(x,y) = 0\,,
  \qquad
  \nabla^c_y\, N_{abcd}(x,y) = 0\,.
  \label{eq:Nconserved}
\end{equation}

\noindent
\textbf{Proof.}
By Eq.~\eqref{eq:Tabconserved} the operator $\hat{t}_{ab}$ satisfies
$\nabla^a \hat{t}_{ab} = 0$ in the Heisenberg picture.
Taking the $x$-divergence of the anticommutator inside the expectation
value in Eq.~\eqref{eq:noisekernel} and using the Leibniz rule for
$\nabla^a_x$---which acts only on the $x$-supported operator---yields
Eq.~\eqref{eq:Nconserved} directly.  No renormalization subtlety arises
because the noise kernel is ultraviolet finite: the identity-operator
counterterms in $\hat{T}^{\mathrm{R}}_{ab} - \hat{T}_{ab}$ cancel in the
connected correlator~\cite{PhillipsHu2001,HuVerdaguer2004}.
\hfill$\square$

\medskip

The Gaussian stochastic tensor $\xi_{ab}[g;\,x)$ is defined by the correlators
%
\begin{equation}
  \langle \xi_{ab}(x)\rangle_s = 0\,,
  \qquad
  \langle \xi_{ab}(x)\,\xi_{cd}(y)\rangle_s
  = N_{abcd}(x,y)\,,
  \label{eq:xicorrelators}
\end{equation}
%
where $\langle\cdot\rangle_s$ denotes the stochastic average.
Since $N_{abcd}$ is symmetric and positive semi-definite
(a consequence of the self-adjointness of $\hat{T}^{\mathrm{R}}_{ab}$),
$\xi_{ab}$ is well-defined.
Taking the divergence of Eq.~\eqref{eq:xicorrelators} and using
Eq.~\eqref{eq:Nconserved}:
%
\begin{equation}
  \langle \nabla^a \xi_{ab}\rangle_s = 0\,,
  \qquad
  \langle \nabla^a_x \xi_{ab}(x)\;\xi_{cd}(y)\rangle_s = 0\,.
  \label{eq:xiconserved}
\end{equation}
%
For a Gaussian field whose mean and two-point function are both annihilated by
$\nabla^a$, all higher cumulants vanish, and hence
$\nabla^a \xi_{ab} = 0$ \emph{with probability one}---it is a deterministic
zero, not merely zero in expectation.

\begin{tcolorbox}[colback=gray!5, colframe=black!50, title={Theorem 1 (Stochastic conservation)}]
Let $\xi_{ab}$ be the Gaussian stochastic source defined by
Eq.~\eqref{eq:xicorrelators} on a semiclassical background
$({\mathcal M}, g_{ab})$ satisfying the semiclassical Einstein equation.
Then $\nabla^a \xi_{ab}[g;\,x) = 0$ as a distributional identity on
$\mathcal{M}$, and the Einstein--Langevin equation
\begin{equation}
  G_{ab}[g+h]
  + \Lambda\, g_{ab}
  - 8\pi G_N\,
  \bigl\langle \hat{T}^{\mathrm{R}}_{ab}[g+h]\bigr\rangle^{(1)}
  = 8\pi G_N\, \xi_{ab}[g]
  \label{eq:EinsteinLangevin}
\end{equation}
is consistent with the contracted Bianchi identity to linear order
in $h_{ab}$.
\end{tcolorbox}

\noindent
This result is standard in the Hu--Verdaguer stochastic gravity
programme~\cite{HuVerdaguer2004,HuVerdaguer2008,CalzettaHu2008}.
What is \emph{new} in our framework is its extension to the case where
the environment carries arithmetic structure (the Zeta-Comb), and the
resulting effective coupling $\kappa_{\mathrm{eff}}(k)$ is complex.
We now show that the Ward identity survives this generalization.

%======================================================================
\subsection{Preservation of the Bianchi identity with complex $\kappa_{\mathrm{eff}}$}
\label{subsec:bianchi}
%======================================================================

Recall from \S\ref{sec:complexkappa} that the effective gravitational
coupling takes the form
%
\begin{equation}
  \kappa_{\mathrm{eff}}(k)
  \;=\;
  \frac{\kappa}{1 - \kappa\,
    \widetilde{D}_{\!R}(k)\big/\mathcal{O}(k)}\,,
  \label{eq:kappaeff}
\end{equation}
%
where $\widetilde{D}_{\!R}(k)$ is the retarded dissipation kernel and
$\mathcal{O}(k)$ encodes the system's response (both defined in
momentum space on the FLRW background).  The imaginary part of
$\kappa_{\mathrm{eff}}$ arises because $\widetilde{D}_{\!R}(k)$ is
complex for causal (retarded) kernels---its real and imaginary parts
are related by Kramers--Kronig dispersion relations.

The key observation is that Eq.~\eqref{eq:kappaeff} is a
\emph{multiplicative} dressing of the bare coupling $\kappa$.
The linearized Einstein tensor on the left-hand side of the
Einstein--Langevin equation still satisfies
%
\begin{equation}
  \nabla^\mu G^{(1)}_{\mu\nu}[g;\, h] = 0
  \label{eq:BianchiLinearized}
\end{equation}
%
identically, as this is a consequence of the contracted Bianchi
identity applied to the background-plus-perturbation metric.
Equation~\eqref{eq:BianchiLinearized} holds for \emph{any} $h_{ab}$,
independently of the field equations; it is a geometric identity,
not a dynamical one.

On the source side, the complex dressing enters as
%
\begin{equation}
  \kappa_{\mathrm{eff}}(k)\, \xi_{ab}(k)
  \;=\;
  \kappa\,\xi_{ab}(k)
  \;+\;
  \kappa^2\,
  \frac{\widetilde{D}_{\!R}(k)}{\mathcal{O}(k)}\,
  \xi_{ab}(k)
  \;+\;
  \mathcal{O}(\kappa^3)\,.
  \label{eq:dressedxi}
\end{equation}
%
At each order in $\kappa$, the stochastic source remains divergence-free
because:
%
\begin{enumerate}
  \item $\nabla^a \xi_{ab} = 0$ by Theorem~1;
  \item The dressing factor
    $\widetilde{D}_{\!R}(k)/\mathcal{O}(k)$ is a \emph{scalar}
    function of the comoving wavenumber $k$ (not a differential
    operator acting on the tensor indices of $\xi_{ab}$), so
    it commutes with the covariant divergence:
    %
    \begin{equation}
      \nabla^a\!
      \Bigl[
        f(k)\,\xi_{ab}(k)
      \Bigr]
      =
      f(k)\,\nabla^a \xi_{ab}(k)
      = 0\,,
      \label{eq:scalarcommute}
    \end{equation}
    %
    where $f(k) = \widetilde{D}_{\!R}(k)/\mathcal{O}(k)$.
\end{enumerate}
%
Thus, promoting $\kappa \to \kappa_{\mathrm{eff}}(k)$ does not
spoil the divergence-free character of the source, and the Bianchi
identity is preserved order by order.

\begin{tcolorbox}[colback=blue!3, colframe=blue!40, title={Corollary (Bianchi compatibility)}]
The complex gravitational coupling $\kappa_{\mathrm{eff}}(k)$ defined
by Eq.~\eqref{eq:kappaeff} is compatible with the contracted Bianchi
identity $\nabla^\mu G_{\mu\nu} = 0$ at every order in the perturbative
expansion in $\varepsilon$ (the multiplicity coupling strength), provided
the bare noise kernel satisfies Eq.~\eqref{eq:Nconserved}.
\end{tcolorbox}

\noindent
\textbf{Remark.}
The argument above relies on the fact that the Zeta-Comb modulation
enters through a \emph{scalar} dressing of the coupling, not through a
tensor-valued modification of the noise kernel's index structure.
This is a consequence of the linear system-environment coupling adopted
in the AL-GFT framework (\S\ref{sec:ALGFT}): the interaction
$S_{\mathrm{int}} \sim \int \phi\, \mathcal{O}$ preserves the
tensorial conservation law because $\mathcal{O}$ is constructed from
scalar contractions of the environment fields.  A nonlinear coupling
(relevant for beyond-Gate-1 extensions) could in principle introduce
tensor dressings, which would require a more careful Ward identity
analysis.

%======================================================================
\subsection{Fluctuation-dissipation relation and energy flow}
\label{subsec:FDR}
%======================================================================

An open quantum system in thermal equilibrium satisfies the
fluctuation-dissipation relation (FDR), which constrains the noise
kernel $N$ in terms of the dissipation kernel $D_R$ and the
temperature $T$ of the environment~\cite{CalzettaHu2008,HuVerdaguer2004}.
In the cosmological setting, the relevant ``temperature'' is the
Gibbons--Hawking temperature $T_{\mathrm{dS}} = H/(2\pi)$ of the
quasi-de~Sitter background.  For a causal (retarded) dissipation kernel,
the FDR in Fourier space reads
%
\begin{equation}
  \widetilde{N}(\omega)
  \;=\;
  \coth\!\Bigl(\frac{\omega}{2\,T_{\mathrm{dS}}}\Bigr)\,
  \mathrm{Im}\,\widetilde{D}_{\!R}(\omega)\,,
  \label{eq:FDRbasic}
\end{equation}
%
where $\omega$ is the frequency conjugate to conformal time.
At high frequencies $\omega \gg T_{\mathrm{dS}}$, the hyperbolic
cotangent approaches unity and $\widetilde{N} \approx
\mathrm{Im}\,\widetilde{D}_{\!R}$, corresponding to the quantum
(zero-temperature) limit.  At low frequencies
$\omega \ll T_{\mathrm{dS}}$, the classical limit
$\widetilde{N} \approx (2\,T_{\mathrm{dS}}/\omega)\,
\mathrm{Im}\,\widetilde{D}_{\!R}$ emerges.

\paragraph{Zeta-Comb generalization.}
In the AL-GFT framework, the environment is not a generic thermal bath
but a structured collection of modes at frequencies set by the Riemann
zeta zeros $\{\gamma_n\}$.  The noise kernel acquires the Zeta-Comb
modulation (derived in \S\ref{sec:zetacomb}):
%
\begin{equation}
  \widetilde{N}(k)
  \;=\;
  \widetilde{N}_0(k)\,
  \biggl[
    1 + 2\varepsilon^2
    \sum_{n=1}^{N_{\mathrm{eff}}}
    e^{-2\sigma\gamma_n^2}\,
    \cos\!\bigl(\gamma_n \ln(k/k_\star) + \varphi_n\bigr)
  \biggr]\,,
  \label{eq:NZetaComb}
\end{equation}
%
where $\widetilde{N}_0(k)$ is the standard (unmodulated) noise kernel,
$\varepsilon$ is the multiplicity coupling, $\sigma$ the resonance width,
$\gamma_n = \mathrm{Im}(\rho_n)$ the $n$-th Riemann zero's imaginary part,
and $\varphi_n = \tfrac{1}{2}\arg\zeta(\rho_n)$.

The dissipation kernel is constrained by the Kramers--Kronig relations,
which are a direct consequence of causality (analyticity of
$\widetilde{D}_{\!R}$ in the upper half of the complex $\omega$-plane):
%
\begin{align}
  \mathrm{Re}\,\widetilde{D}_{\!R}(\omega)
  &= \frac{1}{\pi}\,\mathcal{P}\!\int_{-\infty}^{\infty}
     \frac{\mathrm{Im}\,\widetilde{D}_{\!R}(\omega')}
          {\omega' - \omega}\, d\omega'\,,
  \label{eq:KK1}
  \\[4pt]
  \mathrm{Im}\,\widetilde{D}_{\!R}(\omega)
  &= -\frac{1}{\pi}\,\mathcal{P}\!\int_{-\infty}^{\infty}
     \frac{\mathrm{Re}\,\widetilde{D}_{\!R}(\omega')}
          {\omega' - \omega}\, d\omega'\,.
  \label{eq:KK2}
\end{align}

\noindent
The generalized FDR for the Zeta-Comb environment then takes the form
%
\begin{equation}
  \boxed{
  \widetilde{N}(k)
  \;=\;
  \coth\!\Bigl(\frac{\omega_k}{2\,T_{\mathrm{dS}}}\Bigr)\,
  \mathrm{Im}\,\widetilde{D}_{\!R}(k)
  \;+\;
  \varepsilon^2\,
  \sum_{n}\,
  e^{-2\sigma\gamma_n^2}\;
  \mathcal{H}\!\bigl[\widetilde{N}_0\bigr](k;\,\gamma_n)\,,
  }
  \label{eq:FDRgeneralized}
\end{equation}
%
where $\omega_k$ is the physical frequency associated with wavenumber $k$,
and $\mathcal{H}[\widetilde{N}_0](k;\,\gamma_n)$ denotes the
Hilbert-transform--shifted contribution at beat frequency $\gamma_n$.
The first term is the standard FDR; the second encodes the
\emph{arithmetic correction} from the Zeta-Comb environment.

\paragraph{Energy budget.}
The imaginary part of $\kappa_{\mathrm{eff}}$ governs the energy exchange
between the gravitational system and the arithmetic environment.
Writing $\kappa_{\mathrm{eff}} = |\kappa_{\mathrm{eff}}|\,
e^{i\theta(k)}$, the energy flow rate per logarithmic wavenumber
interval is
%
\begin{equation}
  \frac{d\mathcal{E}}{d\ln k}
  \;=\;
  -\,\mathrm{Im}\,\kappa_{\mathrm{eff}}(k)\;
  \bigl\langle \xi_{ab}(k)\,h^{ab}(k) \bigr\rangle_s
  \;\propto\;
  \varepsilon^2\,
  \sin\theta(k)\,,
  \label{eq:energyflow}
\end{equation}
%
where the phase $\theta(k) = \arg\kappa_{\mathrm{eff}}(k)$ is determined
by the ratio $\mathrm{Im}\,\widetilde{D}_{\!R}/
\mathrm{Re}\,(1 - \kappa\,\widetilde{D}_{\!R}/\mathcal{O})$.
At the sweet-spot parameters $(\sigma \sim 10^{-3},\,
\varepsilon \sim 10^{-2})$ identified in \S\ref{sec:parameterspace},
the phase is small, $\theta \sim \mathcal{O}(\varepsilon^2)$, so the
energy exchange is perturbative and the gravitational sector remains
stable against runaway energy injection from the environment.

\paragraph{Physical interpretation.}
In the language of open quantum systems, the imaginary part
$\mathrm{Im}\,\kappa_{\mathrm{eff}}$ plays the role of a
\emph{gravitational dissipation coefficient}: it quantifies how
efficiently the Zeta-Comb environment can absorb or emit energy
through metric fluctuations.  The Ward identity
(Theorem~1) guarantees that this dissipation respects the
diffeomorphism symmetry---no energy is created or destroyed in the
\emph{total} system (gravity + environment), only redistributed
between the sectors.

This is the gravitational analogue of the optical theorem in
scattering theory: the imaginary part of the forward scattering
amplitude is related to the total cross section.  Here, the imaginary
part of $\kappa_{\mathrm{eff}}$ is related to the total ``cross section''
for graviton--environment interaction, summed over the discrete
Zeta-Comb modes:
%
\begin{equation}
  \mathrm{Im}\,\kappa_{\mathrm{eff}}(k)
  \;=\;
  \kappa^2\,
  \frac{\mathrm{Im}\,\widetilde{D}_{\!R}(k)}{|\mathcal{O}(k)|^2}
  \;+\;
  \mathcal{O}(\kappa^3)\,.
  \label{eq:Imkappa_crosssection}
\end{equation}

%======================================================================
\subsection{Consistency checks}
\label{subsec:checks}
%======================================================================

We summarize the internal consistency of the framework by listing the
constraints that are satisfied simultaneously:

\begin{enumerate}
  \item \textbf{Covariant conservation.}\;
    $\nabla^\mu G_{\mu\nu} = 0$ (geometric, from Bianchi identity);
    $\nabla^a \xi_{ab} = 0$ (dynamical, from Ward identity of the matter
    sector).  Together these guarantee that the perturbed metric
    $g_{ab} + h_{ab}$ satisfies the linearized constraint equations.

  \item \textbf{Causality.}\;
    The retarded dissipation kernel $\widetilde{D}_{\!R}(\omega)$ is
    analytic in the upper half-plane, ensuring that
    $\kappa_{\mathrm{eff}}(\omega)$ also inherits this analyticity
    (since the denominator in Eq.~\eqref{eq:kappaeff} is a polynomial
    in $\widetilde{D}_{\!R}$ and thus shares its domain of analyticity).
    Kramers--Kronig relations~\eqref{eq:KK1}--\eqref{eq:KK2} follow.

  \item \textbf{Unitarity.}\;
    The noise kernel $N_{abcd}$ is positive semi-definite
    (consequence of the self-adjointness of $\hat{T}^{\mathrm{R}}_{ab}$),
    which guarantees that the stochastic source has a well-defined
    probability measure and that the fluctuation-dissipation
    relation~\eqref{eq:FDRbasic} is compatible with the second law of
    thermodynamics.

  \item \textbf{Trace identity for conformal matter.}\;
    If the quantum field is conformally coupled, the trace
    $g^{ab}\xi_{ab} = 0$, so the stochastic source does not contribute
    to the trace anomaly.  This is preserved under
    $\kappa \to \kappa_{\mathrm{eff}}$ since the dressing is scalar.

  \item \textbf{Perturbative stability.}\;
    At the sweet-spot parameters, the correction to $\kappa$ is
    $|\delta\kappa/\kappa| \sim 6.7 \times 10^{-5}$, well within the
    perturbative regime $|\delta\kappa/\kappa| \ll 1$.  The backreaction
    of the Zeta-Comb environment on the background geometry is therefore
    self-consistently small.

  \item \textbf{$\varepsilon \to 0$ recovery.}\;
    In the decoupling limit $\varepsilon \to 0$, the noise kernel
    reduces to $\widetilde{N}_0(k)$, the dressing
    $\kappa_{\mathrm{eff}} \to \kappa$, and the standard semiclassical
    Einstein equation is recovered to arbitrary precision.
\end{enumerate}

\noindent
These six conditions collectively ensure that the complex gravitational
coupling introduced in \S\ref{sec:complexkappa} is not an \emph{ad hoc}
modification of general relativity, but a necessary consequence of
treating gravity as an open quantum system whose environment carries
arithmetic structure.  The Ward identity is the linchpin: it ties
together conservation, causality, and the Bianchi identity into a single
self-consistent package.


%———————————————————————————————————————————
% §5  The GUE Cross-Check
%———————————————————————————————————————————

\section{The GUE Cross-Check}\label{sec:gue}

The central claim of this paper is that the environment
generating the imaginary part of the gravitational coupling
is not generic stochastic noise but is \emph{arithmetically structured}:
its spectral lines are located at the imaginary parts
$\gamma_n$ of the non-trivial Riemann zeta zeros.
Any such claim demands an internal consistency test that is
independent of the overall signal amplitude~$\varepsilon$.
The Gaussian Unitary Ensemble (GUE) statistics of the
$\gamma_n$~\cite{Montgomery1973,Odlyzko1987}
provide exactly such a test.
In this section we show that the \emph{beat spectrum}
of the complex coupling $\kappa_{\mathrm{eff}}(k)$ carries
a pair-correlation fingerprint that distinguishes a
Riemann-zero environment from any environment with
uncorrelated (Poisson) spectral lines.
The test is binary---it requires no fit parameters---and
can be applied to any future dataset (CMB, SGWB, or
interferometric) in which the oscillatory modulation of
Sec.~\ref{sec:zetacomb} is resolved.


%-----------------------------------------------
\subsection{Pair correlation of the beat frequencies}
\label{sec:paircorr}

\subsubsection{From spectral lines to beat frequencies}

The modulation function~$\mathcal{M}(k)$ derived in
Sec.~\ref{sec:zetacomb} is a superposition of
$N_{\mathrm{eff}}$ oscillatory terms with frequencies
$\gamma_n\,(n=1,\dots,N_{\mathrm{eff}})$ in the variable
$\ln(k/k_\star)$.  The effective number of contributing zeros
is set by the Gaussian damping:
\begin{equation}\label{eq:Neff}
  N_{\mathrm{eff}}(\sigma)
  \;=\;
  \sum_{n=1}^{\infty} e^{-2\sigma\,\gamma_n^2}\,,
\end{equation}
which, for the sweet-spot value $\sigma\sim 10^{-3}$,
gives $N_{\mathrm{eff}}\approx 4.2$.
The power spectrum of~$\mathcal{M}(k)$ in the
$u\equiv\ln(k/k_\star)$ domain therefore contains
\begin{equation}\label{eq:Nbeat}
  N_{\mathrm{beat}}
  \;=\;
  \binom{\lfloor N_{\mathrm{eff}}\rfloor}{2}
\end{equation}
pairwise beat frequencies
\begin{equation}\label{eq:deltaij}
  \Delta_{ij} \;=\; |\gamma_i - \gamma_j|\,,
  \qquad 1\le i < j \le N_{\mathrm{eff}}\,.
\end{equation}
At the sweet spot, $\lfloor N_{\mathrm{eff}}\rfloor = 4$
yields $N_{\mathrm{beat}}=6$ resolved beat lines
(with a seventh partially suppressed),
each carrying information about the \emph{spacing}
between pairs of Riemann zeros.

\subsubsection{The Montgomery--Odlyzko pair correlation}

Montgomery's pair correlation conjecture~\cite{Montgomery1973},
confirmed numerically by Odlyzko to extraordinary precision
for zeros near height $T\sim 10^{20}$~\cite{Odlyzko1987,Odlyzko2001},
states that the two-point correlation function of the
normalized zero spacings $\hat\gamma_n = \gamma_n\cdot
\frac{\ln(\gamma_n/2\pi)}{2\pi}$
tends, in the large-height limit, to the GUE form:
\begin{equation}\label{eq:R2GUE}
  R_2^{\mathrm{GUE}}(s)
  \;=\;
  1 - \left(\frac{\sin\pi s}{\pi s}\right)^{\!2}\,.
\end{equation}
Two properties of Eq.~\eqref{eq:R2GUE} are decisive:
\begin{enumerate}
\item \textbf{Level repulsion.}\;
  $R_2^{\mathrm{GUE}}(s)\to 0$ as $s\to 0$,
  meaning that close spacings are quadratically suppressed:
  $R_2^{\mathrm{GUE}}(s) \sim \pi^2 s^2/3 + \mathcal{O}(s^4)$.
  This is in sharp contrast to Poisson statistics,
  for which $R_2^{\mathrm{Poisson}}(s) = 1$ for all~$s$.

\item \textbf{Oscillatory approach to unity.}\;
  For $s\gg 1$, $R_2^{\mathrm{GUE}}(s)=1-(\pi s)^{-2}\sin^2(\pi s)$
  oscillates with decaying amplitude, generating a
  characteristic ``dip--peak'' pattern in the
  pair-correlation histogram.
\end{enumerate}

\subsubsection{Beat-frequency pair correlation as a proxy}

The beat frequencies $\Delta_{ij}$ of Eq.~\eqref{eq:deltaij}
are, by definition, the \emph{differences} between pairs of
spectral lines.  Their empirical distribution, measured
from the Fourier transform of the observed modulation, is a
direct estimator of~$R_2(s)$.
Concretely, define the normalized differences
\begin{equation}\label{eq:snorm}
  s_{ij}
  \;=\;
  \frac{\Delta_{ij}}{\langle\Delta\rangle}\,,
\end{equation}
where $\langle\Delta\rangle$ is the mean spacing among the
first $N_{\mathrm{eff}}$ zeros.
For the first $\sim\!10$ Riemann zeros, the mean spacing is
$\langle\Delta\rangle\approx 3.7$ (in unnormalized units),
and the smallest normalized difference satisfies
$s_{\min}\approx 0.2$--$0.5$, reflecting GUE level repulsion.
A Poisson environment with the same mean density would
populate $s<0.1$ with probability $\sim 10\%$;
a GUE environment suppresses this to $\lesssim 1\%$.


%-----------------------------------------------
\subsection{Level repulsion as a binary diagnostic}
\label{sec:levelrepulsion}

\subsubsection{The diagnostic statistic}

We define a binary diagnostic that requires
no fitting and no knowledge of the signal amplitude~$\varepsilon$:
\begin{equation}\label{eq:Qstat}
  \mathcal{Q}
  \;\equiv\;
  \#\!\bigl\{(i,j):\; s_{ij} < s_{\mathrm{cut}}\bigr\}
  \;\Big/\;
  N_{\mathrm{beat}}\,,
\end{equation}
where $s_{\mathrm{cut}}$ is a fixed threshold
(we adopt $s_{\mathrm{cut}} = 0.1$).
The null hypothesis $H_0$ (Poisson environment)
and alternative $H_1$ (GUE/Riemann-zero environment)
predict:
\begin{center}
\renewcommand{\arraystretch}{1.3}
\begin{tabular}{lcc}
\hline
\textbf{Hypothesis} & $\mathbf{E}[\mathcal{Q}]$ & \textbf{$\mathrm{Var}[\mathcal{Q}]$}\\
\hline
$H_0$: Poisson   & $s_{\mathrm{cut}} = 0.10$
                  & $\sim s_{\mathrm{cut}}(1-s_{\mathrm{cut}})/N_{\mathrm{beat}}$ \\
$H_1$: GUE       & $\displaystyle\int_0^{s_{\mathrm{cut}}}
                     R_2^{\mathrm{GUE}}(s)\,ds
                     \;\approx\; 3\times 10^{-4}$
                  & $\ll \mathrm{Var}_{H_0}$ \\
\hline
\end{tabular}
\end{center}
In words: \emph{any observation of $\mathcal{Q}\approx 0$
is consistent with GUE; $\mathcal{Q}\gtrsim 0.05$ rules it out.}
The test is binary and requires only the extracted
beat frequencies, not their amplitudes.

\subsubsection{Finite-$N_{\mathrm{eff}}$ corrections}

For $N_{\mathrm{eff}} \lesssim 10$, the asymptotic
formula~\eqref{eq:R2GUE} receives corrections of order
$1/N_{\mathrm{eff}}^2$.  Following the finite-$N$ analysis
of Bogomolny and Keating~\cite{BogomolnyKeating1996},
the two-point function for the first~$N$ Riemann zeros
takes the form
\begin{equation}\label{eq:R2finite}
  R_2^{(N)}(s)
  \;=\;
  R_2^{\mathrm{GUE}}(s)
  \;-\;
  \frac{\gamma_0^2 + 2\gamma_1 + c_0}
       {\pi^2\,\bar\rho^{\,2}}\,
  \sin^2(\pi s)
  \;+\;
  \mathcal{O}(\bar\rho^{\,-3})\,,
\end{equation}
where $\bar\rho = \bar\rho(T)$ is the mean density of zeros
at height~$T$, and $\gamma_0$, $\gamma_1$, $c_0$ are
constants from the Hardy--Littlewood prime-pair
conjecture~\cite{BogomolnyKeating1996,Bogomolny2006}.
Crucially, the correction term is
\emph{proportional to $\sin^2(\pi s)$} and therefore
\textbf{vanishes at $s=0$}: level repulsion is preserved
at all finite~$N$.  This ensures that the binary
diagnostic~$\mathcal{Q}$ remains valid even at
$N_{\mathrm{eff}}\sim 4$.

\subsubsection{Nearest-neighbor spacing distribution}

A complementary statistic is the nearest-neighbor spacing
distribution $p(s)$, which for GUE is well approximated by
the Wigner surmise~\cite{Mehta2004}:
\begin{equation}\label{eq:pGUE}
  p_{\mathrm{GUE}}(s)
  \;=\;
  \frac{32}{\pi^2}\,s^2\,
  \exp\!\Bigl(-\frac{4}{\pi}\,s^2\Bigr)\,.
\end{equation}
The quadratic onset $p_{\mathrm{GUE}}(s)\propto s^2$ is the
hallmark of GUE-class level repulsion
($\beta=2$ in the Dyson classification), as opposed to
$p_{\mathrm{GOE}}(s)\propto s$ ($\beta=1$) or
$p_{\mathrm{Poisson}}(s)=e^{-s}$ ($\beta=0$).
With only $N_{\mathrm{eff}}\sim 4$ zeros, the histogram
of nearest-neighbor spacings cannot resolve the full
functional form of Eq.~\eqref{eq:pGUE}, but the
\emph{absence of small spacings} ($s<0.1$) is already
a $2$--$3\sigma$ discriminant between GUE and Poisson
for $N_{\mathrm{beat}}\ge 6$.


%-----------------------------------------------
\subsection{Connection to arithmetic via Rodgers' theorem}
\label{sec:rodgers}

The pair-correlation test acquires deeper significance
through a recent theorem of Rodgers~\cite{Rodgers2018},
which establishes:

\begin{theorem}[Rodgers, 2018]
Assume the Riemann Hypothesis.  Then the GUE Conjecture
for the zeros of $\zeta(s)$ is logically equivalent to
a statement about the distribution of products of primes:
specifically, the likelihood that products of primes
drawn from certain intervals are close to other products
of primes.
\end{theorem}

\noindent
This equivalence transforms our observational program
from a spectral-statistics exercise into a
\emph{test of prime-number distribution}:
\begin{itemize}
\item Measuring $R_2(s)$ from the beat spectrum
  of $\kappa_{\mathrm{eff}}(k)$ and finding it consistent
  with Eq.~\eqref{eq:R2GUE} would constitute
  \textbf{observational evidence for the GUE conjecture},
  obtained not from numerical computation of $\zeta(s)$
  but from a physical measurement.

\item Conversely, finding $R_2(s)$
  \emph{inconsistent} with GUE---for instance,
  Poisson statistics---would rule out the
  Riemann-zero identification of the spectral lines
  and falsify the AL-GFT environment model.
\end{itemize}

\noindent
It is worth emphasizing that Rodgers' theorem is conditional
on RH.  In our framework, this conditionality is absorbed
into the model definition: the AL-GFT environment is
\emph{constructed} from the non-trivial zeros
$\rho_n = \tfrac{1}{2}+i\gamma_n$, which lie on the
critical line by assumption.
A successful GUE cross-check therefore tests the
\emph{internal consistency} of this construction,
not the truth of RH per se.


%-----------------------------------------------
\subsection{Analysis pipeline}
\label{sec:pipeline}

We now describe a concrete, end-to-end pipeline for
extracting the GUE signature from observational data.
The pipeline is designed to be detector-agnostic; we
specify the required inputs and outputs at each stage.

\subsubsection{Stage 1: Spectral extraction}

\paragraph{Input.}
A one-dimensional data vector $d(k)$ representing
the measured quantity (CMB power spectrum $\mathcal{P}(k)$,
or stochastic gravitational-wave background spectral
density $\Omega_{\mathrm{GW}}(f)$) over a range
$[k_{\min},\,k_{\max}]$ with $k_{\max}/k_{\min}\gg 1$
(at least two decades for Fourier resolution).

\paragraph{Procedure.}
\begin{enumerate}
\item Divide out the smooth baseline
  (e.g., $\mathcal{P}_0(k)$ for CMB,
  or the power-law template $\Omega_{\mathrm{PL}}(f)$
  for SGWB) to obtain the residual modulation:
  \begin{equation}\label{eq:residual}
    r(u) \;\equiv\;
    \frac{d(k)}{d_{\mathrm{baseline}}(k)} - 1\,,
    \qquad u = \ln(k/k_\star)\,.
  \end{equation}

\item Compute the Lomb--Scargle periodogram
  $\hat{S}(\omega)$ of $r(u)$ in the log-wavenumber
  domain~$u$.  Peaks in $\hat{S}(\omega)$ at frequencies
  $\omega = \gamma_n$ correspond to individual Riemann-zero
  spectral lines; peaks at $\omega = \Delta_{ij}$
  correspond to beat frequencies.

\item Identify significant peaks above a
  false-alarm probability threshold
  (e.g., FAP $<10^{-3}$, corresponding to
  $\gtrsim 3.3\sigma$ for Gaussian noise).
  Record the set of detected frequencies
  $\{\hat\omega_\alpha\}_{\alpha=1}^{N_{\mathrm{det}}}$.
\end{enumerate}

\subsubsection{Stage 2: Frequency identification}

\paragraph{Procedure.}
\begin{enumerate}
\item Match detected frequencies $\hat\omega_\alpha$
  against the known Riemann zeros
  $\{\gamma_n\}_{n=1}^{N_{\max}}$
  (tabulated to arbitrary precision~\cite{LMFDB}).
  Accept matches satisfying
  $|\hat\omega_\alpha - \gamma_n| <
  \delta\omega_{\mathrm{res}}$,
  where $\delta\omega_{\mathrm{res}}
  \sim \pi/\Delta u$ is the Fourier
  resolution set by the available
  $u$-range $\Delta u = \ln(k_{\max}/k_{\min})$.

\item Form the set of identified zeros
  $\hat{\mathcal{Z}}
  = \{\hat\gamma_{n_1},\dots,\hat\gamma_{n_m}\}$
  and the corresponding set of all pairwise differences
  $\hat{\Delta}_{ij}
  = |\hat\gamma_{n_i} - \hat\gamma_{n_j}|$.

\item Cross-check: verify that beats
  $\hat\Delta_{ij}$ also appear in the
  periodogram~$\hat{S}(\omega)$.
  Concordance between directly detected lines and
  beat-predicted lines strengthens the identification.
\end{enumerate}

\subsubsection{Stage 3: Pair-correlation histogram}

\paragraph{Procedure.}
\begin{enumerate}
\item Normalize the identified spacings:
  $\hat{s}_{ij} = \hat\Delta_{ij}/\langle\hat\Delta\rangle$.

\item Construct the empirical pair-correlation function
  $\hat{R}_2(s)$ by histogramming the $\hat{s}_{ij}$
  with bin width
  $\delta s \sim 1/\sqrt{N_{\mathrm{beat}}}$.

\item Compare $\hat{R}_2(s)$ against
  $R_2^{\mathrm{GUE}}(s)$ [Eq.~\eqref{eq:R2GUE}]
  and $R_2^{\mathrm{Poisson}}(s) = 1$ using a
  Kolmogorov--Smirnov (KS) or Anderson--Darling (AD)
  statistic.

\item Evaluate the binary diagnostic $\mathcal{Q}$
  [Eq.~\eqref{eq:Qstat}] and report:
  \begin{equation}
    \begin{cases}
      \mathcal{Q} < s_{\mathrm{cut}} &
      \Longrightarrow\;\text{consistent with GUE
        (Riemann-zero environment),} \\[4pt]
      \mathcal{Q} \ge s_{\mathrm{cut}} &
      \Longrightarrow\;\text{inconsistent with GUE
        (generic/Poisson environment).}
    \end{cases}
  \end{equation}
\end{enumerate}

\subsubsection{Stage 4: Bootstrap confidence assessment}

\paragraph{Procedure.}
To assess statistical significance with
$N_{\mathrm{eff}}\sim 4$--$10$ zeros:
\begin{enumerate}
\item Generate $10^4$ Monte Carlo realizations of
  $N_{\mathrm{eff}}$ frequencies drawn from a
  Poisson process with the same mean density.
  For each realization, compute the KS distance
  to GUE and record~$\mathcal{Q}$.

\item Generate $10^4$ realizations using the
  \emph{actual} first $N_{\mathrm{eff}}$ Riemann
  zeros with Gaussian frequency errors
  $\delta\omega_{\mathrm{res}}$.
  Compute the same statistics.

\item The $p$-value for GUE rejection under $H_0$
  (Poisson) is the fraction of Poisson realizations
  with KS distance to GUE \emph{smaller} than the
  observed value.  For $N_{\mathrm{beat}}\ge 6$,
  we estimate $p\lesssim 0.05$ is achievable,
  rising to $p\lesssim 0.01$ for
  $N_{\mathrm{eff}}\ge 7$
  (corresponding to $\sigma\lesssim 5\times 10^{-4}$).
\end{enumerate}


%-----------------------------------------------
\subsection{Visibility requirements and the sweet spot}
\label{sec:visibility}

The GUE cross-check requires at least two resolved
spectral lines, i.e.\ $N_{\mathrm{eff}}\ge 2$.
From Eq.~\eqref{eq:Neff}, the condition
$e^{-2\sigma\gamma_2^2}\ge e^{-1}$ (i.e., the second
zero $\gamma_2 = 21.022$ contributes at least $1/e$
of its unsuppressed amplitude) gives
\begin{equation}\label{eq:sigmacut}
  \sigma \;\lesssim\;
  \frac{1}{2\,\gamma_2^2}
  \;\approx\;
  1.1\times 10^{-3}\,.
\end{equation}
A less stringent requirement, $N_{\mathrm{eff}}\ge 2$
(i.e.\ the sum in Eq.~\eqref{eq:Neff} exceeds~2), yields
\begin{equation}\label{eq:sigmavis}
  \sigma \;\lesssim\; 0.002\,.
\end{equation}
This is the \emph{beat-frequency visibility boundary}
identified in Sec.~\ref{sec:paramspace}.  Crucially,
it overlaps with the region of maximal signal amplitude
$|\delta\kappa/\kappa|$
(which is maximized for $\sigma\to 0$ at fixed~$\varepsilon$),
so the sweet spot $\sigma\sim 10^{-3}$ simultaneously
\begin{enumerate}
\item maximizes detectability of the overall modulation, and
\item enables the GUE cross-check with
  $N_{\mathrm{beat}}\ge 6$.
\end{enumerate}
This coincidence is not a fine-tuning but a structural
feature of the framework: the same Gaussian damping
that controls signal amplitude also controls the number
of resolved spectral lines.


%-----------------------------------------------
\subsection{Illustrative numerical example}
\label{sec:gue-example}

To ground the preceding analysis, we evaluate the
beat-frequency statistics for the first four Riemann zeros:
\begin{equation}\label{eq:firstzeros}
  \gamma_1 = 14.1347\,,\quad
  \gamma_2 = 21.0220\,,\quad
  \gamma_3 = 25.0109\,,\quad
  \gamma_4 = 30.4249\,.
\end{equation}
The six pairwise differences are:
\begin{center}
\renewcommand{\arraystretch}{1.2}
\begin{tabular}{ccc}
\hline
$(i,j)$ & $\Delta_{ij}$ & $s_{ij} = \Delta_{ij}/\langle\Delta\rangle$ \\
\hline
(1,2) & $6.887$ & $1.27$ \\
(1,3) & $10.876$ & $2.00$ \\
(1,4) & $16.290$ & $3.00$ \\
(2,3) & $3.989$ & $0.73$ \\
(2,4) & $9.403$ & $1.73$ \\
(3,4) & $5.414$ & $1.00$ \\
\hline
\end{tabular}
\end{center}
where $\langle\Delta\rangle = 8.81$.
The key observation is that
$s_{\min} = 0.73$, achieved by the $(2,3)$ pair.
No normalized spacing falls below~$0.1$---indeed,
none falls below~$0.5$.
This reflects the quadratic level repulsion
of Eq.~\eqref{eq:pGUE}: GUE zeros avoid clustering.

For comparison, four Poisson-distributed frequencies
with the same mean spacing would have
$\mathrm{Prob}(s_{\min}<0.1)\approx 34\%$
for $N_{\mathrm{beat}}=6$ independent pairs.
The observed $s_{\min}=0.73$ thus already rejects
Poisson at the $\sim\!2\sigma$ level from these four
zeros alone.


%-----------------------------------------------
\subsection{Discussion: what the GUE test does and does not prove}
\label{sec:gue-discussion}

\paragraph{What it does.}
A successful GUE cross-check demonstrates that the
spectral lines in $\kappa_{\mathrm{eff}}(k)$ are not
randomly placed but obey the specific correlation
structure of the Riemann zeta zeros.
Via Rodgers' theorem~\cite{Rodgers2018}, this
is equivalent to a statement about the distribution
of prime powers---the deepest known connection between
physics and number theory.
The test is amplitude-independent: even if~$\varepsilon$
is too small for direct SGWB detection, the GUE
signature can in principle be extracted from the
\emph{relative} positions of resolved spectral lines
in any sufficiently precise power-spectrum measurement.

\paragraph{What it does not.}
The GUE test does not prove that the environment
\emph{is} built from Riemann zeros;
it proves only that the environment's spectral statistics
are \emph{consistent with} the GUE universality class.
Other physical systems (quantum billiards, disordered
conductors) share this universality class.
The additional requirement that the individual
frequencies match known values of~$\gamma_n$
(Stage~2 of the pipeline) is needed to distinguish
the Riemann-zero hypothesis from a generic
chaotic-quantum-system hypothesis.
Together, frequency identification \emph{and}
GUE pair correlation constitute a two-pronged test
that, if passed simultaneously, would provide
compelling evidence for the arithmetic structure of
quantum gravity noise.


%———————————————————————————————————————————
% §6  Observational Signatures
%———————————————————————————————————————————

\section{Observational Signatures}\label{sec:observational}

The preceding sections derived the complex gravitational coupling
$\kappa_{\mathrm{eff}}(k)$ (Sec.~\ref{sec:complexkappa}),
its Zeta-Comb modulation spectrum (Sec.~\ref{sec:zetacomb}),
and the GUE internal consistency check (Sec.~\ref{sec:gue}).
In this section we map these theoretical results onto the
space of observable parameters, identify the detector
configurations capable of testing the predictions,
and demonstrate that the parameter sweet spot
$(\sigma\sim 10^{-3},\;\varepsilon\sim 10^{-2})$
simultaneously satisfies existing constraints,
enables GUE cross-checks, and falls within the
projected sensitivity of next-generation
gravitational-wave observatories.


%-----------------------------------------------
\subsection{The two-dimensional parameter space}
\label{sec:paramspace}

\subsubsection{Definition of the control parameters}

The AL-GFT Zeta-Comb modulation derived in
Sec.~\ref{sec:zetacomb} is controlled by two
phenomenological parameters:
\begin{itemize}
\item $\varepsilon$\;---\;the \emph{multiplicity coupling
  strength}, governing the overall amplitude of the
  environment's back-reaction on the gravitational
  sector.  It enters the signal amplitude quadratically:
  $|\delta\kappa/\kappa|\propto\varepsilon^2$.

\item $\sigma$\;---\;the \emph{soft-resonance width},
  a positive dimensionless parameter controlling
  the Gaussian damping of successive Riemann-zero
  contributions via
  $\exp(-\sigma\,\gamma_n^2)$ in $\mathcal{M}(k)$
  and $\exp(-2\sigma\,\gamma_n^2)$ in the signal
  power $|\delta\kappa/\kappa|^2$.
\end{itemize}
All other quantities entering the modulation---the
locations of the zeros $\gamma_n$, the phases
$\phi_n = \tfrac{1}{2}\arg\,\zeta(\tfrac{1}{2}+i\gamma_n)$---are
fixed by number theory and carry no free parameters.
The framework therefore has a \textbf{two-dimensional}
parameter space $(\varepsilon,\sigma)$, making it
maximally constrained for a stochastic-gravity extension.


\subsubsection{Signal amplitude}

From the influence-functional derivation of
Sec.~\ref{sec:complexkappa}, the fractional shift
in the gravitational coupling at wavenumber~$k$ is
\begin{equation}\label{eq:deltakappa}
  \frac{\delta\kappa}{\kappa}(k)
  \;=\;
  \varepsilon^2
  \sum_{n=1}^{\infty}
  e^{-2\sigma\,\gamma_n^2}\,
  e^{i\gamma_n\ln(k/k_\star) + i\phi_n}\,.
\end{equation}
The modulus of this quantity is bounded by
\begin{equation}\label{eq:deltakappa-bound}
  \left|\frac{\delta\kappa}{\kappa}\right|
  \;\le\;
  \varepsilon^2\,
  \sum_{n=1}^{\infty} e^{-2\sigma\,\gamma_n^2}
  \;=\;
  \varepsilon^2\,N_{\mathrm{eff}}(\sigma)\,,
\end{equation}
where $N_{\mathrm{eff}}(\sigma)$ was defined in
Eq.~\eqref{eq:Neff}.  In the regime
$\sigma\ll 1/\gamma_1^2\approx 5\times 10^{-3}$,
$N_{\mathrm{eff}}$ grows (more zeros contribute),
but the sum remains finite for any $\sigma>0$.
For the dominant contribution from $\gamma_1=14.1347$
alone:
\begin{equation}\label{eq:signal-leading}
  \left|\frac{\delta\kappa}{\kappa}\right|_{\gamma_1}
  \;=\;
  \varepsilon^2\,
  e^{-2\sigma\,\gamma_1^2}
  \;=\;
  \varepsilon^2\,
  e^{-399.6\,\sigma}\,.
\end{equation}


\subsubsection{Numerical evaluation at the sweet spot}

At the parameter values $\varepsilon = 0.01$,
$\sigma = 0.001$:
\begin{align}
  e^{-2\sigma\gamma_1^2}
  &= e^{-0.3996} = 0.671\,,
  \label{eq:damp1}\\[4pt]
  e^{-2\sigma\gamma_2^2}
  &= e^{-0.884} = 0.413\,,
  \label{eq:damp2}\\[4pt]
  e^{-2\sigma\gamma_3^2}
  &= e^{-1.251} = 0.286\,,
  \label{eq:damp3}\\[4pt]
  e^{-2\sigma\gamma_4^2}
  &= e^{-1.851} = 0.157\,.
  \label{eq:damp4}
\end{align}
Summing the first ten zeros gives
$N_{\mathrm{eff}} \approx 4.2$, and the signal amplitude
\begin{equation}\label{eq:signal-sweetspot}
  \left|\frac{\delta\kappa}{\kappa}\right|
  \;\approx\;
  (0.01)^2 \times 0.671
  \;=\;
  6.7\times 10^{-5}\,.
\end{equation}
This is the benchmark value used throughout the
subsequent detectability analysis.


%-----------------------------------------------
\subsection{Existing constraints}
\label{sec:constraints}

\subsubsection{Planck 2018: oscillatory features in the
  primordial power spectrum}

The Planck 2018 analysis~\cite{Planck2018X} performed
a dedicated search for oscillatory features in the
primordial scalar power spectrum, parameterized as
\begin{equation}\label{eq:planck-osc}
  \mathcal{P}(k)
  \;=\;
  \mathcal{P}_0(k)\,
  \bigl[1 + A_{\mathrm{osc}}\,
  \sin(\omega_{\mathrm{osc}}\,\ln k + \varphi)\bigr]\,,
\end{equation}
with $A_{\mathrm{osc}}$, $\omega_{\mathrm{osc}}$, and
$\varphi$ as free parameters.
No statistically significant oscillatory signal was
detected; the resulting upper bound is
$A_{\mathrm{osc}} \lesssim 0.03$ (95\% CL) for
$\omega_{\mathrm{osc}}\sim 10$--$30$, which covers
the range spanned by $\gamma_1$--$\gamma_4$.

In our framework, the oscillation amplitude is
\begin{equation}\label{eq:Aosc-map}
  A_{\mathrm{osc}}
  \;\approx\;
  2\,\varepsilon^2\,e^{-2\sigma\gamma_n^2}\,,
\end{equation}
where the factor of~2 accounts for the modulus of the
complex exponential in~$\mathcal{M}(k)$.
Imposing $A_{\mathrm{osc}}<0.03$ at $\gamma_1$ gives:
\begin{equation}\label{eq:planck-bound}
  \varepsilon^2\,e^{-2\sigma\gamma_1^2}
  \;<\;
  0.015
  \qquad\Longrightarrow\qquad
  \varepsilon \;<\;
  0.12\,e^{\sigma\gamma_1^2}\,.
\end{equation}
At $\sigma = 10^{-3}$, this yields
$\varepsilon < 0.15$---well above our benchmark
$\varepsilon = 0.01$.
The sweet spot lies safely within the Planck-allowed
region by more than an order of magnitude.


\subsubsection{CMB non-Gaussianity}

The AL-GFT Gaussian branch (Gate~1) predicts
$f_{\mathrm{NL}}^{\mathrm{local}} = 0$ by construction,
since the environment is Gaussian and the coupling is
linear.  The Planck 2018 constraint
$f_{\mathrm{NL}}^{\mathrm{local}} = -0.9\pm 5.1$
(68\%~CL)~\cite{Planck2018IX} is therefore
automatically satisfied.
Any future detection of $f_{\mathrm{NL}}\neq 0$
would falsify the Gaussian branch and point toward
the non-Gaussian extension (ML-GFT).


\subsubsection{Gravitational-wave propagation}

The imaginary part of $\kappa_{\mathrm{eff}}$
modifies the dispersion relation for gravitational
waves, introducing a frequency-dependent damping.
Current LIGO-Virgo-KAGRA observations (GWTC-4.0)
constrain modifications to the GW dispersion at the level
$\delta c_T/c_T \lesssim 10^{-15}$ for
$f\sim 10$--$10^3\,$Hz.
At these frequencies the AL-GFT modulation is
suppressed by $\exp(-2\sigma\gamma_1^2\ln^2 f/f_\star)$
with $f_\star$ at cosmological scales, rendering
the correction negligibly small
($\ll 10^{-20}$).  There is no tension with
GW propagation bounds.


%-----------------------------------------------
\subsection{The parameter-space landscape}
\label{sec:landscape}

\subsubsection{Four regions of the $(\sigma,\varepsilon)$ plane}

The $(\sigma,\varepsilon)$ parameter space is
structured by four boundaries that carve it into
physically distinct regions (Fig.~\ref{fig:paramspace}):

\begin{enumerate}
\item \textbf{Planck exclusion} (upper boundary).
  The Planck 2018 constraint on oscillatory features,
  Eq.~\eqref{eq:planck-bound}, excludes the region
  \begin{equation}
    \varepsilon \;>\;
    0.12\,e^{\sigma\,\gamma_1^2}
    \;\approx\;
    0.12\,e^{200\sigma}\,.
  \end{equation}
  This is an exponentially rising curve in the
  $(\sigma,\varepsilon)$ plane, becoming less
  constraining as $\sigma$ increases (stronger damping
  reduces the observable amplitude).

\item \textbf{BBO/DECIGO threshold} (lower boundary).
  The projected strain sensitivity of BBO (two correlated
  clusters, $T_{\mathrm{obs}}=3\,$yr) at
  $f\sim 0.1$--$1\,$Hz sets a minimum detectable
  energy-density fraction
  $\Omega_{\mathrm{GW}}^{\min}\sim 10^{-17}$.
  Mapping this to the modulation amplitude via
  $\Omega_{\mathrm{GW}}\propto
  |\delta\kappa/\kappa|^2\,\Omega_{\mathrm{GW},0}$
  (where $\Omega_{\mathrm{GW},0}$ is the standard
  inflationary background), the detection threshold
  requires
  \begin{equation}\label{eq:bbo-threshold}
    |\delta\kappa/\kappa|
    \;\gtrsim\;
    \delta_{\mathrm{det}}\,,
  \end{equation}
  with $\delta_{\mathrm{det}}\sim 10^{-5}$--$10^{-4}$
  depending on the assumed inflationary tensor-to-scalar
  ratio~$r$.

\item \textbf{Beat-frequency visibility boundary}
  (left boundary).
  The GUE cross-check of Sec.~\ref{sec:gue}
  requires $N_{\mathrm{eff}}\ge 2$, which from
  Eq.~\eqref{eq:sigmavis} translates to
  \begin{equation}\label{eq:sigma-vis}
    \sigma \;\lesssim\; 0.002\,.
  \end{equation}
  To the right of this boundary, only a single
  spectral line ($\gamma_1$) contributes, and no
  pair-correlation test is possible.

\item \textbf{Perturbativity bound} (far-left boundary).
  The influence-functional derivation assumes
  $\varepsilon\ll 1$ (weak coupling to the environment).
  We conservatively impose $\varepsilon < 0.1$,
  beyond which higher-order terms in the
  Dyson series may become significant.
\end{enumerate}


\subsubsection{The sweet spot}

The intersection of these four boundaries defines a
\emph{sweet spot} in parameter space:
\begin{equation}\label{eq:sweetspot}
  \boxed{%
    \sigma \sim 10^{-3}\,,
    \qquad
    \varepsilon \sim 10^{-2}\,.
  }
\end{equation}
At this point:
\begin{itemize}
\item The signal amplitude
  $|\delta\kappa/\kappa|\approx 6.7\times 10^{-5}$
  [Eq.~\eqref{eq:signal-sweetspot}] exceeds the
  BBO detection threshold;
\item The Planck constraint is satisfied by a
  factor of~$\sim\!15$;
\item $N_{\mathrm{eff}}\approx 4.2$ zeros contribute,
  enabling the GUE cross-check with
  $N_{\mathrm{beat}} = 6$ beat frequencies;
\item The perturbative expansion is well-controlled
  ($\varepsilon^2 = 10^{-4}$, so the
  next-order correction is $\mathcal{O}(\varepsilon^4)
  \sim 10^{-8}$).
\end{itemize}
The coincidence of these four conditions at a single
point in the $(\sigma,\varepsilon)$ plane is a
\emph{structural feature} of the framework,
not a fine-tuning.  It arises because the same Gaussian
damping $\exp(-2\sigma\gamma_n^2)$ simultaneously
controls three quantities: signal amplitude, number of
active zeros, and departure from the Planck constraint.


%-----------------------------------------------
\subsection{Contour plot of the parameter space}
\label{sec:contour}

Figure~\ref{fig:paramspace} displays the parameter-space
landscape.  The axes are $\sigma$
(horizontal, logarithmic, range $10^{-4}$--$10^{-1}$)
and $\varepsilon$ (vertical, logarithmic,
range $10^{-3}$--$1$).

\begin{figure}[t]
\centering
\includegraphics[width=0.85\textwidth]{parameter_space_final.png}
\caption{%
  Parameter space of the complex gravitational coupling
  in the $(\sigma,\varepsilon)$ plane.
  \textbf{Filled contours:} signal amplitude
  $\log_{10}|\delta\kappa/\kappa|$, computed from
  Eq.~\eqref{eq:deltakappa-bound} summing the first
  100~Riemann zeros.
  \textbf{Solid red curve:} Planck 2018 exclusion
  boundary (oscillatory features,
  $A_{\mathrm{osc}}<0.03$).
  \textbf{Dashed blue line:} BBO/DECIGO detection
  threshold ($|\delta\kappa/\kappa|
  \gtrsim 10^{-5}$).
  \textbf{Dashed green line:} beat-frequency visibility
  boundary ($\sigma = 0.002$,
  $N_{\mathrm{eff}}=2$).
  \textbf{Star:} sweet spot
  $(\sigma,\varepsilon) = (10^{-3},\,10^{-2})$,
  where the signal amplitude is
  $|\delta\kappa/\kappa|\approx 6.7\times 10^{-5}$,
  the GUE cross-check is enabled, and the Planck
  bound is satisfied.
  The allowed region lies below the red curve,
  above the blue line, and to the left of the
  green line.}
\label{fig:paramspace}
\end{figure}

The contour plot encodes several key features:
\begin{enumerate}
\item \emph{Signal amplitude contours} slope
  diagonally from upper-left to lower-right,
  reflecting the competition between
  $\varepsilon^2$ (increasing upward) and
  $e^{-2\sigma\gamma_1^2}$ (decreasing rightward).

\item The \emph{allowed region} (below Planck,
  above BBO, left of beat-visibility) forms a
  wedge whose apex is the sweet spot.

\item For $\sigma > 0.01$, the Planck bound becomes
  so weak ($\varepsilon < 1$) that it is irrelevant,
  but the signal amplitude drops below
  $10^{-8}$---far below any projected detector
  sensitivity.

\item For $\sigma < 10^{-4}$, $N_{\mathrm{eff}}$ grows
  large ($\gtrsim 50$), but the sum of oscillatory
  terms begins to average down the net modulation
  unless the phases $\phi_n$ are coherent over
  many zeros---an unlikely scenario for the first
  $\sim\!50$ Riemann zeros.
\end{enumerate}


%-----------------------------------------------
\subsection{Detector mapping}
\label{sec:detectors}

\subsubsection{Translation to gravitational-wave
  energy density}

The complex gravitational coupling modifies the
tensor transfer function, imprinting the Zeta-Comb
modulation onto the stochastic gravitational-wave
background (SGWB).  The fractional energy density
spectrum acquires an oscillatory component:
\begin{equation}\label{eq:OmegaGW}
  \Omega_{\mathrm{GW}}(f)
  \;=\;
  \Omega_{\mathrm{GW}}^{(0)}(f)
  \;\bigl[1 + \mathcal{M}(f)\bigr]\,,
\end{equation}
where $\Omega_{\mathrm{GW}}^{(0)}(f)$ is the standard
inflationary background and $\mathcal{M}(f)$ is the
modulation function evaluated at $k = 2\pi f/c$
(neglecting the weak $k$-dependence of the transfer
function over the detector bandwidth).

The amplitude of the modulation at frequency~$f$ is
\begin{equation}\label{eq:modamp-f}
  |\mathcal{M}(f)|
  \;\approx\;
  2\,\varepsilon^2\,
  e^{-2\sigma\gamma_1^2}\,
  \bigl|\cos\bigl(
    \gamma_1\ln(f/f_\star) + \phi_1
  \bigr)\bigr|
  \;+\;\text{higher zeros}\,,
\end{equation}
and the induced spectral modulation of the
energy density is
\begin{equation}\label{eq:deltaOmega}
  \delta\Omega_{\mathrm{GW}}
  \;\equiv\;
  |\mathcal{M}|\;\Omega_{\mathrm{GW}}^{(0)}\,.
\end{equation}


\subsubsection{Detector sensitivity overview}

Table~\ref{tab:detectors} summarizes the
gravitational-wave observatories relevant to our
predictions, ordered by target frequency band and
projected timeline.

\begin{table}[t]
\centering
\renewcommand{\arraystretch}{1.3}
\caption{%
  Detector configurations and their relevance to the
  Zeta-Comb gravitational-wave signature.
  $\Omega_{\mathrm{sens}}$ denotes the
  power-law-integrated sensitivity to an SGWB
  for the stated observation time and correlation
  method.  ``Modulation SNR'' estimates the
  signal-to-noise ratio for the oscillatory
  component $\delta\Omega_{\mathrm{GW}}$
  at the sweet-spot parameters, assuming
  $r=0.01$ for the inflationary tensor background.}
\label{tab:detectors}
\begin{tabular}{lccccc}
\hline
\textbf{Detector}
  & \textbf{Band}
  & $\boldsymbol{\Omega_{\mathrm{sens}}}$
  & \textbf{$T_{\mathrm{obs}}$}
  & \textbf{Timeline}
  & \textbf{Mod.\ SNR} \\
\hline
LISA
  & $10^{-4}$--$10^{-1}\,$Hz
  & $\sim 10^{-12}$
  & 4\,yr
  & 2035
  & $\ll 1$ \\
B-DECIGO
  & $0.1$--$10\,$Hz
  & $\sim 10^{-15}$
  & 3\,yr
  & 2030s
  & $\sim 0.1$ \\
DECIGO
  & $0.1$--$10\,$Hz
  & $\sim 10^{-18}$
  & 3\,yr
  & late 2030s
  & $\sim 1$--$3$ \\
BBO
  & $0.03$--$3\,$Hz
  & $\sim 10^{-17}$
  & 5\,yr
  & 2040s--2050s
  & $\sim 3$--$10$ \\
Einstein Tel.\
  & $3$--$10^4\,$Hz
  & $\sim 10^{-13}$
  & 1\,yr
  & $\sim$2035
  & $\ll 1$ \\
Cosmic Expl.\
  & $5$--$10^4\,$Hz
  & $\sim 10^{-13}$
  & 1\,yr
  & late 2030s
  & $\ll 1$ \\
\hline
\end{tabular}
\end{table}


\subsubsection{LISA}

The Laser Interferometer Space Antenna, with a
planned launch in 2035 and a sensitivity band of
$10^{-4}$--$10^{-1}\,$Hz, will be the first
space-based gravitational-wave
observatory~\cite{LISA2024}.
Its SGWB sensitivity, characterized by the
power-law integrated curve, reaches
$\Omega_{\mathrm{GW}}\sim 10^{-12}$ for
$T_{\mathrm{obs}}=4\,$yr.

For the sweet-spot parameters, the modulation
amplitude at $f=10^{-2}\,$Hz is
\begin{equation}
  \delta\Omega_{\mathrm{GW}}^{\mathrm{LISA}}
  \;\sim\;
  6.7\times 10^{-5} \times
  \Omega_{\mathrm{GW}}^{(0)}(10^{-2}\,\text{Hz})\,.
\end{equation}
For standard slow-roll inflation with $r=0.01$,
$\Omega_{\mathrm{GW}}^{(0)}(10^{-2}\,\text{Hz})
\sim 10^{-16}$,
giving
$\delta\Omega_{\mathrm{GW}}\sim 10^{-20}$---eight
orders of magnitude below LISA's reach.
\textbf{LISA cannot detect the Zeta-Comb modulation
for standard inflationary backgrounds.}
However, if an exotic cosmological SGWB of non-inflationary
origin exists at LISA frequencies with
$\Omega\sim 10^{-10}$, the modulation would be
marginally detectable.


\subsubsection{B-DECIGO and DECIGO}

B-DECIGO, the scientific pathfinder for DECIGO planned
for the 2030s, covers $0.1$--$10\,$Hz with
strain sensitivity $\sim 10^{-23}\,\mathrm{Hz}^{-1/2}$
around 1\,Hz.
DECIGO itself targets
$4\times 10^{-24}\,\mathrm{Hz}^{-1/2}$ at 1\,Hz
for a single cluster, improving to
$7\times 10^{-26}\,\mathrm{Hz}^{-1/2}$ with
two co-located clusters using
cross-correlation~\cite{DECIGO2021,DECIGO2024}.

The DECIGO band is the \emph{natural home} for
the Zeta-Comb signature: the frequency
$f\sim 0.1$--$1\,$Hz corresponds to comoving
wavenumbers $k\sim 10^{12}$--$10^{13}\,\mathrm{Mpc}^{-1}$,
far below the Silk-damping scale, where the
primordial tensor spectrum is unaffected by
CMB foreground systematics.

For $r=0.01$ and the DECIGO correlation sensitivity:
\begin{equation}\label{eq:decigo-snr}
  \mathrm{SNR}_{\mathrm{DECIGO}}
  \;\sim\;
  \frac{\delta\Omega_{\mathrm{GW}}}
       {\Omega_{\mathrm{sens}}^{\mathrm{DECIGO}}}
  \;\sim\;
  \frac{6.7\times 10^{-5}\times 10^{-16}}
       {10^{-18}}
  \;\sim\; 1\text{--}3\,,
\end{equation}
placing the signal at the threshold of detectability.
A dedicated template-based search exploiting the
known Zeta-Comb frequency structure could enhance
this by a factor $\sim\!\sqrt{N_{\mathrm{eff}}}
\approx 2$.


\subsubsection{BBO}

The Big Bang Observer, a proposed
constellation of four LISA-like clusters in
heliocentric orbit with arm lengths of
$5\times 10^4\,$km, is specifically designed for
the direct detection of the inflationary SGWB
in the 0.03--3\,Hz band~\cite{BBO2005,Yagi2011}.
Its projected sensitivity
$\Omega_{\mathrm{GW}}\sim 10^{-17}$ after
foreground subtraction makes it the \emph{primary
target detector} for the Zeta-Comb signature.

At the sweet spot:
\begin{equation}\label{eq:bbo-snr}
  \mathrm{SNR}_{\mathrm{BBO}}
  \;\sim\;
  \frac{6.7\times 10^{-5}\times 10^{-16}}
       {10^{-17}}
  \;\sim\; 3\text{--}10\,,
\end{equation}
depending on the exact value of~$r$ and the
efficiency of astrophysical foreground removal
(primarily neutron-star binaries, which must
be individually identified and
subtracted~\cite{Yagi2011}).
This places the signal in the regime of
\emph{confident detection}
($\mathrm{SNR}\gtrsim 3$), with the oscillatory
structure of the modulation providing additional
template-matched enhancement.

\subsubsection{Ground-based next-generation detectors}

The Einstein Telescope and Cosmic Explorer,
with sensitivities exceeding Advanced LIGO by a
factor of~$\sim\!10$ and frequency bands of
3--$10^4\,$Hz, are optimized for compact binary
mergers rather than primordial SGWB
detection~\cite{ET2020,CE2024}.
Their SGWB sensitivity
($\Omega_{\mathrm{GW}}\sim 10^{-13}$) is
insufficient for the Zeta-Comb signal at standard
inflationary amplitudes.
However, these detectors contribute indirectly:
\begin{itemize}
\item By measuring the GW propagation speed to
  $\delta c_T/c_T \sim 10^{-16}$ or better,
  they constrain the real part of $\kappa_{\mathrm{eff}}$
  at high frequencies.
\item By cataloguing $\mathcal{O}(10^5)$ binary
  mergers per year, they enable \emph{statistical}
  tests of the frequency-dependent damping rate
  $\mathrm{Im}\,\kappa_{\mathrm{eff}}(f)$
  averaged over many events.
\end{itemize}


%-----------------------------------------------
\subsection{CMB observational channel}
\label{sec:cmb-channel}

\subsubsection{Power spectrum modulation}

The Zeta-Comb modulation also imprints on the
scalar primordial power spectrum via the noise kernel:
\begin{equation}\label{eq:Pk-mod}
  \mathcal{P}(k)
  \;=\;
  A_s\left(\frac{k}{k_\star}\right)^{n_s-1}
  \bigl[1 + \mathcal{M}(k)\bigr]\,,
\end{equation}
where $\mathcal{M}(k)$ is the same modulation function
entering the gravitational coupling.
The angular power spectrum $C_\ell$ inherits the
oscillatory structure through the transfer functions:
\begin{equation}\label{eq:Cell-mod}
  \delta C_\ell
  \;\propto\;
  \int_0^\infty \frac{dk}{k}\,
  \mathcal{M}(k)\,
  |\Delta_\ell(k)|^2\,,
\end{equation}
where $\Delta_\ell(k)$ is the radiation transfer
function.  The integration averages over many
oscillation cycles for $\gamma_n\gg 1$, reducing
the effective amplitude by a factor
$\sim 1/\sqrt{\ell_{\max}/\ell_{\min}}$.

\subsubsection{Forecasts for LiteBIRD and CMB-S4}

LiteBIRD (launch 2032) will measure the CMB
polarization to $r\sim 10^{-3}$, while CMB-S4
(first light $\sim$2030) will achieve cosmic-variance-limited
$E$-mode measurements to $\ell\sim 5000$.
Neither experiment is optimized for oscillatory
features at $\omega\sim 14$--$30$ in $\ln k$, but
a dedicated search combining TT, TE, and EE
spectra could in principle reach
$A_{\mathrm{osc}}\sim 10^{-3}$ at these
frequencies.

For the sweet-spot parameters,
$A_{\mathrm{osc}}\approx 2\times 6.7\times 10^{-5}
\approx 1.3\times 10^{-4}$,
which is a factor of~$\sim\!8$ below the projected
sensitivity.
\textbf{Direct detection of the Zeta-Comb in CMB
power spectra requires $\varepsilon\gtrsim 0.03$}
(at $\sigma = 10^{-3}$), which remains within
the Planck-allowed region but approaches its boundary.


%-----------------------------------------------
\subsection{The detection roadmap}
\label{sec:roadmap}

We summarize the observational prospects in a
three-stage roadmap, ordered by timeline and
detection likelihood.

\begin{description}
\item[Stage~I (2025--2035): Archival and indirect.]
  \begin{itemize}
  \item Re-analyze Planck 2018 TT/TE/EE data with a
    matched-filter template tuned to the Zeta-Comb
    frequencies $\gamma_1$--$\gamma_4$.  Expected
    outcome: SNR $\lesssim 1$ at the sweet spot, but
    the search establishes the pipeline and sets
    improved model-specific bounds on $\varepsilon$.
  \item LiteBIRD $B$-mode measurement constrains~$r$,
    which is a prerequisite input for the SGWB
    detection forecast (Stage~II).
  \end{itemize}

\item[Stage~II (2030s--2040s): Decihertz window.]
  \begin{itemize}
  \item B-DECIGO ($\sim$2030s) provides the first
    direct SGWB search in the 0.1--10\,Hz band.
    Marginal sensitivity to the Zeta-Comb modulation
    ($\mathrm{SNR}\sim 0.1$) at the sweet spot,
    but useful for foreground characterization and
    pipeline validation.
  \item DECIGO (late 2030s) reaches
    $\mathrm{SNR}\sim 1$--$3$ at the sweet spot.
    A template search exploiting the known Zeta-Comb
    structure could achieve $\mathrm{SNR}\sim 3$--$5$.
    \textbf{First possible detection.}
  \end{itemize}

\item[Stage~III (2040s--2050s): Definitive test.]
  \begin{itemize}
  \item BBO achieves $\mathrm{SNR}\sim 3$--$10$ with
    broad frequency coverage (0.03--3\,Hz).
    The oscillatory pattern across $\sim\!2$ decades
    of frequency enables both:
    \begin{enumerate}
    \item direct identification of individual spectral
      lines ($\gamma_1$, $\gamma_2$, possibly
      $\gamma_3$);
    \item the GUE cross-check via the pair-correlation
      pipeline of Sec.~\ref{sec:pipeline}.
    \end{enumerate}
  \item A detection-plus-GUE-confirmation at BBO would
    constitute observational evidence that the
    statistical structure of the Riemann zeta zeros
    is encoded in the gravitational-wave spectrum---a
    result connecting quantum gravity to number
    theory via a physical measurement.
  \end{itemize}
\end{description}


%-----------------------------------------------
\subsection{The $\sigma$ sweet spot as a structural prediction}
\label{sec:sigmastructure}

The sweet spot $\sigma\sim 10^{-3}$ is not a
post hoc choice but emerges from the structure of the
Riemann zeros themselves.  Specifically:

\begin{enumerate}
\item The first zero $\gamma_1 = 14.1347$ sets the
  scale: $\sigma_{\mathrm{crit}} \equiv
  1/(2\gamma_1^2) \approx 2.5\times 10^{-3}$.
  For $\sigma\gg\sigma_{\mathrm{crit}}$,
  even the leading zero is damped out and the
  modulation vanishes.

\item The beat-frequency visibility threshold
  $\sigma\lesssim 0.002$
  [Eq.~\eqref{eq:sigmavis}] is set by $\gamma_2$:
  $1/(2\gamma_2^2)\approx 1.1\times 10^{-3}$.

\item The ratio $\gamma_2/\gamma_1 \approx 1.49$,
  which controls the separation between
  $\sigma_{\mathrm{crit}}$ and the beat-visibility
  threshold, is a \emph{number-theoretic invariant}.
  No parameter tuning can change it.
\end{enumerate}

\noindent
The conclusion is that any framework embedding the
Riemann zeros into a Gaussian-damped environment will
\emph{generically} find its maximal-information regime
at $\sigma\sim \gamma_1^{-2}$, regardless of the
physical context.  The sweet spot is a prediction
of number theory, not of the model.


%%% References cited in this section:
% \cite{Planck2018X}  — Planck Collaboration, A&A 641, A10 (2020).
% \cite{Planck2018IX} — Planck Collaboration, A&A 641, A9 (2020).
% \cite{LISA2024}     — LISA Collaboration, ESA Definition Study (2024).
% \cite{DECIGO2021}   — Kawamura et al., PTEP 2021, 05A105 (2021).
% \cite{DECIGO2024}   — Kawamura et al., Proc. Sci. 444, 1516 (2024).
% \cite{BBO2005}      — Crowder & Cornish, Phys. Rev. D 72, 083005 (2005).
% \cite{Yagi2011}     — Yagi & Seto, Phys. Rev. D 83, 044011 (2011).
% \cite{ET2020}       — Maggiore et al., JCAP 03, 050 (2020).
% \cite{CE2024}       — Cosmic Explorer Horizon Study (2024).
%% =============================================================
%% §7  The Self-Consistent Bootstrap
%% Paper: "Complex Gravitational Coupling from Arithmetic Quantum Gravity Noise"
%% =============================================================

\section{The Self-Consistent Bootstrap}
\label{sec:bootstrap}

The preceding sections treated the Zeta-Comb environment as \emph{given}:
a noise kernel $N(k)$ characterised by parameters $(\varepsilon,\sigma)$
was imposed, and its observable consequences---complex gravitational
coupling, GUE beat spectrum, gravitational-wave signatures---were derived.
A deeper question remains: \emph{can the environment that dresses $\kappa$
determine its own parameters self-consistently?}

In this section we recast the framework as a nonlinear eigenvalue problem
whose solution, if it exists, would fix $(\varepsilon,\sigma)$ from the
Riemann zero spectrum alone---eliminating them as free parameters and
elevating the theory from a model to a self-consistent bootstrap.
We outline the formal structure, identify the key mathematical obstacles,
and propose a concrete programme for future work.


%% -------------------------------------------------------------
\subsection{From dressed coupling to self-consistency condition}
\label{subsec:dressing-selfconsistency}

Recall from \S\ref{sec:complex-kappa} that integrating out the
arithmetic environment yields an effective gravitational coupling
(Eq.~\eqref{eq:kappa-eff})
%
\begin{equation}
  \kappa_{\mathrm{eff}}(k)
  \;=\;
  \frac{\kappa}
       {1 \;-\; \kappa\,\widetilde{D}_{R}(k)\big/\mathcal{O}(k)}\,,
  \label{eq:kappa-eff-recall}
\end{equation}
%
where $\widetilde{D}_{R}(k)$ is the retarded dissipation kernel and
$\mathcal{O}(k)$ encodes the system response.
Both kernels inherit the Zeta-Comb structure of the noise
(Eq.~\eqref{eq:noise-kernel}):
%
\begin{equation}
  N(k)
  \;=\;
  N_{0}(k)\,\mathcal{M}(k)\,,
  \qquad
  \mathcal{M}(k)
  \;=\;
  1 \;+\; 2\,\varepsilon\!\sum_{n=1}^{N_{\mathrm{eff}}}
          e^{-\sigma\gamma_{n}^{2}}\,
          \cos\!\bigl(\gamma_{n}\ln(k/k_{\star})+\phi_{n}\bigr)
  \;+\; \mathcal{O}(\varepsilon^{2})\,,
  \label{eq:modulation-recall}
\end{equation}
%
where $\gamma_{n}$ are the imaginary parts of the non-trivial Riemann
zeta zeros, $\sigma$ is the soft-resonance wid%% =============================================================
%% §8  Discussion
%% Paper: "Complex Gravitational Coupling from Arithmetic Quantum Gravity Noise"
%% =============================================================

\section{Discussion}
\label{sec:discussion}

The central result of this paper is that any causal quantum
environment coupled to gravity renders the gravitational
coupling complex:
$\kappa_{\mathrm{eff}}(k) = \mathrm{Re}\,\kappa_{\mathrm{eff}}
+ i\,\mathrm{Im}\,\kappa_{\mathrm{eff}}$,
with the imaginary part encoding dissipation into the environment
and the real part acquiring a dispersive correction constrained
by Kramers--Kronig causality.
This is a theorem---it follows from the influence-functional
formalism of Hu and Verdaguer~\cite{HuVerdaguer2004,HuVerdaguer2008}
applied to the Einstein--Langevin equation, and holds regardless
of the environment's microscopic structure.
What distinguishes our framework is the \emph{choice} of
environment: the Arithmetic-Langevin GFT (AL-GFT) noise kernel,
whose frequencies are set by the imaginary parts $\gamma_{n}$
of the non-trivial Riemann zeta zeros.

In this section we place the results in context, develop the
connection to the Berry--Keating programme, and assess the
broader implications for quantum gravity, number theory, and
observational cosmology.


%% -------------------------------------------------------------
\subsection{The Berry--Keating connection}
\label{subsec:BK-connection}

\subsubsection{From spectral conjecture to physical environment}

The Hilbert--P\'olya conjecture asserts that the non-trivial
zeros of $\zeta(s)$ are eigenvalues of a self-adjoint operator
on a suitable Hilbert space~\cite{HilbertPolya}.
Berry and Keating~\cite{BerryKeating1999} sharpened this to
a semiclassical conjecture: the operator should quantise the
classical Hamiltonian $H = xp$, whose periodic orbits have
actions $S_{p} = \hbar\ln p$ labelled by the primes.
Bender, Brody, and M\"uller~\cite{Bender2017BKHamiltonian}
constructed a $\mathcal{PT}$-symmetric realisation
(Eq.~\eqref{eq:BK-hamiltonian}), and
Sierra~\cite{Sierra2014Rindler} embedded the zeros as discrete
eigenvalues of a Dirac fermion in Rindler spacetime subject to
delta-function potentials at positions $\ln p$.

Our framework adds a new layer to this programme.
Rather than seeking a single Hamiltonian whose eigenvalues
\emph{are} the $\gamma_{n}$, we use the $\gamma_{n}$ as
input data for a physical environment and ask what
\emph{observable consequences} follow.
The Zeta-Comb environment can therefore be viewed as a
\emph{spectral probe}: it translates the abstract question
``Are the $\gamma_{n}$ eigenvalues of a physical operator?''
into the concrete question ``Does the gravitational-wave
background carry the GUE correlations expected of such
eigenvalues?''

\subsubsection{Trace-formula interpretation}

The Gutzwiller trace formula connects the density of states
of a quantum-chaotic system to a sum over classical periodic
orbits~\cite{Gutzwiller1971}:
%
\begin{equation}
  d(E)
  \;=\;
  \bar{d}(E)
  \;+\;
  \frac{1}{\pi}\,\mathrm{Re}
  \sum_{\mathrm{p.o.}}
  A_{p}\,e^{iS_{p}/\hbar}\,.
  \label{eq:gutzwiller}
\end{equation}
%
For the Berry--Keating Hamiltonian, the primitive periodic
orbits are labelled by primes $p$, with actions
$S_{p} = \hbar\ln p$ and periods $T_{p} = \ln p$.
The oscillatory part of $d(E)$ then reproduces the explicit
formula of Riemann--von~Mangoldt for the zero-counting
function $N(T)$~\cite{Edwards1974}.

In our framework, an analogous structure appears in the
noise kernel.
The Zeta-Comb modulation~\eqref{eq:modulation-recall} is a
sum over zeros weighted by Gaussian-damped amplitudes:
%
\begin{equation}
  \mathcal{M}(k)
  \;\sim\;
  1 \;+\;
  \sum_{n}\,e^{-\sigma\gamma_{n}^{2}}\,
  e^{i\gamma_{n}\ln(k/k_{\star})}\,,
  \label{eq:noise-trace}
\end{equation}
%
which is the Fourier dual of the Gutzwiller sum: where the
trace formula sums over \emph{orbits} (primes) to produce a
density of \emph{states} (zeros), our noise kernel sums over
\emph{zeros} to produce an observable \emph{spectrum} in
momentum space.
The Gaussian damping $e^{-\sigma\gamma_{n}^{2}}$ plays the
role of the semiclassical amplitude $A_{p}$, controlling how
many terms contribute.
The duality
%
\begin{equation}
  \text{Gutzwiller: }
  \sum_{\text{primes}} \to \text{zeros}
  \qquad\longleftrightarrow\qquad
  \text{Zeta-Comb: }
  \sum_{\text{zeros}} \to \text{observables}
  \label{eq:duality}
\end{equation}
%
closes the loop: the explicit formula encodes primes into zeros;
the Zeta-Comb decodes zeros into gravitational-wave signatures.
This is the sense in which our framework makes the
Berry--Keating programme \emph{empirically accessible}---not
by proving the Riemann Hypothesis, but by creating a physical
context in which its consequences could be measured.


%% -------------------------------------------------------------
\subsection{Rodgers' theorem and the arithmetic content of the
            beat spectrum}
\label{subsec:rodgers-arithmetic}

The GUE cross-check developed in \S\ref{sec:GUE} acquires
additional significance through the theorem of
Rodgers~\cite{Rodgers2018,Rodgers2024MMJ}.
Conditioned on the Riemann Hypothesis, Rodgers proved that the
GUE Conjecture for zeta zeros is \emph{logically equivalent} to
a precise statement about the distribution of primes: namely,
the asymptotic evaluation of correlations among products of
primes drawn from short intervals.

In our context, this equivalence has a striking physical
translation.
The beat spectrum $\Delta\gamma_{mn} = |\gamma_{m}-\gamma_{n}|$,
which we proposed as a binary GUE diagnostic
(\S\ref{sec:GUE}), is not merely a statistical test---it is,
via Rodgers' theorem, a \emph{direct probe of prime
correlations}.
If the beat spectrum measured in a gravitational-wave
background exhibits GUE level repulsion ($R_{2}(u) \to 0$ as
$u \to 0$), Rodgers' theorem guarantees that the underlying
zero spacings encode the Hardy--Littlewood predictions for
prime $k$-tuples.
Conversely, detecting Poisson statistics ($R_{2}(u) \to 1$)
would refute the GUE Conjecture and, by Rodgers' equivalence,
would imply that prime counts in short intervals violate
their conjectured variance.

This chain of implications---gravitational-wave data
$\to$ beat spectrum $\to$ GUE/Poisson $\to$ prime
distribution---is, to our knowledge, the first proposed
pathway from an astrophysical measurement to a statement
in analytic number theory.


%% -------------------------------------------------------------
\subsection{Relation to stochastic gravity and open quantum
            systems}
\label{subsec:stochastic-gravity-context}

The result that gravity becomes an open quantum system when
coupled to a quantum environment is not new; it is the
foundation of stochastic semiclassical
gravity~\cite{HuVerdaguer2004,HuVerdaguer2008,CalzettaHu2008}.
The Einstein--Langevin equation adds a stochastic source
$\xi_{\mu\nu}$ to the semiclassical Einstein equation, with
the noise kernel $N_{\mu\nu\alpha\beta}$ determined by the
vacuum fluctuations of the stress-energy tensor of quantum
matter fields.

Our contribution is not to the formalism itself but to the
\emph{choice of environment}.
Standard applications of stochastic gravity use free-field
or thermal environments, producing noise kernels that are
smooth in momentum space.
The AL-GFT environment, by contrast, produces a noise kernel
with discrete oscillatory structure---the
Zeta-Comb---inherited from the arithmetic properties of the
GFT microscopic degrees of freedom.
This discrete structure is what generates the complex
$\kappa_{\mathrm{eff}}$ with observable signatures, rather
than a featureless renormalisation of Newton's constant.

Recent work on gravitational decoherence of quantum
superpositions~\cite{Blencowe2013,Kanno2024} and on
gravitationally-induced entanglement
degradation~\cite{Anastopoulos2021} confirms the physical
relevance of treating gravity as an open system.
Our framework extends this programme to the inflationary
epoch, where the environment is not thermal matter but the
discrete quantum-geometric structure of spacetime itself.


%% -------------------------------------------------------------
\subsection{Falsifiability and the role of the sweet spot}
\label{subsec:falsifiability}

A theory that cannot fail is not a theory.
The framework developed here makes three classes of
falsifiable predictions, ordered by observational horizon:

\begin{enumerate}
  \item \textbf{Near-term (Planck re-analysis, 2026--2028).}
    Log-periodic oscillations in the CMB power spectrum at
    frequencies $\gamma_{n}\!/\!(2\pi)$.
    The matched-filter search uses \emph{no free parameters}
    beyond $(\varepsilon,\sigma)$ in the allowed range; the
    frequencies themselves are fixed by number theory.
    A null result at $\varepsilon \gtrsim 0.01$,
    $\sigma \sim 10^{-3}$ would shrink the viable parameter
    space.

  \item \textbf{Medium-term (BBO/DECIGO, 2030s--2040s).}
    A stochastic gravitational-wave background with
    amplitude $|\delta\kappa/\kappa| \sim 7 \times 10^{-5}$
    at the sweet spot
    $(\varepsilon,\sigma) = (10^{-2},\, 10^{-3})$.
    The GUE cross-check---level repulsion in the beat
    spectrum---is a binary yes/no test requiring no fitting.
    Detection of Poisson statistics would falsify the
    Zeta-Comb hypothesis while leaving the general
    complex-$\kappa$ theorem intact.

  \item \textbf{Long-term (self-consistent bootstrap).}
    The nonlinear eigenvalue problem of
    \S\ref{sec:bootstrap} either admits a non-trivial
    $\sigma^{\star}$ in the observationally allowed window
    or it does not.
    If it does, $(\varepsilon,\sigma)$ are no longer free
    parameters---they are predictions.
    If it does not, the Zeta-Comb environment decouples and
    the framework reduces to standard stochastic gravity
    with no arithmetic content.
\end{enumerate}

The existence of the sweet spot at
$\sigma \sim 10^{-3}$---where the beat-frequency visibility
threshold ($\sigma \lesssim 0.002$) overlaps with the
BBO/DECIGO sensitivity band---was not engineered.
It emerged from the structure of the first Riemann zeros
($\gamma_{1} = 14.1347$, $\gamma_{2} = 21.022$) and the
Gaussian weight function $e^{-\sigma\gamma_{n}^{2}}$.
This coincidence, discussed in \S\ref{sec:parameter-space},
constitutes a form of \emph{theoretical naturalness}: the
region where the framework is most testable is also where
its signatures are strongest.


%% -------------------------------------------------------------
\subsection{Limitations and caveats}
\label{subsec:limitations}

Several limitations should be stated plainly.

\paragraph{Perturbative regime.}
All results assume $\varepsilon \ll 1$ and
$|\kappa\,\widetilde{D}_{R}| \ll |\mathcal{O}|$,
i.e.\ the environment is a weak perturbation on the
gravitational dynamics.
The Kramers--Kronig analysis, the linearised FDT, and
the signal-amplitude estimates all break down if these
conditions are violated.
For $\varepsilon \gtrsim 0.1$, higher-order corrections
to the influence functional---including non-Gaussian
terms in the noise---would need to be included, taking
the theory beyond the Gaussian AL-GFT branch (Gate~1).

\paragraph{Environment specification.}
The Zeta-Comb is motivated by the AL-GFT microscopic model
(\S\ref{sec:intro}; Gate~1 of the CEQG-RG-Langevin
Blueprint), but the mapping from GFT microscopic
couplings to the effective parameters
$(\varepsilon,\sigma)$ involves approximations---notably
the linear-coupling and Gaussian-environment
assumptions---that may receive corrections from the
full GFT dynamics.
Gate~2 (Wetterich RG flow from $k = M_{\mathrm{Pl}}$
to $k = H_{0}$) will address this by promoting
$\sigma$ to a running coupling $\sigma(k)$.

\paragraph{Cosmological model dependence.}
The observational signatures are computed on an FLRW
background with Starobinsky-like inflation.
Alternative inflationary models change the transfer
function from primordial noise to CMB observables and
could shift the optimal frequency window.
However, the \emph{existence} of log-periodic oscillations
at Riemann-zero frequencies is model-independent; only
their amplitude is affected.

\paragraph{No proof of the Riemann Hypothesis.}
Nothing in this paper proves or assumes the Riemann
Hypothesis.
The $\gamma_{n}$ are used as empirical input---the
first $\sim 10^{13}$ zeros have been verified
numerically~\cite{Platt2021,Odlyzko1987}.
The GUE cross-check tests whether the \emph{measured}
beat spectrum matches the GUE prediction, which is a
contingent physical question, not a mathematical proof.
However, if the self-consistent bootstrap
(\S\ref{sec:bootstrap}) were shown to close only when
the zeros satisfy RH, this would constitute a
\emph{physical} argument for the hypothesis---one that
would need to be sharpened into a rigorous proof by
number-theoretic methods.


%% -------------------------------------------------------------
\subsection{Broader implications}
\label{subsec:broader-implications}

\subsubsection{For quantum gravity phenomenology}

The framework demonstrates that the open-quantum-systems
perspective on gravity is not merely a formal tool but a
\emph{phenomenological programme} with testable
consequences.
The key enabling insight is that the noise kernel of
stochastic gravity, usually treated as smooth and
unstructured, can carry discrete information from the
UV completion of gravity.
Any approach to quantum gravity that predicts a specific
noise spectrum---whether from loop quantum gravity,
string theory, or causal set theory---could in principle
be tested by the same methods developed here.
The general complex-$\kappa$ theorem
(\S\ref{sec:complex-kappa}) is agnostic about the
environment's origin; the Zeta-Comb is one instantiation,
but the framework is a template.

\subsubsection{For the Langlands programme and arithmetic
              quantum field theory}

The AL-GFT construction places automorphic data
(specifically, the $\gamma_{n}$ and their pair
correlations) inside a physical noise kernel.
If the non-Gaussian extension---Modular-Langevin GFT
(ML-GFT)---proves consistent, it would promote this
to a role for automorphic representations (Maass forms,
Hecke operators) in the bispectrum.
This aligns with the broader vision of the Langlands
programme as a unifying structure not only for
mathematics but for the symmetries of physical
theories~\cite{Langlands1970,Frenkel2007}.

\subsubsection{For observational cosmology}

The practical deliverable is a new class of
\emph{template waveforms} for CMB and gravitational-wave
data analysis: log-periodic oscillations at
number-theoretically determined frequencies.
Unlike generic oscillatory features (which require
fitting both amplitude and frequency), the Zeta-Comb
frequencies are fixed, reducing the search to a
one-parameter family in $\sigma$.
This dramatically reduces the look-elsewhere effect
in template searches, improving the effective
signal-to-noise ratio for a given observation time.
We anticipate that purpose-built matched-filter
pipelines, applied to Planck residual maps and future
BBO/DECIGO data, will provide the first empirical
constraints on arithmetic quantum gravity noise.


%% =============================================================
%% Closing
%% =============================================================

\bigskip
\noindent
\emph{Whether the cosmos confirms or refutes this arithmetic
structure, the question itself---Do the Riemann zeros
shape the fabric of spacetime?---is now empirically
meaningful.
The answer lies in the data.}
th, and $\varepsilon$
the multiplicity coupling strength.

The Kramers--Kronig relation links $D_{R}$ to $N$ via
a Hilbert transform (\S\ref{sec:FDT}), so the self-energy entering
\eqref{eq:kappa-eff-recall} is itself a functional of
$(\varepsilon,\sigma,\{\gamma_{n}\})$.
Schematically,
%
\begin{equation}
  \kappa_{\mathrm{eff}}
  \;=\;
  \kappa_{\mathrm{eff}}\!\bigl[\varepsilon,\,\sigma,\,
  \{\gamma_{n}\};\;\kappa\bigr]\,.
  \label{eq:kappa-functional}
\end{equation}
%
A \emph{self-consistent} theory demands that the environment producing
the noise be compatible with the gravitational dynamics it modifies.
Concretely, if the dressed coupling $\kappa_{\mathrm{eff}}$ back-reacts
on the environment---for instance by shifting the effective mass of
environment modes or by altering the background geometry on which they
propagate---then the environment parameters $(\varepsilon,\sigma)$ are
not free but must satisfy a fixed-point equation:
%
\begin{equation}
\boxed{%
  \kappa_{\mathrm{eff}}\!\bigl[\varepsilon(\sigma),\,\sigma,\,
  \{\gamma_{n}\};\;\kappa\bigr]
  \;\stackrel{!}{=}\;
  \kappa_{\mathrm{phys}}\,,
}
  \label{eq:bootstrap-condition}
\end{equation}
%
where $\kappa_{\mathrm{phys}} = \sqrt{16\pi G_{N}}$ is the
physically measured gravitational coupling and the notation
$\varepsilon(\sigma)$ emphasises that, once the bootstrap
closes, $\varepsilon$ becomes a derived function of $\sigma$.


%% -------------------------------------------------------------
\subsection{Nonlinear eigenvalue formulation}
\label{subsec:nonlinear-eigenvalue}

To make the bootstrap condition~\eqref{eq:bootstrap-condition}
tractable, we truncate the Zeta-Comb to the first $N$ zeros and
define the \emph{weight vector}
%
\begin{equation}
  v_{n}(\sigma) \;\equiv\; e^{-2\sigma\gamma_{n}^{2}}\,,
  \qquad n = 1,\ldots,N\,,
  \label{eq:weight-vector}
\end{equation}
%
which governs the contribution of each zero to the observable
signal (cf.\ Eq.~\eqref{eq:signal-amplitude}).
The self-consistency condition then takes the matrix form
%
\begin{equation}
  \mathbf{K}(\sigma)\,\mathbf{v}(\sigma) \;=\; \mathbf{0}\,,
  \label{eq:bootstrap-matrix}
\end{equation}
%
where $\mathbf{K}(\sigma)$ is an $N\times N$ matrix constructed from
the dressed propagator, the Kramers--Kronig--transformed noise kernel,
and the Ward-identity constraints derived in \S\ref{sec:conservation}.
Explicitly, the matrix elements are
%
\begin{equation}
  K_{mn}(\sigma)
  \;=\;
  \delta_{mn}
  \;-\;
  \varepsilon^{2}\,v_{m}(\sigma)\,
  \Bigl[
    \mathcal{H}_{mn}\!\bigl(\{\gamma\}\bigr)
    \;+\;
    i\,\mathcal{R}_{mn}\!\bigl(\{\gamma\}\bigr)
  \Bigr]\,,
  \label{eq:K-matrix-elements}
\end{equation}
%
where $\mathcal{H}_{mn}$ encodes the Hilbert-transform couplings
between zeros $m$ and $n$ (arising from the dissipation kernel), and
$\mathcal{R}_{mn}$ captures the reactive part constrained by the
fluctuation--dissipation theorem.

Equation~\eqref{eq:bootstrap-matrix} is a \emph{nonlinear eigenvalue
problem} (NLEP) in the parameter $\sigma$: the matrix $\mathbf{K}$
depends on $\sigma$ both through the weight vector $\mathbf{v}(\sigma)$
and through the $\sigma$-dependence of the Hilbert-transform integrals.
The problem is to find values $\sigma^{\star}$ such that
$\mathbf{K}(\sigma^{\star})$ has a non-trivial null space:
%
\begin{equation}
  \det\,\mathbf{K}(\sigma^{\star}) \;=\; 0\,.
  \label{eq:det-condition}
\end{equation}
%
This is structurally analogous to the gap equation in BCS
superconductivity, where the pairing gap $\Delta$ appears both in the
quasiparticle spectrum and in the self-energy that generates $\Delta$
itself.  Here, $\sigma$ plays the role of $\Delta$: it sets the
effective number of contributing zeros ($N_{\mathrm{eff}}$), which in
turn determines the noise kernel, which feeds back into the condition
on~$\sigma$.


%% -------------------------------------------------------------
\subsection{Structure of the solution space}
\label{subsec:solution-space}

Several properties of the NLEP~\eqref{eq:bootstrap-matrix} can be
established without solving it explicitly.

\paragraph{Existence of a trivial fixed point.}
At $\varepsilon = 0$, the matrix $\mathbf{K}$ reduces to the identity
and $\det\mathbf{K} = 1 \neq 0$.
The trivial solution $\varepsilon = 0$ (decoupled environment)
always satisfies \eqref{eq:bootstrap-condition} but produces no
observable signal.
A non-trivial bootstrap therefore requires $\varepsilon \neq 0$,
and the question is whether \eqref{eq:det-condition} admits a
solution at finite $\varepsilon$.

\paragraph{Monotonicity of the weight vector.}
Since $\gamma_{n+1} > \gamma_{n}$ and $\sigma > 0$, the weights
$v_{n}(\sigma)$ form a strictly decreasing sequence:
%
\begin{equation}
  v_{1}(\sigma) \;>\; v_{2}(\sigma) \;>\; \cdots \;>\; v_{N}(\sigma)
  \;>\; 0\,.
  \label{eq:v-monotone}
\end{equation}
%
This ensures that the lowest zeros ($\gamma_{1} = 14.1347\ldots$,
$\gamma_{2} = 21.0220\ldots$) dominate the bootstrap,
providing a natural truncation hierarchy controlled by $\sigma$.
At the sweet spot $\sigma \sim 10^{-3}$, the effective number of
zeros contributing is $N_{\mathrm{eff}} \approx 4.2$,
with 7 pairwise beat frequencies (Table~\ref{tab:sweet-spot}).

\paragraph{Spectral flow with $\sigma$.}
Define $\lambda_{1}(\sigma) \equiv$ smallest singular value of
$\mathbf{K}(\sigma)$.
As $\sigma \to 0$, all weights $v_{n} \to 1$ and the off-diagonal
couplings $\mathcal{H}_{mn}$ grow without bound (too many zeros
contribute), driving $\lambda_{1} \to -\infty$.
As $\sigma \to \infty$, all $v_{n} \to 0$ and
$\mathbf{K} \to \mathbf{I}$, so $\lambda_{1} \to 1$.
By continuity, $\lambda_{1}(\sigma^{\star}) = 0$ for at least one
$\sigma^{\star} \in (0,\infty)$, provided the off-diagonal couplings
are sufficiently strong.
This \emph{intermediate-value argument} does not guarantee uniqueness,
but it provides a topological existence proof for the non-trivial
bootstrap.


%% -------------------------------------------------------------
\subsection{Connection to the Berry--Keating programme}
\label{subsec:berry-keating}

The nonlinear eigenvalue problem~\eqref{eq:bootstrap-matrix}
acquires additional structure when viewed through the lens
of the Berry--Keating conjecture~\cite{BerryKeating1999,
Bender2017BKHamiltonian}.
Berry and Keating proposed that the Riemann zeros are eigenvalues
of a quantisation of the classical Hamiltonian $H = xp$, with a
suitable boundary condition encoding the prime numbers.
Bender, Brody, and M\"uller~\cite{Bender2017BKHamiltonian}
subsequently constructed a $\mathcal{PT}$-symmetric Hamiltonian
%
\begin{equation}
  \hat{H}_{\mathrm{BK}}
  \;=\;
  \frac{1}{1 - e^{-i\hat{p}}}
  \bigl(\hat{x}\hat{p} + \hat{p}\hat{x}\bigr)
  \frac{1}{1 - e^{-i\hat{p}}}\,,
  \label{eq:BK-hamiltonian}
\end{equation}
%
whose eigenvalues, subject to a regularity condition, coincide with
the $\gamma_{n}$.

In our framework, the Zeta-Comb environment is \emph{by construction}
built from the $\gamma_{n}$.
If a Berry--Keating--type Hamiltonian governs the environment modes,
then the self-consistency condition~\eqref{eq:bootstrap-matrix}
becomes a statement about the spectral properties of
$\hat{H}_{\mathrm{BK}}$:
\emph{the gravitational bootstrap closes if and only if the
environment spectrum is that of a self-adjoint (or
$\mathcal{PT}$-symmetric) operator whose eigenvalues are
the Riemann zeros.}
This establishes a physical bridge between the Riemann Hypothesis
and quantum gravity: the bootstrap has a non-trivial solution
$\sigma^{\star}$ precisely when the $\gamma_{n}$ satisfy
the spectral rigidity imposed by the GUE universality class---a
known consequence of RH via Montgomery--Odlyzko pair
correlation~\cite{Montgomery1973,Odlyzko1987}.

Conversely, if the bootstrap were observed to close experimentally
(via the gravitational-wave signatures of
\S\ref{sec:parameter-space}), it would constitute indirect
\emph{physical evidence} for the spectral interpretation of the
Riemann zeros.


%% -------------------------------------------------------------
\subsection{Numerical strategy and convergence}
\label{subsec:numerical-strategy}

We outline a concrete algorithm for solving the
NLEP~\eqref{eq:bootstrap-matrix} numerically.

\begin{enumerate}
  \item \textbf{Truncation.}
    Fix $N_{\max}$ zeros (we recommend $N_{\max} = 50$, well
    within Odlyzko's tabulated range).
    For $\sigma \gtrsim 10^{-4}$, the Gaussian suppression
    $v_{n} \sim e^{-2\sigma\gamma_{n}^{2}}$ renders contributions
    from $n > N_{\max}$ negligible ($v_{50} < 10^{-12}$ at
    $\sigma = 10^{-3}$).

  \item \textbf{Matrix assembly.}
    For each trial $\sigma$, compute $\mathbf{v}(\sigma)$ from
    \eqref{eq:weight-vector} and assemble $\mathbf{K}(\sigma)$
    via numerical Hilbert transforms of the noise kernel restricted
    to the first $N_{\max}$ zeros.
    The Ward-identity constraint
    ($k^{\mu}\,\mathrm{Im}\,\kappa_{\mathrm{eff}} = 0$;
    \S\ref{sec:conservation}) is imposed as a
    projection onto the transverse-traceless subspace at each step.

  \item \textbf{Determinant scan.}
    Evaluate $\det\mathbf{K}(\sigma)$ on a logarithmic grid
    $\sigma \in [10^{-5},\, 10^{-1}]$.
    Sign changes in $\mathrm{Re}\,\det\mathbf{K}$ locate
    candidate roots $\sigma^{\star}$.

  \item \textbf{Refinement.}
    Polish each candidate using a Newton--Raphson iteration on
    the smallest singular value $\lambda_{1}(\sigma)$,
    with Jacobian estimated by finite differences
    $\partial_{\sigma}\lambda_{1} \approx
    [\lambda_{1}(\sigma+\delta) - \lambda_{1}(\sigma-\delta)]
    / 2\delta$.

  \item \textbf{Consistency check.}
    Verify that $\sigma^{\star}$ lies within the observationally
    allowed region of parameter space
    (Fig.~\ref{fig:parameter-space}):
    beat-frequency visibility requires $\sigma \lesssim 0.002$,
    and Planck 2018 constraints impose
    $\varepsilon < 0.21$ at $\sigma = 10^{-3}$
    (Eq.~\eqref{eq:planck-bound}).
    A solution $\sigma^{\star} \sim 10^{-3}$ falling in the
    sweet spot would be a striking confirmation of the
    framework's internal coherence.
\end{enumerate}

The computational cost is dominated by the Hilbert-transform
integrals in Step~2, each requiring $\mathcal{O}(N_{\max}^{2})$
quadrature evaluations.
For $N_{\max} = 50$, a single $\sigma$-scan with 200 grid points
requires $\sim 5 \times 10^{5}$ function evaluations---feasible
on a laptop in minutes.


%% -------------------------------------------------------------
\subsection{Open questions and future work}
\label{subsec:open-questions}

The self-consistent bootstrap raises several questions that
we leave for future investigation.

\begin{enumerate}
  \item \textbf{Uniqueness.}
    Does \eqref{eq:det-condition} admit a unique non-trivial
    $\sigma^{\star}$, or a discrete set?
    Multiple solutions would correspond to distinct
    ``phases'' of the arithmetic environment, each with a
    different effective number of active zeros.
    Stability analysis (sign of
    $\partial_{\sigma}\lambda_{1}|_{\sigma^{\star}}$) would
    identify the physical branch.

  \item \textbf{Dynamical origin of $\varepsilon$.}
    In the present formulation, $\varepsilon$ enters as an
    overall coupling strength.
    A fully closed bootstrap should determine
    $\varepsilon^{\star}$ alongside $\sigma^{\star}$.
    This likely requires extending the NLEP to a
    two-parameter system
    $\mathbf{K}(\varepsilon,\sigma)\,\mathbf{v} = \mathbf{0}$,
    supplemented by a normalisation condition on
    $\kappa_{\mathrm{phys}}$.

  \item \textbf{Relation to the Hilbert--P\'olya programme.}
    If the bootstrap selects a unique $\sigma^{\star}$, the
    resulting dressed propagator
    $G_{\mathrm{eff}}(k;\sigma^{\star})$ defines a
    \emph{physical operator} whose spectral properties are
    constrained by the Riemann zeros.
    Establishing whether this operator is unitarily equivalent
    to a Berry--Keating--type Hamiltonian would provide a
    concrete realisation of the Hilbert--P\'olya conjecture
    within quantum gravity.

  \item \textbf{Renormalisation-group flow.}
    The parameter $\sigma$ has been treated as scale-independent.
    Embedding the bootstrap in the Gate~2 Wetterich flow
    (\S\ref{sec:gate2-outlook}) would promote $\sigma$ to a
    running coupling $\sigma(k)$, with the NLEP replaced by a
    functional fixed-point equation.
    The interplay between the discrete Zeta-Comb structure
    and continuous RG flow is an unexplored frontier.

  \item \textbf{Implications of failure.}
    If no non-trivial $\sigma^{\star}$ exists within the
    observationally allowed window, the framework admits only
    the trivial solution $\varepsilon = 0$---meaning the
    arithmetic environment decouples and leaves no observable
    imprint.
    This would \emph{falsify} the predictive power of the
    Zeta-Comb model while leaving the general theorem
    (any causal environment produces complex $\kappa$;
    \S\ref{sec:complex-kappa}) intact.
\end{enumerate}

\noindent
The self-consistent bootstrap thus serves a dual purpose:
it is both the mechanism by which the framework could achieve
maximal predictive power---fixing all parameters from number
theory---and the sharpest internal test of whether arithmetic
quantum gravity noise is more than a mathematical curiosity.
Its resolution, whether by the numerical programme above or
by analytic methods connecting to the Berry--Keating Hamiltonian,
constitutes the natural next chapter of this work.

%% =============================================================
%% §8  Discussion
%% Paper: "Complex Gravitational Coupling from Arithmetic Quantum Gravity Noise"
%% =============================================================

\section{Discussion}
\label{sec:discussion}

The central result of this paper is that any causal quantum
environment coupled to gravity renders the gravitational
coupling complex:
$\kappa_{\mathrm{eff}}(k) = \mathrm{Re}\,\kappa_{\mathrm{eff}}
+ i\,\mathrm{Im}\,\kappa_{\mathrm{eff}}$,
with the imaginary part encoding dissipation into the environment
and the real part acquiring a dispersive correction constrained
by Kramers--Kronig causality.
This is a theorem---it follows from the influence-functional
formalism of Hu and Verdaguer~\cite{HuVerdaguer2004,HuVerdaguer2008}
applied to the Einstein--Langevin equation, and holds regardless
of the environment's microscopic structure.
What distinguishes our framework is the \emph{choice} of
environment: the Arithmetic-Langevin GFT (AL-GFT) noise kernel,
whose frequencies are set by the imaginary parts $\gamma_{n}$
of the non-trivial Riemann zeta zeros.

In this section we place the results in context, develop the
connection to the Berry--Keating programme, and assess the
broader implications for quantum gravity, number theory, and
observational cosmology.


%% -------------------------------------------------------------
\subsection{The Berry--Keating connection}
\label{subsec:BK-connection}

\subsubsection{From spectral conjecture to physical environment}

The Hilbert--P\'olya conjecture asserts that the non-trivial
zeros of $\zeta(s)$ are eigenvalues of a self-adjoint operator
on a suitable Hilbert space~\cite{HilbertPolya}.
Berry and Keating~\cite{BerryKeating1999} sharpened this to
a semiclassical conjecture: the operator should quantise the
classical Hamiltonian $H = xp$, whose periodic orbits have
actions $S_{p} = \hbar\ln p$ labelled by the primes.
Bender, Brody, and M\"uller~\cite{Bender2017BKHamiltonian}
constructed a $\mathcal{PT}$-symmetric realisation
(Eq.~\eqref{eq:BK-hamiltonian}), and
Sierra~\cite{Sierra2014Rindler} embedded the zeros as discrete
eigenvalues of a Dirac fermion in Rindler spacetime subject to
delta-function potentials at positions $\ln p$.

Our framework adds a new layer to this programme.
Rather than seeking a single Hamiltonian whose eigenvalues
\emph{are} the $\gamma_{n}$, we use the $\gamma_{n}$ as
input data for a physical environment and ask what
\emph{observable consequences} follow.
The Zeta-Comb environment can therefore be viewed as a
\emph{spectral probe}: it translates the abstract question
``Are the $\gamma_{n}$ eigenvalues of a physical operator?''
into the concrete question ``Does the gravitational-wave
background carry the GUE correlations expected of such
eigenvalues?''

\subsubsection{Trace-formula interpretation}

The Gutzwiller trace formula connects the density of states
of a quantum-chaotic system to a sum over classical periodic
orbits~\cite{Gutzwiller1971}:
%
\begin{equation}
  d(E)
  \;=\;
  \bar{d}(E)
  \;+\;
  \frac{1}{\pi}\,\mathrm{Re}
  \sum_{\mathrm{p.o.}}
  A_{p}\,e^{iS_{p}/\hbar}\,.
  \label{eq:gutzwiller}
\end{equation}
%
For the Berry--Keating Hamiltonian, the primitive periodic
orbits are labelled by primes $p$, with actions
$S_{p} = \hbar\ln p$ and periods $T_{p} = \ln p$.
The oscillatory part of $d(E)$ then reproduces the explicit
formula of Riemann--von~Mangoldt for the zero-counting
function $N(T)$~\cite{Edwards1974}.

In our framework, an analogous structure appears in the
noise kernel.
The Zeta-Comb modulation~\eqref{eq:modulation-recall} is a
sum over zeros weighted by Gaussian-damped amplitudes:
%
\begin{equation}
  \mathcal{M}(k)
  \;\sim\;
  1 \;+\;
  \sum_{n}\,e^{-\sigma\gamma_{n}^{2}}\,
  e^{i\gamma_{n}\ln(k/k_{\star})}\,,
  \label{eq:noise-trace}
\end{equation}
%
which is the Fourier dual of the Gutzwiller sum: where the
trace formula sums over \emph{orbits} (primes) to produce a
density of \emph{states} (zeros), our noise kernel sums over
\emph{zeros} to produce an observable \emph{spectrum} in
momentum space.
The Gaussian damping $e^{-\sigma\gamma_{n}^{2}}$ plays the
role of the semiclassical amplitude $A_{p}$, controlling how
many terms contribute.
The duality
%
\begin{equation}
  \text{Gutzwiller: }
  \sum_{\text{primes}} \to \text{zeros}
  \qquad\longleftrightarrow\qquad
  \text{Zeta-Comb: }
  \sum_{\text{zeros}} \to \text{observables}
  \label{eq:duality}
\end{equation}
%
closes the loop: the explicit formula encodes primes into zeros;
the Zeta-Comb decodes zeros into gravitational-wave signatures.
This is the sense in which our framework makes the
Berry--Keating programme \emph{empirically accessible}---not
by proving the Riemann Hypothesis, but by creating a physical
context in which its consequences could be measured.


%% -------------------------------------------------------------
\subsection{Rodgers' theorem and the arithmetic content of the
            beat spectrum}
\label{subsec:rodgers-arithmetic}

The GUE cross-check developed in \S\ref{sec:GUE} acquires
additional significance through the theorem of
Rodgers~\cite{Rodgers2018,Rodgers2024MMJ}.
Conditioned on the Riemann Hypothesis, Rodgers proved that the
GUE Conjecture for zeta zeros is \emph{logically equivalent} to
a precise statement about the distribution of primes: namely,
the asymptotic evaluation of correlations among products of
primes drawn from short intervals.

In our context, this equivalence has a striking physical
translation.
The beat spectrum $\Delta\gamma_{mn} = |\gamma_{m}-\gamma_{n}|$,
which we proposed as a binary GUE diagnostic
(\S\ref{sec:GUE}), is not merely a statistical test---it is,
via Rodgers' theorem, a \emph{direct probe of prime
correlations}.
If the beat spectrum measured in a gravitational-wave
background exhibits GUE level repulsion ($R_{2}(u) \to 0$ as
$u \to 0$), Rodgers' theorem guarantees that the underlying
zero spacings encode the Hardy--Littlewood predictions for
prime $k$-tuples.
Conversely, detecting Poisson statistics ($R_{2}(u) \to 1$)
would refute the GUE Conjecture and, by Rodgers' equivalence,
would imply that prime counts in short intervals violate
their conjectured variance.

This chain of implications---gravitational-wave data
$\to$ beat spectrum $\to$ GUE/Poisson $\to$ prime
distribution---is, to our knowledge, the first proposed
pathway from an astrophysical measurement to a statement
in analytic number theory.


%% -------------------------------------------------------------
\subsection{Relation to stochastic gravity and open quantum
            systems}
\label{subsec:stochastic-gravity-context}

The result that gravity becomes an open quantum system when
coupled to a quantum environment is not new; it is the
foundation of stochastic semiclassical
gravity~\cite{HuVerdaguer2004,HuVerdaguer2008,CalzettaHu2008}.
The Einstein--Langevin equation adds a stochastic source
$\xi_{\mu\nu}$ to the semiclassical Einstein equation, with
the noise kernel $N_{\mu\nu\alpha\beta}$ determined by the
vacuum fluctuations of the stress-energy tensor of quantum
matter fields.

Our contribution is not to the formalism itself but to the
\emph{choice of environment}.
Standard applications of stochastic gravity use free-field
or thermal environments, producing noise kernels that are
smooth in momentum space.
The AL-GFT environment, by contrast, produces a noise kernel
with discrete oscillatory structure---the
Zeta-Comb---inherited from the arithmetic properties of the
GFT microscopic degrees of freedom.
This discrete structure is what generates the complex
$\kappa_{\mathrm{eff}}$ with observable signatures, rather
than a featureless renormalisation of Newton's constant.

Recent work on gravitational decoherence of quantum
superpositions~\cite{Blencowe2013,Kanno2024} and on
gravitationally-induced entanglement
degradation~\cite{Anastopoulos2021} confirms the physical
relevance of treating gravity as an open system.
Our framework extends this programme to the inflationary
epoch, where the environment is not thermal matter but the
discrete quantum-geometric structure of spacetime itself.


%% -------------------------------------------------------------
\subsection{Falsifiability and the role of the sweet spot}
\label{subsec:falsifiability}

A theory that cannot fail is not a theory.
The framework developed here makes three classes of
falsifiable predictions, ordered by observational horizon:

\begin{enumerate}
  \item \textbf{Near-term (Planck re-analysis, 2026--2028).}
    Log-periodic oscillations in the CMB power spectrum at
    frequencies $\gamma_{n}\!/\!(2\pi)$.
    The matched-filter search uses \emph{no free parameters}
    beyond $(\varepsilon,\sigma)$ in the allowed range; the
    frequencies themselves are fixed by number theory.
    A null result at $\varepsilon \gtrsim 0.01$,
    $\sigma \sim 10^{-3}$ would shrink the viable parameter
    space.

  \item \textbf{Medium-term (BBO/DECIGO, 2030s--2040s).}
    A stochastic gravitational-wave background with
    amplitude $|\delta\kappa/\kappa| \sim 7 \times 10^{-5}$
    at the sweet spot
    $(\varepsilon,\sigma) = (10^{-2},\, 10^{-3})$.
    The GUE cross-check---level repulsion in the beat
    spectrum---is a binary yes/no test requiring no fitting.
    Detection of Poisson statistics would falsify the
    Zeta-Comb hypothesis while leaving the general
    complex-$\kappa$ theorem intact.

  \item \textbf{Long-term (self-consistent bootstrap).}
    The nonlinear eigenvalue problem of
    \S\ref{sec:bootstrap} either admits a non-trivial
    $\sigma^{\star}$ in the observationally allowed window
    or it does not.
    If it does, $(\varepsilon,\sigma)$ are no longer free
    parameters---they are predictions.
    If it does not, the Zeta-Comb environment decouples and
    the framework reduces to standard stochastic gravity
    with no arithmetic content.
\end{enumerate}

The existence of the sweet spot at
$\sigma \sim 10^{-3}$---where the beat-frequency visibility
threshold ($\sigma \lesssim 0.002$) overlaps with the
BBO/DECIGO sensitivity band---was not engineered.
It emerged from the structure of the first Riemann zeros
($\gamma_{1} = 14.1347$, $\gamma_{2} = 21.022$) and the
Gaussian weight function $e^{-\sigma\gamma_{n}^{2}}$.
This coincidence, discussed in \S\ref{sec:parameter-space},
constitutes a form of \emph{theoretical naturalness}: the
region where the framework is most testable is also where
its signatures are strongest.


%% -------------------------------------------------------------
\subsection{Limitations and caveats}
\label{subsec:limitations}

Several limitations should be stated plainly.

\paragraph{Perturbative regime.}
All results assume $\varepsilon \ll 1$ and
$|\kappa\,\widetilde{D}_{R}| \ll |\mathcal{O}|$,
i.e.\ the environment is a weak perturbation on the
gravitational dynamics.
The Kramers--Kronig analysis, the linearised FDT, and
the signal-amplitude estimates all break down if these
conditions are violated.
For $\varepsilon \gtrsim 0.1$, higher-order corrections
to the influence functional---including non-Gaussian
terms in the noise---would need to be included, taking
the theory beyond the Gaussian AL-GFT branch (Gate~1).

\paragraph{Environment specification.}
The Zeta-Comb is motivated by the AL-GFT microscopic model
(\S\ref{sec:intro}; Gate~1 of the CEQG-RG-Langevin
Blueprint), but the mapping from GFT microscopic
couplings to the effective parameters
$(\varepsilon,\sigma)$ involves approximations---notably
the linear-coupling and Gaussian-environment
assumptions---that may receive corrections from the
full GFT dynamics.
Gate~2 (Wetterich RG flow from $k = M_{\mathrm{Pl}}$
to $k = H_{0}$) will address this by promoting
$\sigma$ to a running coupling $\sigma(k)$.

\paragraph{Cosmological model dependence.}
The observational signatures are computed on an FLRW
background with Starobinsky-like inflation.
Alternative inflationary models change the transfer
function from primordial noise to CMB observables and
could shift the optimal frequency window.
However, the \emph{existence} of log-periodic oscillations
at Riemann-zero frequencies is model-independent; only
their amplitude is affected.

\paragraph{No proof of the Riemann Hypothesis.}
Nothing in this paper proves or assumes the Riemann
Hypothesis.
The $\gamma_{n}$ are used as empirical input---the
first $\sim 10^{13}$ zeros have been verified
numerically~\cite{Platt2021,Odlyzko1987}.
The GUE cross-check tests whether the \emph{measured}
beat spectrum matches the GUE prediction, which is a
contingent physical question, not a mathematical proof.
However, if the self-consistent bootstrap
(\S\ref{sec:bootstrap}) were shown to close only when
the zeros satisfy RH, this would constitute a
\emph{physical} argument for the hypothesis---one that
would need to be sharpened into a rigorous proof by
number-theoretic methods.


%% -------------------------------------------------------------
\subsection{Broader implications}
\label{subsec:broader-implications}

\subsubsection{For quantum gravity phenomenology}

The framework demonstrates that the open-quantum-systems
perspective on gravity is not merely a formal tool but a
\emph{phenomenological programme} with testable
consequences.
The key enabling insight is that the noise kernel of
stochastic gravity, usually treated as smooth and
unstructured, can carry discrete information from the
UV completion of gravity.
Any approach to quantum gravity that predicts a specific
noise spectrum---whether from loop quantum gravity,
string theory, or causal set theory---could in principle
be tested by the same methods developed here.
The general complex-$\kappa$ theorem
(\S\ref{sec:complex-kappa}) is agnostic about the
environment's origin; the Zeta-Comb is one instantiation,
but the framework is a template.

\subsubsection{For the Langlands programme and arithmetic
              quantum field theory}

The AL-GFT construction places automorphic data
(specifically, the $\gamma_{n}$ and their pair
correlations) inside a physical noise kernel.
If the non-Gaussian extension---Modular-Langevin GFT
(ML-GFT)---proves consistent, it would promote this
to a role for automorphic representations (Maass forms,
Hecke operators) in the bispectrum.
This aligns with the broader vision of the Langlands
programme as a unifying structure not only for
mathematics but for the symmetries of physical
theories~\cite{Langlands1970,Frenkel2007}.

\subsubsection{For observational cosmology}

The practical deliverable is a new class of
\emph{template waveforms} for CMB and gravitational-wave
data analysis: log-periodic oscillations at
number-theoretically determined frequencies.
Unlike generic oscillatory features (which require
fitting both amplitude and frequency), the Zeta-Comb
frequencies are fixed, reducing the search to a
one-parameter family in $\sigma$.
This dramatically reduces the look-elsewhere effect
in template searches, improving the effective
signal-to-noise ratio for a given observation time.
We anticipate that purpose-built matched-filter
pipelines, applied to Planck residual maps and future
BBO/DECIGO data, will provide the first empirical
constraints on arithmetic quantum gravity noise.


%% =============================================================
%% Closing
%% =============================================================

\bigskip
\noindent
\emph{Whether the cosmos confirms or refutes this arithmetic
structure, the question itself---Do the Riemann zeros
shape the fabric of spacetime?---is now empirically
meaningful.
The answer lies in the data.}

%%% References cited in this section (to be merged with master .bib):
% \cite{Montgomery1973}  — H. L. Montgomery, Proc. Symp. Pure Math. 24 (1973) 181.
% \cite{Odlyzko1987}     — A. M. Odlyzko, Math. Comp. 48 (1987) 273.
% \cite{Odlyzko2001}     — A. M. Odlyzko, "The 10^22-nd zero of the Riemann zeta function" (2001).
% \cite{BogomolnyKeating1996} — E. B. Bogomolny and J. P. Keating, Phys. Rev. Lett. 77 (1996) 1472.
% \cite{Bogomolny2006}   — E. B. Bogomolny et al., J. Phys. A 39 (2006) 10743.
% \cite{Mehta2004}        — M. L. Mehta, Random Matrices, 3rd ed. (Academic Press, 2004).
% \cite{Rodgers2018}      — B. Rodgers, "Arithmetic consequences of the GUE conjecture for zeta zeros."
% \cite{LMFDB}            — The LMFDB Collaboration, https://www.lmfdb.org.



%%---------------------------------------------------------------------
%% Additional bibliography entries for this section:
%%
%% @article{Montgomery1973,
%%   author  = {Montgomery, H. L.},
%%   title   = {The pair correlation of zeros of the zeta function},
%%   journal = {Proc.\ Symp.\ Pure Math.},
%%   volume  = {24},
%%   pages   = {181--193},
%%   year    = {1973},
%% }
%%
%% @article{Odlyzko1987,
%%   author  = {Odlyzko, A. M.},
%%   title   = {On the distribution of spacings between zeros of the
%%              zeta function},
%%   journal = {Math.\ Comp.},
%%   volume  = {48},
%%   pages   = {273--308},
%%   year    = {1987},
%% }
%%
%% @phdthesis{Rodgers2013,
%%   author  = {Rodgers, B.},
%%   title   = {The statistics of the zeros of the {R}iemann
%%              zeta-function and related topics},
%%   school  = {UCLA},
%%   year    = {2013},
%% }
%%
%% @article{BerryKeating1999,
%%   author  = {Berry, M. V. and Keating, J. P.},
%%   title   = {The {R}iemann zeros and eigenvalue asymptotics},
%%   journal = {SIAM Rev.},
%%   volume  = {41},
%%   pages   = {236--266},
%%   year    = {1999},
%% }
%%
%% @article{BenderBrodyMueller2017,
%%   author  = {Bender, C. M. and Brody, D. C. and M{\"u}ller, M. P.},
%%   title   = {Hamiltonian for the zeros of the {R}iemann zeta function},
%%   journal = {Phys.\ Rev.\ Lett.},
%%   volume  = {118},
%%   pages   = {130201},
%%   year    = {2017},
%%   eprint  = {1608.03679},
%% }
%%
%% @article{PlanckInflation2018,
%%   author  = {{Planck Collaboration}},
%%   title   = {{P}lanck 2018 results. {X}. {C}onstraints on inflation},
%%   journal = {Astron.\ Astrophys.},
%%   volume  = {641},
%%   pages   = {A10},
%%   year    = {2020},
%%   eprint  = {1807.06211},
%% }
%%
%% @article{RudnickSarnak1996,
%%   author  = {Rudnick, Z. and Sarnak, P.},
%%   title   = {Zeros of principal {$L$}-functions and random matrix
%%              theory},
%%   journal = {Duke Math.\ J.},
%%   volume  = {81},
%%   pages   = {269--322},
%%   year    = {1996},
%% }
%%---------------------------------------------------------------------



%%---------------------------------------------------------------------
%% Bibliography entries referenced in this section (BibTeX keys):
%%
%% @article{HuVerdaguer2008,
%%   author  = {Hu, B. L. and Verdaguer, E.},
%%   title   = {Stochastic Gravity: Theory and Applications},
%%   journal = {Living Rev. Relativ.},
%%   volume  = {11},
%%   pages   = {3},
%%   year    = {2008},
%%   eprint  = {0802.0658},
%% }
%%
%% @article{FeynmanVernon1963,
%%   author  = {Feynman, R. P. and Vernon, F. L.},
%%   title   = {The Theory of a General Quantum System Interacting
%%              with a Linear Dissipative System},
%%   journal = {Ann.\ Phys.},
%%   volume  = {24},
%%   pages   = {118--173},
%%   year    = {1963},
%% }
%%
%% @book{CalzettaHu2008,
%%   author    = {Calzetta, E. A. and Hu, B. L.},
%%   title     = {Nonequilibrium Quantum Field Theory},
%%   publisher = {Cambridge University Press},
%%   year      = {2008},
%% }
%%
%% @book{Titchmarsh1948,
%%   author    = {Titchmarsh, E. C.},
%%   title     = {Introduction to the Theory of Fourier Integrals},
%%   publisher = {Oxford University Press},
%%   edition   = {2nd},
%%   year      = {1948},
%% }
%%
%% @misc{KramersKronig,
%%   note = {See e.g.\ \cite{Nussenzveig1972} for a comprehensive
%%           treatment.},
%% }
%%
%% @book{Nussenzveig1972,
%%   author    = {Nussenzveig, H. M.},
%%   title     = {Causality and Dispersion Relations},
%%   publisher = {Academic Press},
%%   year      = {1972},
%% }
%%
%% @article{Oriti2012,
%%   author  = {Oriti, D.},
%%   title   = {Group field theory as the microscopic description of
%%              the quantum spacetime fluid},
%%   journal = {Proc.\ Sci.},
%%   year    = {2012},
%%   eprint  = {0710.3276},
%% }
%%---------------------------------------------------------------------


%% ——— Bibliography stub ———
\begin{thebibliography}{99}

\bibitem{HuVerdaguer2008}
B.~L.~Hu and E.~Verdaguer,
``Stochastic Gravity: Theory and Applications,''
Living Rev.\ Rel.\ \textbf{11}, 3 (2008)
[arXiv:0802.0658 [gr-qc]].

\bibitem{HuVerdaguer2004}
B.~L.~Hu and E.~Verdaguer,
``Stochastic Gravity: Theory and Applications,''
Living Rev.\ Rel.\ \textbf{7}, 3 (2004)
[arXiv:gr-qc/0307032].

\bibitem{CalzettaHu1994}
E.~Calzetta and B.~L.~Hu,
``Noise and fluctuations in semiclassical gravity,''
Phys.\ Rev.\ D \textbf{49}, 6636 (1994).

\bibitem{FeynmonVernon1963}
R.~P.~Feynman and F.~L.~Vernon,
``The theory of a general quantum system interacting with a linear dissipative system,''
Annals Phys.\ \textbf{24}, 118 (1963).

\bibitem{KramersKronig}
R.~de~L.~Kronig, ``On the theory of dispersion of X-rays,''
J.~Opt.\ Soc.\ Am.\ \textbf{12}, 547 (1926);
H.~A.~Kramers, ``La diffusion de la lumière par les atomes,''
Atti Cong.\ Intern.\ Fisici, Como (1927).

\bibitem{GibbonsHawkingPerry1978}
G.~W.~Gibbons, S.~W.~Hawking and M.~J.~Perry,
``Path Integrals and the Indefiniteness of the Gravitational Action,''
Nucl.\ Phys.\ B \textbf{138}, 141 (1978).

\bibitem{Mazur1990}
P.~O.~Mazur and E.~Mottola,
``The conformal factor and the cosmological constant,''
Nucl.\ Phys.\ B (Proc.\ Suppl.)\ \textbf{18}, 214 (1990).

\bibitem{ConformalInstability2022}
H.-X.~Feng,
``Local conformal instability and local non-collapsing in the Ricci flow,''
Annals Phys.\ \textbf{440}, 168831 (2022)
[arXiv:2201.10732 [gr-qc]].

\bibitem{GravPI_UCSB}
D.~Grabovsky,
``Gravitational Path Integrals,''
UCSB Lecture Notes (2023),
\url{https://web.physics.ucsb.edu/~davidgrabovsky/files-notes/231C Notes.pdf}.

\bibitem{Montgomery1973}
H.~L.~Montgomery,
``The pair correlation of zeros of the zeta function,''
in \emph{Analytic Number Theory}, Proc.\ Symp.\ Pure Math.\ \textbf{24}, 181--193 (1973).

\bibitem{Odlyzko1987}
A.~M.~Odlyzko,
``On the distribution of spacings between zeros of the zeta function,''
Math.\ Comp.\ \textbf{48}, 273 (1987).

\bibitem{MontgomeryWiki}
``Montgomery's pair correlation conjecture,''
Wikipedia (2024),
\url{https://en.wikipedia.org/wiki/Montgomery\%27s_pair_correlation_conjecture}.

\bibitem{MathWorldMontgomery}
E.~W.~Weisstein,
``Montgomery's Pair Correlation Conjecture,''
MathWorld,
\url{https://mathworld.wolfram.com/MontgomerysPairCorrelationConjecture.html}.

\bibitem{BerryKeating1999}
M.~V.~Berry and J.~P.~Keating,
``The Riemann zeros and eigenvalue asymptotics,''
SIAM Rev.\ \textbf{41}, 236 (1999).

\bibitem{BBM2017}
C.~M.~Bender, D.~C.~Brody and M.~P.~M\"uller,
``Hamiltonian for the Zeros of the Riemann Zeta Function,''
Phys.\ Rev.\ Lett.\ \textbf{118}, 130201 (2017)
[arXiv:1608.03679 [quant-ph]].

\bibitem{Rodgers2013}
B.~Rodgers,
``Arithmetic consequences of the GUE conjecture for zeta zeros,''
\url{http://user.math.uzh.ch/rodgers/ArithmeticGUE.pdf}.

\bibitem{ALGFTGate1}
R.~O.~Van~Gelder,
``Gate 1 Formalization: Arithmetic--Langevin Group Field Theory,''
Citizen Gardens Research Initiative, Technical Report (2026).

\bibitem{Planck2018X}
Planck Collaboration,
``Planck 2018 results.\ X.\ Constraints on inflation,''
Astron.\ Astrophys.\ \textbf{641}, A10 (2020)
[arXiv:1807.06211].

\bibitem{BBO2011}
K.~Yagi and N.~Seto,
``Detector configuration of DECIGO/BBO and identification of cosmological neutron-star binaries,''
Phys.\ Rev.\ D \textbf{83}, 044011 (2011).

\bibitem{DECIGO2021}
S.~Kawamura \emph{et al.},
``Current status of space gravitational wave antenna DECIGO and B-DECIGO,''
PTEP \textbf{2021}, 05A105 (2021)
[arXiv:2006.13545].

\bibitem{BristolRH}
University of Bristol,
``Quantum physics sheds light on Riemann hypothesis,''
\url{https://www.bristol.ac.uk/maths/research/highlights/riemann-hypothesis/}.

\bibitem{QuantaRH}
N.~Wolchover,
``Physicists Attack Math's \$1,000,000 Question,''
Quanta Magazine (2017),
\url{https://www.quantamagazine.org/quantum-physicists-attack-the-riemann-hypothesis-20170404/}.

\end{thebibliography}

\end{document}